\RequirePackage[thinlines]{easybmat} %-- muss aufgrund eines Bugs vor etex und tikz geladen werden
\documentclass[a4paper, twoside, headsepline, index=totoc,toc=listof, fontsize=10pt, cleardoublepage=empty, headinclude, DIV=13, BCOR=13mm, titlepage,]{scrartcl}
\usepackage{scrtime} % KOMA, Uhrzeit ermoeglicht
\usepackage{scrpage2} % wie fancyhdr, nur optimiert auf KOMA-Skript, leicht andere Syntax
\usepackage{etoolbox}
\usepackage{letltxmacro}

% Compiler abhaengige IFs
\usepackage{ifxetex,ifluatex}
\newif\ifxetexorluatex
\ifxetex
  \xetexorluatextrue
\else
  \ifluatex
    \xetexorluatextrue
  \else
    \xetexorluatexfalse
  \fi
\fi

%--Farbdefinitionen
%-- muss vor tikz geladen werden
\usepackage[usenames, table, x11names]{xcolor}
\definecolor{dark_gray}{gray}{0.45}
\definecolor{light_gray}{gray}{0.6}

%--Zum Zeichnen
%-- muss vor polyglossia bzw. babel geladen werden
\usepackage{tikz}
\usepackage{tikz-cd}
\usetikzlibrary{external}
\tikzset{>=latex}
\usetikzlibrary{shapes,arrows,intersections}
\usetikzlibrary{calc,3d}
\usetikzlibrary{decorations.pathreplacing,decorations.markings}
\usetikzlibrary{angles}

%-- Konfiguration von tikzexternalize
\tikzexternalize[prefix=tikz/, up to date check=diff]

\ifxetexorluatex
\tikzset{external/system call={xelatex \tikzexternalcheckshellescape %-- verwende LuaLaTeX, wegen dynamischer Speicherallokation
    -halt-on-error -interaction=batchmode -jobname "\image" "\texsource"}}
\else
\tikzset{external/system call={pdflatex \tikzexternalcheckshellescape 
    -halt-on-error -interaction=batchmode -jobname "\image" "\texsource"}}
\fi

%-- tikzexternalize fuer tikzcd deaktivieren, da inkompatibel
\AtBeginEnvironment{tikzcd}{\tikzexternaldisable}
\AtEndEnvironment{tikzcd}{\tikzexternalenable}

%-- um Inkompatibilitaeten von quotes und polyglossia bzw. babel zu vermeiden
\tikzset{
  every picture/.append style={
    execute at begin picture={\shorthandoff{"}},
    execute at end picture={\shorthandon{"}}
  }
}
\usetikzlibrary{quotes}
\usepackage{pgfplots}
\usepgfplotslibrary{colormaps}
\pgfplotsset{compat=1.10}



%-- Mathe
\usepackage{mathtools} % beinhaltet amsmath
\mathtoolsset{showonlyrefs, centercolon}
\newtagform{brackets}[\textbf]{[}{]}
\usetagform{brackets}
\usepackage{wasysym}
\usepackage{amssymb} %zusätzliche Symbole
\usepackage{latexsym} %zusätzliche Symbole
\usepackage{stmaryrd} %für Blitz
\usepackage{nicefrac} %schräge Brüche
\usepackage{cancel} %Befehle zum Durchstreichen
\usepackage{mathdots}
% \usepackage{bm}
\DeclareSymbolFont{bbold}{U}{bbold}{m}{n}
\DeclareSymbolFontAlphabet{\mathbbold}{bbold}
\newcommand{\mathds}[1]{\mathbb{#1}} %Um Kompatibilitaet mit frueheren Benutzung von dsfont herzustellen
\newcommand{\ind}{\mathbbold{1}} % charakteristische-Funktion-Eins
\def\mathul#1#2{\color{#1}\underline{{\color{black}#2}}\color{black}} %farbiges Untersteichen im Mathe-Modus

%-- Alles was mit Schrift und XeteX zu tun hat
\ifxetexorluatex
    % XeLaTeX or LuaTeX
	\ifxetex
		\usepackage{mathspec} %beinhaltet fontspec 
	\else
		\usepackage[no-math]{fontspec}
	\fi
    \usepackage{polyglossia} %babel-ersatz
    \setmainlanguage[spelling=new,babelshorthands=false]{german}
	\newcommand\glqq{"}
	\newcommand\grqq{"}
	\defaultfontfeatures{Mapping=tex-text, WordSpace={1.4}} %
    \setmainfont[Ligatures=Common, BoldFont={* Bold}, ItalicFont={* Light Italic},SmallCapsFont={Linux Libertine O}, SmallCapsFeatures={Letters=SmallCaps}]{Source Sans Pro}
	\setsansfont[Scale=MatchLowercase,Ligatures=Common, BoldFont={* Medium}]{Ubuntu}
	\ifxetex
		\setallmonofonts[Scale=MatchLowercase, ItalicFont={*}]{Consolas} 
	\else
		\setmonofont[Scale=MatchLowercase, ItalicFont={*}]{Consolas}
	\fi
	\usepackage{xltxtra}
	\usepackage{fontawesome}
	\usepackage{microtype}
\else
    % default: pdfLaTeX
    \usepackage[ngerman]{babel}
    \usepackage[T1]{fontenc}
    \usepackage[utf8]{inputenc}
	\usepackage{textcomp} %verhindert ein paar Fehler bei den Fonts
    \usepackage[babel=true, tracking=true]{microtype}
	\usepackage{lmodern}
	\renewcommand{\familydefault}{\sfdefault} 
\fi


%--Mixed
\usepackage[neverdecrease]{paralist}
\usepackage[german=quotes]{csquotes}
\usepackage{makeidx}
\usepackage{booktabs}
\usepackage{wrapfig}
\usepackage{float}
\usepackage[margin=10pt, font=small, labelfont={sf, bf}, format=plain, indention=1em]{caption}
\captionsetup[wrapfigure]{name=Abb. }
\usepackage{stackrel}
\usepackage{ifthen}
\usepackage{multicol}
\usepackage{qrcode}
\flushbottom

%--Hyphenation
\hyphenation{di-men-si-o-na-le}


%--Unterstreichung
\usepackage[normalem]{ulem}
\setlength{\ULdepth}{1.8pt}

%--Indexverarbeitung
\newcommand{\bet}[1]{\textbf{#1}}
\newcommand{\Index}[1]{\textbf{#1}\index{#1}}
\makeindex
\setindexpreamble{{\noindent \itshape Die \emph{Seitenzahlen} sind mit Hyperlinks zu den entsprechenden Seiten versehen, also anklickbar} \ifxetexorluatex \faHandUp\fi \par \bigskip}
\renewcommand{\indexpagestyle}{scrheadings}


%--Marginnote & todonotes
\usepackage{marginnote}
\renewcommand*{\marginfont}{\color{dark_gray} \itshape \footnotesize }
\usepackage[textsize=small]{todonotes}

%--Todonotes mit tikz/externalize kompatibel machen
\LetLtxMacro{\oldtodo}{\todo}
\renewcommand{\todo}[2][]{\tikzexternaldisable\oldtodo[#1]{#2}\tikzexternalenable}
\LetLtxMacro{\oldmissingfigure}{\missingfigure}
\renewcommand{\missingfigure}[2][]{\tikzexternaldisable\oldmissingfigure[{#1}]{#2}\tikzexternalenable}

%--Konfiguration von Hyperref pdfstartview=FitH, 
\usepackage[hidelinks, pdfpagelabels,  bookmarksopen=true, bookmarksnumbered=true, linkcolor=black, urlcolor=SkyBlue2, plainpages=false,pagebackref, citecolor=black, hypertexnames=true, pdfauthor={Jannes Bantje}, pdfborderstyle={/S/U}, linkbordercolor=SkyBlue2, colorlinks=false]{hyperref}

% \let\orighref\href
\ifxetexorluatex
\newcommand{\hrefsym}[2]{\href{#1}{\texttt{#2\,\faExternalLink}}}
\newcommand{\hrefsymmail}[2]{\href{#1}{\texttt{\faEnvelope\,#2}}}
\renewcommand{\url}[1]{\hrefsym{#1}{\nolinkurl{#1}}}
\fi
\providecommand{\hrefsym}{\href}
% \usepackage{cleveref}
% \crefformat{footnote}{#2\footnotemark[#1]#3}

%--Römische Zahlen
\newcommand{\RM}[1]{\MakeUppercase{\romannumeral #1{}}}

%--Punkte (nach hyperref)
\usepackage{ellipsis}

%%--Abkürzungen etc.
\newcommand{\light}{\text{\Large $\lightning$}}

%-- Definitionen von weiteren Mathe-Befehlen
\DeclareMathOperator{\re}{Re} %Realteil
\DeclareMathOperator{\im}{Im} %Imaginaerteil
\DeclareMathOperator{\id}{id} %identische Abbildung
\DeclareMathOperator{\Sp}{Sp} %Spur
\DeclareMathOperator{\sgn}{sgn} %Signum
\DeclareMathOperator{\alt}{Alt} %Alternierende n-linearFormen
\DeclareMathOperator{\End}{End} %Endomorphismen
\DeclareMathOperator{\Vol}{Vol} %Volumen
\DeclareMathOperator{\dom}{dom} %Domain
\DeclareMathOperator{\Hom}{Hom} %Homomorphismen
\DeclareMathOperator{\bild}{Bild} %Bild
\DeclareMathOperator{\Kern}{Kern }
\DeclareMathOperator{\rg}{rg} %Rang
\DeclareMathOperator{\Rg}{Rg} %Rang
\DeclareMathOperator{\diam}{diam} %Durchmesser
\DeclareMathOperator{\dist}{dist} %Distanz
\DeclareMathOperator{\grad}{grad} %Gradient
\DeclareMathOperator{\dive}{div} %Gradient
\DeclareMathOperator{\rot}{rot} %Rotation
\DeclareMathOperator{\hess}{Hess} %Hesse-Matrix
\DeclareMathOperator{\koker}{Koker} %Kokern
\DeclareMathOperator{\aut}{Aut} %Automorphismen
\DeclareMathOperator{\ord}{ord} %Ordnung
\DeclareMathOperator{\ggT}{ggT} %ggT
\DeclareMathOperator{\kgV}{kgV} %kgV
\DeclareMathOperator{\Gr}{Gr} %Gerade
\DeclareMathOperator{\Kr}{Kr} %Kreis
\DeclareMathOperator{\Char}{char} %Charakteristik
\DeclareMathOperator{\Aut}{Aut} %Automorphismen
% \DeclareMathOperator{\D}{D} %formale Ableitung
\newcommand{\D}{\ensuremath{\mathrm{D}\mkern-1.5mu}}
\newcommand{\mathd}{\ensuremath{\mathrm{d}\mkern-1.0mu}}
\newcommand{\Tmap}{\ensuremath{\mathrm{T}\mkern-1.0mu}}
\newcommand{\opL}{\ensuremath{\mathrm{L}\mkern-0.6mu}}
\DeclareMathOperator{\cov}{cov} %Kovarianz
\DeclareMathOperator{\Gal}{Gal} %Galois
\DeclareMathOperator{\supp}{supp} %Träger
\DeclareMathOperator{\Sym}{Sym} %Symmetrische Gruppe
\DeclareMathOperator{\Zyl}{Zyl} %Zylinder
\DeclareMathOperator{\Mod}{Mod} %Moduln
\DeclareMathOperator{\EPK}{EPK} %Einpunktkompaktifizierung
\DeclareMathOperator{\conj}{conj} %Konjugation
\DeclareMathOperator{\ann}{ann} %Annulator
\DeclareMathOperator{\tor}{tor} %Torsionsmodul
\DeclareMathOperator{\Int}{Int} %Interior/Inneres
\DeclareMathOperator{\Br}{Br} %Brauerergruppe
\DeclareMathOperator{\GL}{GL} %allgemeine lineare Gruppe
\DeclareMathOperator{\SL}{SL} %spezielle lineare Gruppe
\DeclareMathOperator{\SU}{SU} %spezielle unitäre Gruppe
\DeclareMathOperator{\Tr}{Tr} %Spur/Trace
\DeclareMathOperator{\conv}{conv} %Konvexe Teilmenge
% \DeclareMathOperator{\tr}{tr}
\newcommand{\tr}{\mathrm{tr}}
\newcommand{\w}{\mathrm{w}}
\newcommand{\weakT}[1]{\ensuremath{\mathcal{T}_{#1}^{\mkern+1.0mu\text{\raisebox{0.4ex}{$\mathrm{w}$}}}}}
\newcommand{\weakTstar}[1]{\ensuremath{\mathcal{T}_{#1}^{\mkern+1.0mu\text{\raisebox{0.4ex}{$\mathrm{w}$}}^*}}}
\newcommand{\rand}[1]{\ensuremath{\partial^{\scriptscriptstyle #1}}}


% -- Kategorien
\DeclareMathOperator{\Ob}{Ob} %Objekte
\DeclareMathOperator{\Mor}{Mor} %Morphismen

% \DeclareMathOperator{\TOP}{TOP} %Topologische Räume
\newcommand{\TOP}{\textsc{Top}}
% \DeclareMathOperator{\HTOP}{HTOP} %Topologische Räume + Homotopieäquivalenz
\newcommand{\HTOP}{\textsc{HTop}}
% \DeclareMathOperator{\VR}{VR} %Vektorräume
\newcommand{\VR}{\textsc{VR}}
% \DeclareMathOperator{\MOD}{MOD} %Moduln
\newcommand{\MOD}{\textsc{Mod}}
% \DeclareMathOperator{\MENGEN}{MENGEN} %Mengen
\newcommand{\MENGEN}{\textsc{Mengen}}
% \DeclareMathOperator{\MAN}{MAN} %Mannigfaltigkeiten
\newcommand{\MAN}{\textsc{Man}}
% \DeclareMathOperator{\GRUPPEN}{GRUPPEN} %Gruppen
\newcommand{\GRUPPEN}{\textsc{Gruppen}}
% \DeclareMathOperator{\ABELGRUPPEN}{ABEL.GRUPPEN} %abelsche Gruppen
\newcommand{\ABELGRUPPEN}{\textsc{Abel.Gruppen}}
% \DeclareMathOperator{\KAT}{KAT} %Kategorien
\newcommand{\KAT}{\textsc{Kat}}
% \DeclareMathOperator{\FUN}{FUN} %Funktoren
\newcommand{\FUN}{\textsc{Fun}}


\newcommand{\op}{\mathrm{op}}

\newcommand{\Ker}{\ensuremath{\Kern \,}}
\newcommand{\zirkel}{\rotatebox[origin = c]{-90}{$\varangle$}}
\newcommand{\wirkung}{\!\curvearrowright\!}

\newcommand{\oE}{\OE}

%--Beweisende
\newcommand{\bewende}{\ifmmode \tag*{$\square$} \else \hfill $\square$ \fi}

%--Volumen
\newcommand{\vol}[1]{\ensuremath{\Vol_{#1}}}

%--Skalarprodukt
\DeclarePairedDelimiterX\sprod[2]{\langle}{\rangle}{#1\,\delimsize\vert\,#2}

%--Betrag, Gaußklammer
\DeclarePairedDelimiter{\abs}{\lvert}{\rvert}
\DeclarePairedDelimiter{\ceil}{\lceil}{\rceil}

%--Norm
\DeclarePairedDelimiter\doppelstrich{\Vert}{\Vert}
\newcommand{\norm}[2][\relax]{
\ifx#1\relax \ensuremath{\doppelstrich*{#2}}
\else \ensuremath{\doppelstrich*{#2}_{#1}}
\fi}


%--Umklammern
\DeclarePairedDelimiter\enbrace{(}{)}
\DeclarePairedDelimiter\benbrace{[}{]}
\DeclarePairedDelimiter\lenbrace{<}{>}

%--Mengen
\DeclarePairedDelimiterX\mengenA[1]{\lbrace}{\rbrace}{#1}
\DeclarePairedDelimiterX\mengenB[2]{\lbrace}{\rbrace}{#1\, \delimsize\vert \, #2}

\makeatletter
\newcommand{\set}[2][\relax]{
\ifx#1\relax \ensuremath{
\mengenA*{#2}}
\else \ensuremath{%
  \mengenB*{#1}{#2}}
\fi}
\makeatother


%--offen und abgeschlossen
\newcommand{\off}{\! \stackrel[\text{offen}]{}{\subset} \!}
\newcommand{\abg}{\! \stackrel[\text{abg.}]{}{\subset} \!}

%--Differential
\newcommand{\diff}[2]{\ensuremath{\frac{{\partial #1}}{{\partial #2}} }}
\newcommand{\diffd}[2]{\ensuremath{\frac{\mathd #1}{\mathd #2} }}

%--Diagonale Linien für Matrizen
\newcommand\tikzmark[1]{%
  \tikzexternaldisable\tikz[overlay,remember picture,baseline] \node [anchor=base] (#1) {};\tikzexternalenable}

\newcommand\MyLine[3][]{%
  \tikzexternaldisable\begin{tikzpicture}[overlay,remember picture]
    \draw[#1] (#2.south) -- (#3.north);
  \end{tikzpicture}\tikzexternalenable}
  
%--Jordankasten
\newcommand{\Jord}{ \tikzexternaldisable \begin{BMAT}{|ccc|}{|ccc|} %
				0 & & \\ %
				1\tikzmark{z} &  \diagdown & \\ %
				& \tikzmark{u}1 \MyLine[thick]{z}{u} &  0 %
			\end{BMAT} \tikzexternalenable}
			
%--Jordankasten mit Lambda
\newcommand{\JordLa}[1][\empty]{%
	\tikzexternaldisable\ifthenelse{\equal{#1}{\empty}}
		{\begin{BMAT}{|ccc|}{|ccc|} %
				\lambda & & \\ %
				1\tikzmark{z} &  \diagdown & \\ %
				& \tikzmark{u}1 \MyLine[thick]{z}{u} &  \lambda  %
			\end{BMAT}}
		{\begin{BMAT}{|ccc|}{|ccc|} %
				\lambda_{#1}  & & \\ %
				1\tikzmark{z} &  \diagdown & \\ %
				& \tikzmark{u}1 \MyLine[thick]{z}{u} &  \lambda_{#1}  %
			\end{BMAT}}\tikzexternalenable}

%--Inhaltsverzeichnis
\usepackage[tocindentauto]{tocstyle}
\usetocstyle{KOMAlike}			

\newcommand{\fach}{Einführung in die Algebra}
\newcommand{\semester}{WiSe 2013}
\newcommand{\homepage}{http://wwwmath.uni-muenster.de/reine/u/topos/lehre/WS2013-2014/Algebra/algebra.html}
\newcommand{\prof}{Prof.\,Dr.\,Arthur Bartels}
\providecommand{\verfasser}{Jannes Bantje}

%--Konfiguration von scrheadings
\setheadsepline{1pt}[\color{light_gray}]
\pagestyle{scrheadings}
%\fancyhf[H, F]{}
\clearscrheadfoot

\providecommand{\shortFach}{\fach}
\lehead{ \includegraphics[height=0.5 cm,keepaspectratio]{../!config/Bilder/Logo_WWU_Muenster_light_gray.pdf}}
\rehead{\rule{0cm}{0.5cm}\footnotesize \sffamily \color{light_gray} \verfasser -- Mitschrift \shortFach}
\lohead{\rule{0cm}{0.5cm} \footnotesize \sffamily \color{light_gray} Stand: \today \; \thistime[:]}
\rohead{\includegraphics[height=0.5 cm,keepaspectratio]{../!config/Bilder/fb10logo_gray.pdf}}


\ofoot[{ \color{dark_gray} \LARGE \sffamily \thepage}]{{ \color{dark_gray} \LARGE \sffamily \thepage}} %hier wir auch der plain Stil bearbeitet!
\automark{section}
\ifoot{ \color{dark_gray} \small \leftmark}

%--Metadaten
\author{\verfasser}
\titlehead{\includegraphics[height=1.5cm, keepaspectratio]{../!config/Bilder/Logo_WWU_Muenster.pdf}%
\hfill \includegraphics[height=1.3cm, keepaspectratio]{../!config/Bilder/fb10logo.pdf}}
\title{Skript \fach}
\subtitle{Mitschrift der Vorlesung  \enquote{\fach} von \prof}

\publishers{
Aktuelle Version verfügbar bei: \bigskip\\
\normalsize
\begin{minipage}[t]{0.45\textwidth}
	\centering{
	\qrcode[height=3.5cm, version=4]{https://github.com/JaMeZ-B/latex-wwu}
	\medskip\\
	\includegraphics[height=0.4cm, keepaspectratio]{../!config/Bilder/github_octo.pdf} 
	\includegraphics[height=0.4cm, keepaspectratio]{../!config/Bilder/GitHub_Logo.pdf} (inklusive Sourcecode)\\
	{\small\url{https://github.com/JaMeZ-B/latex-wwu}}}
\end{minipage}
\hfill
\begin{minipage}[t]{0.45\textwidth}
	\centering{
	\qrcode[height=3.5cm, version=4]{btsync://B6WH2DISQ5QVYIRYIEZSF4ZR2IDVKPN3I?n=latex_share}
	\medskip\\
	\includegraphics[height=0.4cm, keepaspectratio]{../!config/Bilder/bt_sync_logo.pdf} 
	{\large \textbf{Bittorrent} Sync} \\
	\texttt{B6WH2DISQ5QVYIRYIEZSF4ZR2IDVKPN3I}
	}
\end{minipage}
}


\begin{document}
\pagenumbering{Roman}
\maketitle
\begin{abstract}
\section*{Vorwort --- Mitarbeit am Skript}
Dieses Dokument ist eine Mitschrift aus der Vorlesung \enquote{\fach, \semester}, gelesen von \prof. Der Inhalt entspricht weitestgehend dem Tafelanschrieb. Für die
Korrektheit des Inhalts übernehme ich keinerlei Garantie! Für Bemerkungen und Korrekturen -- und seien es nur Rechtschreibfehler -- bin ich sehr dankbar. 
Korrekturen lassen sich prinzipiell auf drei Wegen einreichen: 
\begin{itemize}
	\item Persönliches Ansprechen in der Uni, Mails an \hrefsymmail{mailto:\mail}{\mail} (gerne auch mit annotieren PDFs) 
	\item \emph{Direktes} Mitarbeiten am Skript: Den Quellcode poste ich auf GitHub (siehe oben), also stehen vielfältige Möglichkeiten der Zusammenarbeit zur Verfügung:
	Zum Beispiel durch Kommentare am Code über die Website und die Kombination Fork + Pull Request. Wer sich verdient macht oder ein Skript zu einer Vorlesung, die 
	ich nicht besuche, beisteuern will, dem gewähre ich gerne auch Schreibzugriff.
	
	Beachten sollte man dabei, dass dazu ein Account bei \url{github.com} notwendig ist, der allerdings ohne Angabe von persönlichen Daten angelegt werden kann. 
	Wer bei GitHub (bzw. dem zugrunde liegenden Open-Source-Programm \enquote{\texttt{git}}) -- verständlicherweise -- Hilfe beim Einstieg braucht, dem helfe ich gerne 
	weiter. Es gibt aber auch zahlreiche empfehlenswerte Tutorials im Internet.\footnote{zB. \url{https://try.github.io/levels/1/challenges/1}, ist auf Englisch, aber dafür 
	interaktives LearningByDoing}
	\item \emph{Indirektes} Mitarbeiten: \TeX-Dateien per Mail verschicken. 
	
	Dies ist nur dann sinnvoll, wenn man einen ganzen Abschnitt ändern möchte (zB. einen alternativen Beweis geben), da ich die Änderungen dann per Hand einbauen muss! Ich freue mich aber auch über solche Beiträge!
\end{itemize}
\section*{Vorlesungshomepage}
{\centering 
\begin{minipage}[c][][c]{\textwidth}
	\centering \qrcode[height=3cm]{\homepage} \medskip\\
	\footnotesize \url{\homepage}
\end{minipage}
\par}

\end{abstract}

\setcounter{page}{1}
\tableofcontents
\cleardoubleoddemptypage
\pagenumbering{arabic}
\setcounter{page}{1}



\section{Gruppen} % (fold)
\label{sec:1}

\subsection[Definition: Monoid]{Definition} % (fold)
\label{sub:11}
Ein \Index{Monoid} ist eine Menge $M$ mit einer Verknüpfung $M \times M \to M, (a,b) \mapsto a \cdot b$ so dass folgende Axiome erfüllt sind:
\begin{enumerate}[(i)]
	\item $\forall a,b,c \in M : (a \cdot b )\cdot c = a \cdot (b \cdot c)$
	\item $\exists e \in M : \forall a \in M : a \cdot e = a = e \cdot a$
\end{enumerate} 
Gilt $\forall a,b \in M$ $a \cdot b = b \cdot a$ so heißt $M$ \Index{kommutativ} oder \Index{abelsch}. Beispiel: 
\begin{multicols}{2}
	\begin{enumerate}[(i)]
		\item $(\mathds{N},+)$
		\item $(\mathds{Z},+)$
		\item $(\mathds{N}, \cdot )$
		\item $\enbrace*{N \cup \{ \infty\}, +} $ mit $a + \infty = \infty = \infty + a$ $\forall a \in \mathds{N} \cup \{ \infty \}$
	\end{enumerate}
\end{multicols}
Bemerkung: Das neutrale Element $e$ ist eindeutig.
% subsection 11 (end)

\subsection[Definition: Gruppe]{Definition} % (fold)
\label{sub:12}
Eine \Index{Gruppe} ist ein Monoid $G$ in dem jedes Element ein Inverses besitzt:
\[
	\forall a \in G : \enspace \exists b \in G \text{ mit } a \cdot b = e = b \cdot a
\] 
% subsection 12 (end)

\subsection[Bemerkung: Eindeutigkeit des Inversen]{Bemerkung} % (fold)
\label{sub:13}
Das Inverse zu $a$ ist eindeutig und wird mit $a ^{-1}$ bezeichnet.
% subsection 13 (end)

\subsection[Bemerkung: Lösbarkeit einfacher Gleichungen in Gruppen]{Bemerkung} % (fold)
\label{sub:14}
Sei $G$ eine Gruppe. $a, b \in G$. Dann besitzt die Gleichung $a \cdot x = b$ eine eindeutige Lösung, nämlich $x = a ^{-1} \cdot b$
% subsection 14 (end)

\subsection[Beispiele für Gruppen]{Beispiele} % (fold)
\label{sub:15}
	\begin{enumerate}[i)]
		\item $(\mathds{Z}, +)$
		\item $ (\mathds{Q},+)$
		\item $(\mathds{Q}^\times, \cdot )$ \hfill {\small \color{light_gray}(Menge der Einheiten)}
		\item $Gl(V)$ \hfill {\small \color{light_gray} (Menge aller invertierbaren Endomorphismen über $V$)}
		\item $Gl(n,K)$ \hfill {\small \color{light_gray} (Menge aller invertierbaren $n \times n$-Matrizen über den Körper $K$)}
		\item \label{sym_grp} Sei $X$ eine Menge. Dann heißt $S_X := \set[\sigma : X \to X]{\sigma \text{ ist bijektiv}} $ die \Index{symmetrische Gruppe} von $X$. 
		( $S_n := S_{ \{ 1, \ldots , n\}} $) 
		\item Die Symmetriegruppe $S_\square$ des Quadrats
		\item Die Symmetriegruppe $S_\bigcirc$ des Kreises
		\item Sei $K$ ein Körper, dann ist die Menge $\aut(K)$ der \bet{Körperautomorphismen} \index{Automorphismen!Körperautomorphismen} von $K$ eine Gruppe
		\item Sei $p = X^n + a_{n-1} X^{n-1} + \ldots  + a_0 \in K[X]$ ein Polynom. Wir werden später die Nullstellen von $p$ untersuchen, in dem wir $p$ eine Gruppe 
		zuordnen
		\item Sei $\varphi$ eine Kategorie und $M$ ein Objekt von $\varphi$. Dann ist die Menge $\aut_\varphi(M)$ aller invertierbaren Morphismen $f : M  \to M$ in $\varphi$
		eine Gruppe.
	\end{enumerate}
% subsection 15 (end)

\subsection{Verknüpfungstafeln} % (fold)
\label{sub:16}
Verknüpfungstafeln für die Gruppen $\enbrace*{\nicefrac{\mathds{Z}}{4 \mathds{Z}}, +} $ und $(S_{\triangle}, \circ)$
\begin{table}[h]
	\centering{
	\begin{tabular}{c|cccc}
	$+$ & $[0]$ & $[1]$ & $[2]$ & $[3]$\\
	\hline
	$[0]$ & $[0]$ & $[1]$ & $[2]$ & $[3]$\\
	
	$[1]$ & $[1]$ & $[2]$ & $[3]$ & $[0]$\\
	
	$[2]$ & $[2]$ & $[3]$ & $[0]$ & $[1]$\\
	
	$[3]$ & $[3]$ & $[0]$ & $[1]$ & $[2]$\\
	\hline
	\end{tabular}
	\qquad 
	\begin{tabular}{c|cccccc}
	$\circ $ & $\id$ & $R_1$ & $R_2$ & $S_A$ & $S_B$ & $S_C$\\
	\hline
	$\id$ & $\id$ & $R_1$ & $R_2$ & $S_A$ & $S_B$ & $S_C$\\
	
	$R_1$ & $R_1$ & $R_2$ & $\id$ & $S_C$ & $S_A$ & $S_B$\\
	
	$R_2$ & $R_2$ & $\id$ & $R_1$ & $S_B$ & $S_C$ & $S_A$\\
	
	$S_A$ & $S_A$ & $S_B$ & $S_C$ & $\id$ & $R_1$ & $R_2$\\
	
	$S_B$ & $S_B$ & $S_C$ & $S_A$ & $R_2$ & $\id$ & $R_1$\\
	
	$S_C$ & $S_C$ & $S_A$ & $S_B$ & $R_1$ & $R_2$ & $\id$\\
	\hline
	\end{tabular}
	\caption{Verknüpfungstafeln für $\enbrace*{\nicefrac{\mathds{Z}}{4 \mathds{Z}}, +} $ und $(S_{\triangle}, \circ )$}
	}
\end{table}
% subsection 16 (end)

\subsection[Definition: Mächtigkeit, Ordnung einer Gruppe]{Definition} % (fold)
\label{sub:17}
Die \Index{Ordnung} einer Gruppe $G$ ist ihre Mächtigkeit $\abs*{G} $. 
% subsection 17 (end)

\subsection[Definition: Untergruppe]{Definition} % (fold)
\label{sub:18}
Sei $G$ eine Gruppe. Eine Teilmenge $U \subseteq G$ heißt eine \Index{Untergruppe} falls gilt:
\begin{multicols}{3}
	\begin{enumerate}[(i)]
		\item $e \in U$
		\item $\forall a,b \in U$ ist $a \cdot b \in U$
		\item $\forall a \in U : a ^{-1} \in U$
	\end{enumerate}
\end{multicols}
% subsection 18 (end)

\subsection[Definition: Gruppenhomomorphismus]{Definition} % (fold)
\label{sub:19}
Seien $G$ und $H$ Gruppen. Eine Abbildung $\varphi : G \to H$ heißt ein \Index{Gruppenhomomorphismus}, falls gilt:
\[
	\forall a,b \in G : \varphi(a \cdot b) = \varphi(a) \cdot \varphi(b)
\] 
% subsection 19 (end)

\subsection[Bemerkung: Eigenschaften von Gruppenhomomorphismen]{Bemerkung} % (fold)
\label{sub:110}
Ist $\varphi$ ein Gruppenhomomorphismus, so gilt: 
\begin{multicols}{2}
	\begin{enumerate}[(i)]
		\item $\varphi(e) =e$
		\item $\forall a \in G : \varphi \enbrace*{ a ^{-1}} = \varphi(a) ^{-1} $
	\end{enumerate}
\end{multicols}
% subsection 110 (end)

\subsection[Beispiele für Gruppenhomomorphismen]{Beispiel} % (fold)
\label{sub:111}
\begin{enumerate}[(i)]
	\item Sei $n \in \mathds{N}$. Dann ist $\varphi_n : \mathds{Z} \to \mathds{Z}$ mit $\varphi_n(k)= n \cdot k$ ein Gruppenhomomorphismus. \marginnote{gemeint: $(\mathds{Z},+)$}
	\item Sei $G$ eine Gruppe. Zu $a \in G$ sei $L_a : G \to G$ mit $L_a(g) = a \cdot g$ \marginnote{$L_a$ ist kein Gruppenhomomorphismus!}
	Dann ist $G \xrightarrow{\varphi}  S_G$, $a \mapsto L_a$ ein injektiver Gruppenhomomorphismus. Behauptung:
	\[
		\varphi(a \cdot b) = \varphi(a) \cdot \varphi(b)
	\]
	Beweis:
	\[
		a \cdot b \cdot g =L_{ab} (g) = L_a \circ L_b (g) = L_a \big( L_b(g)\big) = a \cdot b \cdot g
	\]
	Injektivität: Seien $a,b \in G$ mit $L_a = \varphi(a) = \varphi(b) = L_b$. Es gilt $a \cdot g = L_a (g) = L_b (g) = b g \Rightarrow  a = b$
	\item Sei $G$ eine Gruppe und $g \in G$. Dann ist $\mathrm{conj}_g : G \to G$ mit $\mathrm{conj}_g (a) := g \cdot a \cdot g ^{-1}$ ein Automorphismus (bijektiver
	Gruppenhomomorphismus $G \to G$) von $G$.
	
	Automorphismen dieser Form heißen \bet{innere Automorphismen} \index{Automorphismen!innere}.
\end{enumerate}
% subsection 112 (end)

\subsection[Definition: Bild und Kern eines Gruppenhomomorphismus]{Definition} % (fold)
\label{sub:112}
Sei $\varphi : G \to H$ ein Gruppenhomomorphismus. Dann heißt $\bild (\varphi) = \set[\varphi(g)]{g \in G} $ das \Index{Bild} von $\varphi$ und 
$\Kern(\varphi) = \set[g]{\varphi(g) = e_H} $ der \Index{Kern} von $\varphi$.

Bemerkung: \begin{itemize}
	\item $\Kern (\varphi)$ und $\bild ( \varphi)$ sind Untergruppen von $G$ bzw. $H$.
	\item $\varphi$ surjektiv $\Leftrightarrow \bild (\varphi) = H$. $\varphi$ injektiv $\Leftrightarrow \Kern \varphi = \{ e_G\}$
\end{itemize}
% subsection 113 (end)

\subsection[Beispiel: Das kartesische Produkt von Gruppen ist auch eine Gruppe]{Beispiel} % (fold)
\label{sub:113}
Seien $G$ und $H$ Gruppen. Dann ist auch das kartesische Produkt $G \times H$ eine Gruppe
\[
	(g,h) \cdot (g', h') = (g g', h h') , e_{G\times H} = (e_G, e_H), (g,h) ^{-1} = (g ^{-1} , h ^{-1})
\]
% subsection 114 (end)

\subsection[Lemma: Endliche Monoide, die links- und rechtskürzbar sind, sind Gruppen]{Lemma} % (fold)
\label{sub:114}
Sei $M$ ein endlicher Monoid. Es gelte
\begin{align*}
	\forall a,b,b' \in G : a \cdot b = a \cdot b' &\Rightarrow b=b' \\
	b \cdot a = b' \cdot a &\Rightarrow b=b'
\end{align*}
Dann ist $M$ eine Gruppe. Beweis: Übung
% subsection 115 (end)

\subsection[Definition: Linksnebenklasse und Index]{Definition} % (fold)
\label{sub:115}
Sei $G$ eine Gruppe, $U \subseteq G$ eine Untergruppe. Eine \Index{Linksnebenklasse} von $U$ ist eine Teilmenge von $G$ der Gestalt $a\, U := \set[a \cdot u]{u \in U} $
wobei $a \in G$. Die Menge aller Linksnebenklassen bezeichnen wir mit  $\nicefrac{G}{U} := \set[a\, U]{a \in G} $. Die Anzahl der Linksnebenklassen heißt der \Index{Index} 
$[G:U]$ von $U$ in $G$.

Beispiel: $G= \mathds{Z}, U = 2\mathds{Z} = \set[n \in \mathds{Z}]{n \text{ ist gerade}} $. Dann benutzen wir additive Schreibweise $n+U$ für die Linksnebenklasse von $U$.
Dann:
\begin{enumerate}[i)]
	\item $2\mathds{Z} = 0+2\mathds{Z} = 2+2\mathds{Z} = 4+2\mathds{Z}= -2+2\mathds{Z}= \ldots $
	\item $1+2\mathds{Z} = 3+2\mathds{Z} = 5+2\mathds{Z}=-1+2\mathds{Z} = \ldots $
	\item $\mathds{Z} = 2\mathds{Z} \, \dot\cup \, (1+2\mathds{Z})$
\end{enumerate}
% subsection 116 (end)

\subsection[Lemma: Äquivalente Aussagen zur Gleichheit von Nebenklassen]{Lemma} % (fold)
\label{sub:116}
Seien $a \,U$ und $b\,U$ Linksnebenklassen von $U$ in $G$. Dann sind äquivalent: 
\begin{multicols}{2}
	\begin{enumerate}[i)]
		\item $a\,U=b\,U$
		\item $a\,U \cap b\,U \not= \emptyset$
		\item $b ^{-1} a \in U$
		\item $a \in b\,U$
	\end{enumerate}
\end{multicols}
\minisec{Beweis}
\begin{description}
	\item[i)$\Rightarrow $ii)] Ist $a\, U=b\, U$ dann ist $a\, U \cap b\, U = a\, U$. Da $e \in U$ ist $a=a \cdot e \in a\,U$, also $b \,U \cap a \,U \not= \emptyset$
	\item[ii)$\Rightarrow $iii)] Sei $g \in a\, U \cap b\, U$. Dann $g= a \cdot u = b \cdot v$ mit $u,v \in U$. $\Rightarrow  b ^{-1} \cdot a = v \cdot  u ^{-1} \in U$.
	\item[iii)$\Rightarrow $iv)] Ist $b ^{-1} \cdot a \in U$, so folgt $a= b \cdot (b ^{-1} a) \in b\, U$
	\item[iv)$\Rightarrow $i)] Sei $a \in b\, U$. Dann $a=b \cdot v$ mit $v \in U$. Es folgt 
	\[
		a \,U = \set[a \cdot u]{u \in U} = \set[b \cdot v \cdot u]{u \in U} = \set[b \cdot u']{u' \in U} = b\, U \bewende
	\]  
\end{description}
% subsection 117 (end)

\subsection[Lemma: Je zwei Linksnebenklassen sind gleichmächtig]{Lemma} % (fold)
\label{sub:117}
Je zwei Linksnebenklassen von $U$ in $G$ sind gleichmächtig, $\abs*{a \, U} = \abs*{b \, U}$.
\minisec{Beweis}
Für $a \in G$ ist $U \to a\, U, u \mapsto a \cdot u$ eine bijektive Abbildung ($x \mapsto a ^{-1} x$ ist die inverse Abbildung). \bewende
% subsection 118 (end)

\subsection{Satz von Lagrange} % (fold)
\label{sub:118}
Sei $G$ eine endliche Gruppe und $U$ eine Untergruppe von $G$. Dann gilt: 
\[
	\abs*{G} =  [G:U] \cdot \abs*{U}  
\]
\minisec{Beweis}
$G$ ist die disjunkte Vereinigung der Linksnebenklassen (Da $g \in g U$, ist $G$ die Vereinigung der Linksnebenklassen, nach dem 1. Lemma ist diese disjunkt). Also
\[
	\abs*{G} = \abs*{\dot \bigcup_{aU \in \nicefrac{G}{U}} a U}  = \sum_{a U \in \nicefrac{G}{U}} \abs*{aU} = \sum_{a U \in \nicefrac{G}{U}} \abs*{U} = \abs*{\nicefrac{G}{H}}
	\cdot \abs*{U} = [G:U] \cdot \abs*{U}       \bewende
\]
% subsection 119 (end)

\subsection[Bemerkung zu Rechtsnebenklassen]{Bemerkung} % (fold)
\label{sub:119}
Analog zu Linksnebenklassen kann man die \Index{Rechtsnebenklassen} $U a = \set[u \cdot  a]{u \in U} $ betrachten. Die Menge der Rechtsnebenklassen bezeichnen wir mit $G \backslash U$.

Übung: 
\begin{enumerate}[(i)]
	\item $\abs*{\nicefrac{G}{U}} = \abs*{G \backslash U}  $
	\item $\abs*{U} = \abs*{U \,a}  $
	\item $U a = U b \Leftrightarrow U a \cap U b \not= \emptyset \Leftrightarrow a \in U b \Leftrightarrow a b ^{-1} \in U$
\end{enumerate} 
% subsection 120 (end)

\subsection[Definition: Normalteiler]{Definition} % (fold)
\label{sub:120}
Eine Untergruppe $N \le G$ heißt ein \Index{Normalteiler}, wenn $a N = N a$ ist für alle $a \in G$. 

Bemerkung: Ist $G$ abelsch, so ist jede Untergruppe ein Normalteiler.
% subsection 121 (end)

\subsection[Lemma: Kriterium für Normalteiler]{Lemma} % (fold)
\label{sub:121}
Eine Untergruppe $N \le G$ ist genau dann ein Normalteiler, wenn gilt:
\[
	\forall a \in G : a N a ^{-1} = N
\]
(Hier ist $a N a ^{-1} = \set[a n a ^{-1}]{n \in N} $)

Beispiel: Ist $\varphi$ ein Gruppenhomomorphismus, so ist $\Kern \varphi \subseteq G$ ein Normalteiler.
% subsection 122 (end)

\subsection[Lemma: Für einen Normalteiler $N$ hat $\nicefrac{G}{N}$ Gruppenstruktur]{Lemma} % (fold)
\label{sub:122}
Sei $N \le G$ ein Normalteiler. dann wird auf $\nicefrac{G}{N}$ durch $aN \cdot bN := abN$ eine wohldefinierte Gruppenstruktur erklärt. Ihr neutrales Element ist $eN=N$.
Es gilt $(aN)^{-1} = a ^{-1}N$. Durch $\pi : G \to \nicefrac{G}{N}, a \mapsto aN$ ist ein Gruppenhomomorphismus mit $\Kern \pi = N$ definiert.
% subsection 123 (end)

\subsection[Definition: Faktorgruppe]{Definition} % (fold)
\label{sub:123}
$\nicefrac{G}{N}$ heißt \Index{Faktorgruppe} oder \Index{Restklassengruppe} von $G$ nach $N$.

Beispiel: $G=\mathds{Z}, N = n \mathds{Z}, \nicefrac{G}{\mathds{N}}= \nicefrac{\mathds{Z}}{n \mathds{Z}}$.
\minisec{Beweis der Lemmas}
Zur Wohldefiniertheit von $aN \cdot bN := abN$ ist zu zeigen: Ist $a_1 N = a_2 N$ und $b_1 N = b_2 N$ so ist $a_1 b_1N = a_2 b_2 N$. Sei also $a_1 N = a_2 N$ und 
$b_1 N = b_2 N$. Es folgt $a:= a_2 ^{-1} a_1 \in N$ und $b := b_2 ^{-1} b_1 \in N$. Dann ist 
\[
	a_1 \, b_1 \, N = a_2 \,a_2 ^{-1} \,a_1 \,b_2\, b_2 ^{-1}\, b_1\, N = a_2\, a\, b_2\, b\, N = a_2\, b_2\, ( \underbrace{b_2 ^{-1}\, a\, b_2}_{\in N})\, \underbrace{b}_{\in N}\, N = a_2\, b_2\, N \bewende
\]
% subsection 124 (end)

\subsection{Homomorphiesatz} % (fold)
\label{sub:124}
\begin{wrapfigure}[5]{R}{1cm}
	\begin{tikzcd}
		G \dar[two heads, swap]{\pi } \rar{\varphi}& H \\
		\nicefrac{G}{N} \urar[dashed, swap]{\overline{\varphi} }
	\end{tikzcd}
\end{wrapfigure}
Sei $\varphi : G \to H$ ein Gruppenhomomorphismus und $N \unlhd G$ ein Normalteiler. Gilt $N \subseteq \Kern \varphi$, so gibt es einen eindeutigen Gruppenhomomorphismus
$\overline{\varphi} : \nicefrac{G}{N} \to H $ mit $\varphi = \overline{\varphi} \circ \pi$.

Es gilt $\Kern \overline{\varphi} = \pi ( \Kern \varphi) $ und $\bild \overline{\varphi} = \bild \varphi $. Insbesondere ist $\overline{\varphi} $ ein Isomorphismus, falls
$\varphi$ surjektiv ist und $\Kern \varphi = N$ ist.

\minisec{Beweis}
Wenn $\overline{\varphi} $ existiert, dann ist $\overline{\varphi} (aN) = \overline{\varphi}( \pi (a)) = \varphi(a) $, also ist $\overline{\varphi} $ eindeutig.

Zu Existenz definieren wir $\overline{\varphi}(aN) := \varphi(a) $. zZ: $a_1 \, N = a_2  \, N \Rightarrow \varphi(a_1)= \varphi(a_2)$. Ist $a_1 N = a_2 N$, dann 
$a_2 ^{-1} a_1 \in N \subseteq \Kern \varphi$, also $\varphi(a_1)= \varphi(a_2 a_2 ^{-1} a_1)$. \bewende
% subsection 125 (end)

\subsection[Korollar aus dem Homomorphiesatz]{Korollar} % (fold)
\label{sub:126}
Sei $\varphi : G \to H$ ein Gruppenhomomorphismus. Dann gilt $\bild \varphi \cong \nicefrac{G}{\Kern \varphi}$. 
\minisec{Beweis}
Mit $N = \Kern \varphi$ und $H= \bild \varphi$ folgt dies aus dem Homomorphiesatz.
% subsection 126 (end)

\subsection[Bemerkung: Kurze exakte Folge]{Bemerkung} % (fold)
\label{sub:127}
Seien $N,G$ und $Q$ Gruppen mit Gruppenhomomorphismen 
\[
	\begin{tikzcd}
		N \rar[hook]{i} & G \rar[two heads]{q} & Q
	\end{tikzcd}
	\tag{$\star$}
\]
Gilt $i$ injektiv, $q$ surjektiv,
$\Kern q = \bild i$ so heißt $(\star)$ eine \Index{kurze exakte Folge} (oder Sequenz). Es gilt dann $Q\cong \nicefrac{G}{\Kern q} \cong \nicefrac{G}{\bild i}$.
% subsection 127 (end)

\subsection[Definition: einfache Gruppe]{Definition} % (fold)
\label{sub:128}
Eine Gruppe $G$, die keinen Normalteiler außer $\{e\}$ und $G$ besitzt, heißt \Index{einfach}. 
% subsection 128 (end)

\subsection[Bemerkung: Die einfachen endlichen Gruppen sind vollständig klassifiziert]{Bemerkung} % (fold)
\label{sub:129}
Die einfachen endlichen Gruppen sind vollständig klassifiziert. Es gibt mehrere unendliche Serien solcher Gruppen (zyklische Gruppen von Primzahlordnung, 
alternierende Gruppen $A_n$ für $n \ge 5$, Gruppen vom Lie-Typ) und 26 sporadische einfache Gruppen. Die größte der sporadischen Gruppen hat etwa die Ordnung 
$8 \cdot 10^{53}$ und hat den schönen Namen das \glqq Monster\grqq.
% subsection 129 (end)
% section 1 (end)
\newpage

\section{Zyklische Gruppen} % (fold)
\label{sec:2}

\subsection[Definition: Von $x$ erzeugte Untergruppe]{Definition} % (fold)
\label{sub:21}
Sei $G$ eine Gruppe und $x \in G$. Dann heißt $\langle x \rangle := \set[x^n]{n \in \mathds{Z}} $ die von $x$ \bet{erzeugte Untergruppe}\index{Untergruppe!erzeugte}.
% subsection 21 (end)

\subsection[Bemerkung: Eigenschaften der Untergruppe $\langle x\rangle$]{Bemerkung} % (fold)
\label{sub:22}
\begin{enumerate}[(i)]
	\item $\langle x\rangle$ ist eine Untergruppe $e=x^0, \enspace x^n \cdot x^m = x^{n+m}, \enspace {(x^n)}^{-1}= x^{-n}$
	\item $\langle x\rangle$ ist die kleinste Untergruppe von $G$, die $x$ enthält
	\item $\langle x\rangle$ ist kommutativ: $x^n \cdot x^m = x^m \cdot  x^n$
	\item $\varphi : \mathds{Z} \to G$, $\varphi(n) := x^n$ ist ein Gruppenhomomorphismus mit $\bild \varphi = \langle x\rangle$
\end{enumerate}
% subsection 21 (end)

\subsection[Definition: zyklische Gruppe]{Definition} % (fold)
\label{sub:23}
Eine Gruppe $G$ heißt \bet{zyklisch}\index{Gruppe!zyklische}, falls sie von einem Element erzeugt wird, d.h. falls es $x \in G$ gibt mit $\langle x\rangle =G$.
% subsection 223 (end)

\subsection[Beispiel: Erzeuger von $\mathds{Z}$]{Beispiel} % (fold)
\label{sub:24}
$\mathds{Z}$ wird erzeugt von $1 \in \mathds{Z}$ oder auch $-1 \in \mathds{Z}$, aber nicht von $2 \in \mathds{Z}$.
% subsection 24 (end)

\subsection[Lemma 1: Charakterisierung von zyklisch durch surjektiven Gruppenhomomorphismus]{Lemma 1} % (fold)
\label{sub:25}
Sei $G$ eine Gruppe. Dann ist $G$ genau dann zyklisch, wenn es einen surjektiven Gruppenhomomorphismus $\varphi : \mathds{Z} \twoheadrightarrow G$ gibt.
\minisec{Beweis}
Ist $G=\langle x\rangle$ zyklisch, so ist $\varphi : \mathds{Z} \to G $ mit $\varphi(n) := x^n$ ein surjektiver Gruppenhomomorphismus. Ist $\varphi : \mathds{Z} \to G$ ein surjektiver
Gruppenhomomorphismus so gilt $G = \langle x\rangle$ mit $x := \varphi(1)$
% subsection 24 (end)

\subsection{Satz: Klassifikationssatz für zyklische Gruppen} % (fold)
\label{sub:26}
Sei $G$ eine zyklische Gruppe. Dann ist 
\[
	G \cong \begin{cases}
		\mathds{Z}, &\text{ falls } \abs*{G}= \infty \\
		\nicefrac{\mathds{Z}}{m \mathds{Z}}, &\text{ falls } \abs*{G}=m < \infty
	\end{cases}
\]
\minisec{Beweis}
Sei $G$ eine zyklische Gruppe. Aus Lemma 1 (\ref{sub:25}) folgt: $\exists$ surjektiver Gruppenhomomorphismus $\varphi : \mathds{Z} \twoheadrightarrow G$. Also $G \cong \nicefrac{\mathds{Z}}{\Kern \varphi}$ nach dem Homomorphiesatz (\ref{sub:124}). Aus Lemma 2 (\ref{sub:28}) folgt $\exists m \in \mathds{N}$ mit $\Kern \varphi = m \mathds{Z}$. Ist $m=0$, so folgt
$\Kern \varphi  =\{0\}$, also $G \cong \nicefrac{\mathds{Z}}{\{0\}} \cong \mathds{Z}$. Andernfalls $G \cong \nicefrac{\mathds{Z}}{m \mathds{Z}}$. Ist $m\not= 0$, also 
$ G \cong \nicefrac{\mathds{Z}}{m \mathds{Z}}$ so folgt $m= \abs*{\nicefrac{\mathds{Z}}{m \mathds{Z}}} = \abs*{G}$. \bewende
% subsection 25 (end)

\subsection[Bemerkung: Vergleich von zyklischen Gruppen über die Ordnung]{Bemerkung} % (fold)
\label{sub:27}
Zyklische Gruppen werden durch ihre Ordnung klassifiziert: Für zyklische Gruppen $H$ und $G$ gilt: 
\[
	G \cong H \Leftrightarrow \abs*{G}= \abs*{H}  
\]
% subsection 26 (end)

\subsection[Lemma 2: Jede Untergruppe von $\mathds{Z}$ ist zyklisch]{Lemma 2} % (fold)
\label{sub:28}
Sei $U \subseteq  \mathds{Z}$ eine Untergruppe. Dann gibt es $m \in \mathds{N}$ mit $U = m \mathds{Z} = \langle m\rangle$. Insbesondere ist $U$ zyklisch.
\minisec{Beweis}
Ist $U= \{0\}$, so setze $m := 0$. Sei also $U \not= \{0\}$. Dann gibt es $n \in U, n \not= 0$. Nun ist $n>0$ oder $-n>0$. $U$ enthält also mindestens ein $n$ mit $n>0$. 
Sei $m$ das minimale positive Element in $U$. 

Behauptung: $U= m \mathds{Z}$. Wegen $m \in U$ ist $m \mathds{Z}= \langle m\rangle \subseteq U$. Sei $a \in U$. Division mit Rest ergibt $a= q \cdot m + r$ mit 
$q \in \mathds{Z}, r \in \set{0, \ldots , m-1} $. Es folgt 
\[
	r= \underbrace{a}_{\in U}- \underbrace{q \cdot m}_{\in U} \in U
\]
Da $r < m$ und $m \in U$ das minimale positive Element war, folgt $r=0$. Also $a= q \cdot m \in U$. \bewende
% subsection 27 (end)

\subsection[Proposition: Jede Untergruppe einer zyklischen Gruppe ist zyklisch]{Proposition} % (fold)
\label{sub:29}
Jede Untergruppe einer zyklischen Gruppe ist zyklisch.
\minisec{Beweis}
Sei $G$ zyklisch und $U \subseteq G$. Mit Lemma 1 (\ref{sub:25}) folgt: $\exists \varphi : \mathds{Z} \twoheadrightarrow G$. Dann ist $U' := \varphi ^{-1}(U)$ eine Untergruppe von $\mathds{Z}$. Lemma
2 besagt $U'= m \mathds{Z}$. Insbesondere ist $U'$ zyklisch. Damit ist auch $U = \varphi(U')$ zyklisch. \bewende
% subsection 28 (end)

\subsection[Definition: Ordnung eines Gruppenelements]{Definition} % (fold)
\label{sub:210}
Sei $G$ eine Gruppe und $a \in G$. Dann heißt $\ord(a) := \abs*{\langle a\rangle} $ die \Index{Ordnung} von $a$.
% subsection 29 (end)

\subsection[Bemerkung, wann $\ord(a)$ endlich ist]{Bemerkung} % (fold)
\label{sub:211}
$\ord (a)$ ist genau dann endlich, wenn es ein positives $n \in \mathds{N}$ gibt mit $a^n = e$. In diesem Fall ist $\ord (a)$ die kleinste positive Zahl $n$ mit $a^n=e$.
% subsection 210 (end)

\subsection[Satz: Ordnung der Gruppen wird von $\ord(a)$ geteilt]{Satz} % (fold)
\label{sub:212}
Sei $G$ eine endliche Gruppe und $a \in G$. Dann teilt $\ord(a)$ die Ordnung von $G$.
\minisec{Beweis}
Satz von Lagrange mit $U:= \langle a\rangle$. \bewende
% subsection 211 (end)

\subsection[Satz über Gruppen von Primzahlordnung]{Satz} % (fold)
\label{sub:213}
Sei $G$ eine Gruppe von Primzahlordnung. Dann gilt
\begin{enumerate}[i)]
	\item $G$ ist zyklisch. Insbesondere $G \cong \nicefrac{\mathds{Z}}{p \mathds{Z}}$ mit $p = \abs*{G} $
	\item $\{e\}$ und $G$ sind die einzigen Untergruppen von $G$. Insbesondere ist $G$ einfach.
\end{enumerate}
\minisec{Beweis}
Sei $a \in G \setminus \{e\}$. Da die Ordnung von $a$ die Ordnung von $G$ teilt, ist $\ord(a) \in \set{1, \abs*{G} } $. Da $a \not= e$ ist, ist $\ord(a) \not=1$. Also 
$\ord(a)= \abs*{G} $. Also $\langle a\rangle =G$.\\
ii) folgt mit Lagrange \bewende
% subsection 212 (end)
% section 2 (end)
\newpage

\section{Gruppenwirkungen} % (fold)
\label{sec:3}

\subsection[Definition: Wirkung]{Definition} % (fold)
\label{sub:31}
Sei $G$ eine Gruppe und $X$ eine Menge. Eine \Index{Wirkung} von $G$ auf $X$ ist eine Abbildung $G \times X \to X$, $(g,x) \mapsto g \cdot x$ die folgende Axiome erfüllt:
\begin{enumerate}[(i)]
	\item $\forall x \in X : e \cdot x = x$
	\item $\forall g,h \in G, \forall x \in X : g \cdot (h \cdot x) = (g \cdot h) \cdot x$
\end{enumerate}
% subsection 31 (end)

\subsection[Bemerkung: Anderer Namen für Wirkungen]{Bemerkung} % (fold)
\label{sub:32}
Oft spricht man auch von einer \Index{Operation} oder \Index{Aktion}.
% subsection 32 (end)

\subsection[Bemerkung zu Wirkungen und Gruppenhomomorphismen]{Bemerkung} % (fold)
\label{sub:33}
Sei $G \times X \to X$, $(g,x) \mapsto g x$ eine Wirkung. Für $g \in G$ sei $\psi_g : X \to X, \psi_g(x) := g x$. Dann definiert $g \mapsto \psi_g$ einen 
Gruppenhomomorphismus \marginnote{siehe \ref{sub:15} \ref{sym_grp})} $G \to S_X$. ($\psi_g \in S_X$ da $\psi_{g ^{-1}} = (\psi_g) ^{-1}$). Sei umgekehrt $\alpha : G \to S_X$ ein Gruppenhomomorphismus. Dann erhalten wir 
eine Wirkung $G \times X \to X$ durch $g \cdot x := \underbrace{\alpha(g)}_{\in S_X} (x)$.
% subsection 33 (end)

\subsection[Beispiele für Wirkungen]{Beispiele} % (fold)
\label{sub:34}
\begin{enumerate}[i)]
	\item $S_X$ wirkt auf $X$ durch $\sigma \cdot x = \sigma(x)$ für $\sigma \in S_X, x \in X$
	\item Die Symmetriegruppe $S_{\square}$ des Quadrats wirkt kanonisch auf
	\begin{itemize}
		\item der Menge der Punkte des Quadrats
		\item der Menge der Ecken des Quadrats
		\item der Menge der Kanten des Quadrats
		\item der Menge der Paare von Kanten
	\end{itemize}
	\item Sei $G$ eine Gruppe. Dann ist $G \times G \to G$, $(g,x) \mapsto g \cdot_G x$ eine Wirkung von $G$ auf $X:= G$. Dies ist die 
	\Index{Linkstranslationswirkung}. 
	
	$g \cdot_R x := x g$, dann $g \cdot_R ( h \cdot_R x) = g \cdot_R (x h) = (x h) g = x \cdot ( h \cdot g) = (h \cdot g) \cdot_R x$ \\lässt sich durch Inverse kurieren.
	
	Die \Index{Rechtstranslationswirkung} ist definiert durch $G \times G \to G$, $(g,x) \mapsto x \cdot g ^{-1} $
	\item \label{konj_wirk} Die \Index{Konjugationswirkung} von $G$ auf $X:= G$ ist gegeben durch $G \times G \to G$, $(g,x) \mapsto g \cdot x \cdot g ^{-1}$.
	\item Sei $U \subseteq G$ eine Untergruppe. Dann wirkt $G$ durch Linkstranslation auf den Linksnebenklassen $\nicefrac{G}{U}$. 
	\[
		G \times \nicefrac{G}{U} \to \nicefrac{G}{U} \enspace, \quad (g,aU) \mapsto g a U
	\]
\end{enumerate}
% subsection 34 (end)

\subsection[Definition: Bahn und Standgruppe]{Definition} % (fold)
\label{sub:35}
Sei $G \times X \to X$ eine Wirkung. Sei $x \in X$. 
\begin{enumerate}[i)]
	\item $G x := \set[g \cdot x]{g \in G} \subseteq X $ heißt die \Index{Bahn} oder der \Index{Orbit} von $x$ unter $G$.
	\item $G_x := \set[g]{g \cdot x = x} \subseteq G $ heißt die \Index{Standgruppe} oder \Index{Isotopiegruppe} von $x$ unter $G$. 
\end{enumerate}
% subsection 35 (end)

\subsection[Bemerkung: $G_x$ ist eine Untergruppe von $G$]{Bemerkung} % (fold)
\label{sub:36}
$G_x$ ist eine Untergruppe von $G$, wie man sich leicht überlegt.
% subsection 36 (end)

\subsection[Lemma über die Gleichheit von Bahnen]{Lemma} % (fold)
\label{sub:37}
Sei $G \times X \to X$ eine Wirkung. Seien $x,y \in X$. Dann gilt $Gx= Gy \Longleftrightarrow Gx \cap Gy \not= \emptyset$
\minisec{Beweis}
\begin{description}
	\item[\glqq$\Rightarrow $\grqq] Ist klar, da $x \in Gx$, also $Gx \not= \emptyset$
	\item[\glqq$\Leftarrow$\grqq] Sei $z \in Gx \cap Gy$. Dann $z= gx= hy$ mit $g,h \in G$. Also $(h ^{-1} g) \cdot x = y$. Dann gilt
	\[
		G \cdot y = G (h ^{-1} g) \cdot x \subseteq G \cdot x
	\]
	Genauso $Gx \subseteq Gy$. \bewende
\end{description}
% subsection 37 (end)

\subsection[Korollar: $X$ ist die disjunkte Vereinigung der Bahnen]{Korollar} % (fold)
\label{sub:38}
$X$ ist die disjunkte Vereinigung der Bahnen.
% subsection 38 (end)

\subsection[Lemma: Bijektion $\nicefrac{G}{G_x} \to Gx$]{Lemma} % (fold)
\label{sub:39}
Sei $G \times X \to X$ eine Wirkung. Sei $x \in X$. Dann definiert 
\[
	\begin{tikzcd}
		\nicefrac{G}{G_x} \rar{\varphi} & Gx
	\end{tikzcd} \quad , \qquad  g G_x \mapsto g \cdot x
\]
eine Bijektion.
\minisec{Beweis}
\begin{description}
	\item[$\varphi$ ist wohldefiniert:] Sei $gG_x = hG_x$. Dann $h ^{-1} g \in G_x$. Also $g \cdot x = h (h ^{-1} g) \cdot x = h \cdot x$
	\item[$\varphi$ surjektiv:] nach Definition der Bahn $Gx$.
	\item[$\varphi$ injektiv:] Sei $\varphi(gG_x) = \varphi(h G_x)$. Also $g \cdot x = h \cdot x$. Dann $h ^{-1} \cdot g \cdot x = x$, also $h ^{-1} g \in G_x$. Daher
	$hG_x= gG_x$. \bewende
\end{description}
% subsection 39 (end)

\subsection{Satz (Bahnengleichung)} % (fold)
\label{sub:310}\hypertarget{bahnengleichung}{ }
Sei $G \times X \to X$ eine Wirkung auf einer endlichen Menge $X$. Sei $x_1, \ldots ,x_n$ ein \Index{Vertretersystem} der Bahnen, d.h. jede Bahn enthält genau ein $x_i$. Dann gilt
\[
	\abs*{X} = \sum_{i=1}^{n} \abs*{Gx_i}= \sum_{i=1}^{n} \left[G:G_{x_i}\right]  
\]
\minisec{Beweis}
Da $X$ die disjunkte Vereinigung der Bahnen $Gx_1, \ldots , Gx_n$ ist, gilt $\abs*{X} = \sum_{i=1}^{n} \abs*{Gx_i}  $. Da es eine Bijektion $\nicefrac{G}{G_{x_i}} \to Gx_i$
gibt, ist $\abs*{Gx_i} = \abs*{\nicefrac{G}{G_{x_i}}} = \left[ G : G_{x_i} \right]  $. \bewende
% subsection 310 (end)

\subsection[Definition: Zentralisator, Zentrum und Normalisator]{Definition} % (fold)
\label{sub:311}
Sei $S \subseteq G$ eine Teilmenge. Dann heißt 
\[
	\begin{array}{rll}
	Z_S &:= \set[g \in G]{g \cdot s=s \cdot g \enspace\forall s \in S} & \quad \text{der \Index{Zentralisator} von $S$}\\
	Z_G &:= \set[g \in G]{g \cdot h = h \cdot g \enspace\forall h \in G} & \quad \text{das \Index{Zentrum} der Gruppe $G$}\\
	N_S &:= \set[g \in G]{gS=Sg} & \quad \text{der \Index{Normalisator} von $S$ in $G$}
\end{array}
\]
% subsection 311 (end)

\subsection[Bemerkung: Eigenschaften von von Zentralisator, Normalisator und Zentrum]{Bemerkung} % (fold)
\label{sub:312}
\begin{enumerate}[(i)]
	\item $Z_S$ und $N_S$ sind Untergruppen von $G$
	\item $Z_G$ ist ein Normalteiler.
	\item $Z_G$ ist abelsch.
	\item Ist $U$ eine Untergruppe, so ist $U \subseteq N_U$ ein Normalteiler. Genauso ist $N_U$ die größte Untergruppe von $G$, die $U$ als Normalteiler enthält.
\end{enumerate}
% subsection 312 (end)

\subsection[Bemerkung um den nächsten Satz besser zu verstehen]{Bemerkung} % (fold)
\label{sub:313}
$g \cdot s=s \cdot g \iff s= g \cdot s \cdot g ^{-1}$ \marginnote{Wer hätte das gedacht\ldots}
% subsection 313 (end)

\subsection{Satz (Klassengleichung)} % (fold)
\label{sub:314}
Sei $G$ eine endliche Gruppe und $g_1, \ldots , g_n$ ein Vertretersystem für die Konjugationswirkung \marginnote{siehe \ref{sub:34} \ref{konj_wirk})} auf $G \setminus Z_G$. Dann gilt 
\[
	\abs*{G} = \abs*{Z_G} + \sum_{i=1}^{n} \left[ G: Z_{\{g_i\}}\right]  
\]
\minisec{Beweis}
Die \hyperlink{bahnengleichung}{Bahnengleichung (\ref{sub:310})} für die Konjugationswirkung von $G$ auf $G \setminus Z_G$ liefert
\[
	\abs*{G} - \abs*{Z_G} =   \abs*{G \setminus Z_G} = \sum_{i=1}^{n} \left[ G : Z_{\{g_i\}}\right] \bewende
\]
% subsection 314 (end)

\subsection[Definition: $p$-Gruppe]{Definition} % (fold)
\label{sub:315}
Sei $p$ eine Primzahl. Eine \bet{$p$-Gruppe}\index{p-Gruppe@$p$-Gruppe} ist eine endliche Gruppe, deren Ordnung eine Potenz von $p$ ist.
% subsection 316 (end)

\subsection[Beispiel: $p$-Gruppen]{Beispiel} % (fold)
\label{sub:316}
\begin{enumerate}[(i)]
	\item $\nicefrac{\mathds{Z}}{4 \mathds{Z}}$, $\nicefrac{\mathds{Z}}{3 \mathds{Z}} \times \nicefrac{\mathds{Z}}{3 \mathds{Z}}$
	\[
		H_{\nicefrac{\mathds{Z}}{2 \mathds{Z}}} = \set[{\begin{pmatrix}
			1 & a & c \\
			0 & 1 & b \\
			0 & 0 & 1
		\end{pmatrix}}]{a,b,c \in \nicefrac{\mathds{Z}}{2 \mathds{Z}}} 
	\]
	$\abs*{H_{\nicefrac{\mathds{Z}}{2 \mathds{Z}}}} = 2^3 = 8 $
\end{enumerate}
% subsection 317 (end)

\subsection[Korollar: $p$-Gruppen haben ein nichttriviales Zentrum]{Korollar} % (fold)
\label{sub:317}
Sei $P$ eine $p$-Gruppe, dann folgt aus der Klassengleichung, dass das Zentrum von $P$ nichttrivial ist.
\minisec{Beweis}
Sei $p^m = \abs*{P} $. Mit der Klassengleichung (\ref{sub:314}) folgt für ein Vertretersystem\index{Vertretersystem} $x_i$, $i=1,\ldots ,n$:
\[
	p^m = \abs*{P} = \abs*{Z_P} + \sum_{i=1}^{n} \left[P : Z_{x_i}\right]  
\]
Da die $x_i$ nicht im Zentrum von $P$ liegen, ist $Z_{x_i} \not= P$ $\forall x_i$. Daher teilt $p$ den Index $[P:Z_{x_i}] = \frac{\abs*{P} }{\abs*{Z_{x_i}}}$. Also gilt: $p$ teilt
auch $\abs*{Z_P} $, also $\abs*{Z_P} >1 \Rightarrow $ $Z_P$ ist nicht $\set{e} $. \bewende
% subsection 318 (end)
% section 3 (end)
\newpage

\section{Sylow-Gruppen} % (fold)
\label{sec:4}

\subsection[Definition: $p$-Sylow-Gruppe]{Definition} % (fold)
\label{sub:41}
Sei $G$ eine endliche Gruppe. Eine Untergruppe $P$ von $G$ heißt \bet{$p$-Sylow-Gruppe}\index{p-Sylow-Gruppe@$p$-Sylow-Gruppe}, wenn $P$ eine $p$-Gruppe ist und $[G:P]$ teilerfremd zu $p$ ist. 
% subsection 41 (end)

\subsection[Bemerkung über die Ordnung von $p$-Sylowgruppen]{Bemerkung} % (fold)
\label{sub:42}
Ist $\abs*{G}= p^n \cdot l $ mit $p \nmid l$, so hat jede $p$-Sylow-Gruppe die Ordnung $p^n$.
% subsection 42 (end)

\subsection[Beispiele für $p$-Sylowgruppen]{Beispiele} % (fold)
\label{sub:43}
\begin{enumerate}[(i)]
	\item $\nicefrac{\mathds{Z}}{12 \mathds{Z}}$ enthält eine $2$-Sylowgruppe.
	\item $S_3$ enthält drei $2$-Sylowgruppen
	\[
		\set{e, (1 2)}, \set{e,(1 3)}, \set{e, (2 3)}   
	\]
	und eine $3$-Sylowgruppe
	\[
		\set{e, (1 2 3), (1 3 2)} 
	\]
\end{enumerate}
% subsection 43 (end)

\subsection{Satz (Sylow)} % (fold)
\label{sub:44}
Sei $G$ eine endliche Gruppe und $p$ eine Primzahl. Dann gelten folgende 3 Aussagen
\begin{enumerate}[(1)]
	\item Zu jeder $p$-Untergruppe $U$ von $G$ gibt es mindestens eine $p$-Sylow-Gruppe $P$ mit $U \subseteq P$. Insbesondere gibt es mindestens eine $p$-Sylow-Gruppe.
	\item Sind $P$ und $Q$ zwei $p$-Sylow-Gruppen, dann gibt es ein $g \in G$ mit $g P g ^{-1} =Q$. Je zwei $p$-Sylow-Gruppen sind also konjugiert zueinander.
	\item Für die Anzahl $n_p$ der $p$-Sylow-Gruppen gilt: $n_p \mid \abs*{G} $ und $n_p \equiv 1 \mod p$
\end{enumerate}
Beweis im Internet unter: \\ \url{http://wwwmath.uni-muenster.de/reine/u/topos/lehre/WS2013-2014/Algebra/algebra.html}
% subsection 44 (end)

\subsection[Bemerkung zur Bezeichnung der Sylow-Sätze]{Bemerkung} % (fold)
\label{sub:45}
Aussage $(k)$ heißt auch $k$-ter Sylowsatz.
% subsection 45 (end)

\subsection[Korollar: Eine $p$-Sylowgruppen ist Normalteiler gdw. sie die einzige $p$-Sylowgruppe ist]{Korollar aus (2)} % (fold)
\label{sub:46}
Eine $p$-Sylow-Gruppe in $G$ ist ein Normalteiler genau dann, wenn sie die einzige $p$-Sylow-Gruppe von $G$ ist.
\minisec{Beweis}
Sei $P$ eine $p$-Sylow-Gruppe und $g \in G$. Dann ist $g P g ^{-1}$ eine weitere $p$-Sylow-Gruppe in $G$. Wenn $P$ die einzige $p$-Sylowgruppe ist, gilt also $gP g ^{-1}=P$.
Für die andere Richtung seien $P_1$ und $P_2$ $p$-Sylowgruppen in $G$. Sei $P_1$ ein Normalteiler von $G$. Nach (2) gibt es ein $g \in G$ mit $P_2 = g P_1 g ^{-1} = P_1$, 
da $P_1$ Normalteiler. \bewende 
% subsection 46 (end)

\subsection[Satz: Isomorphie einer Gruppe der Ordnung $p \cdot q$, $p,q$ prim]{Satz} % (fold)
\label{sub:47}
Seien $p$ und $q$ Primzahlen mit $p < q$ und $p \nmid (q-1)$. Dann ist jede Gruppe der Ordnung $p \cdot q$ isomorph zu 
\[
	\nicefrac{\mathds{Z}}{p \mathds{Z}} \times \nicefrac{\mathds{Z}}{q \mathds{Z}}
\]
\minisec{Beweis}
Sei $G$ die Gruppe mit Ordnung $p \cdot q$. Sei $n_p$ die Anzahl der $p$-Sylow-Gruppen. Wir wissen:
\begin{multicols}{2}
	\begin{enumerate}[(a)]
		\item $n_p \mid p \cdot q$ 
		\item $n_p \equiv 1 \mod p$
	\end{enumerate}
\end{multicols}
Da $p$ und $q$ prim sind, folgt $n_p \mid p$ oder $n_p \mid q$. Wegen (b) ist $n_p \mid p$ ausgeschlossen, also gilt $n_p \in \set{1,q} $. Aber $n_p =q$ ist unmöglich, denn
$q \equiv 1 \mod p \Leftrightarrow p \mid q-1$ \light zur Annahme. \\
Es gibt also eine normale $p$-Sylow-Gruppe $S_p$. Sei $n_q$ die Anzahl der $q$-Sylow-Gruppen
\begin{multicols}{2}
	\begin{enumerate}[(a')]
		\item $n_q \mid p \cdot q$
		\item $n_q \equiv 1 \mod q$
	\end{enumerate}
\end{multicols}
Wie oben folgt $n_q \in \set{1,p} $. \\
Wegen $q > p$ ist $p \not\equiv 1 \mod q$. Daher bleibt nur $n_q = 1$. Folglich gibt es eine $q$-Sylow-Gruppe $S_q$ als Normalteiler in $G$. Außerdem ist 
$S_q \cap S_p = \set{e}$. Betrachte nun 
\[
	\varphi : S_p \times S_q \to G \quad , \quad (a,b) \mapsto a \cdot b
\]
Für $a \in S_p$ und $b \in S_q$ gilt:
\begin{align*}
	a \cdot \underbrace{(b \cdot a ^{-1} \cdot b ^{-1})}_{\in S_p} \in S_p \quad &\text{und auch} \quad  (a \cdot b \cdot a ^{-1}) \cdot b ^{-1} \in S_q  \\
	\Rightarrow & a \cdot b \cdot a ^{-1} \cdot b ^{-1} \in S_p \cap S_q \\
	\Rightarrow & a \cdot b \cdot a ^{-1} \cdot b ^{-1} = e \\
	\Rightarrow  &ab=ba
\end{align*}

$\Rightarrow \varphi$ ist ein Gruppenhomomorphismus.\\
Wegen Primzahlordnung gilt $S_p \cong \nicefrac{\mathds{Z}}{p \mathds{Z}}$, $S_q \cong \nicefrac{\mathds{Z}}{q \mathds{Z}}$. Sei $a \in S_p$ und $b \in S_q$ mit $\varphi(a,b)=ab =e$
\[
	a b = e \Leftrightarrow \underbrace{a}_{\in S_q} = \underbrace{b ^{-1} }_{\in S_p} \Rightarrow  a,b \in S_q \cap S_q
\]
folglich gilt $a=b=e$. $\Rightarrow \varphi$ ist injektiv. Da $\abs*{S_p \times S_q} = \abs*{G}= p \cdot q  $, ist $\varphi$ auch surjektiv. \bewende
% subsection 47 (end)
% section 4 (end)
\newpage
\section{Polynome} % (fold)
\label{sec:5}

\subsection[Definition: Polynom, Grad, Leitkoeffizient]{Definition} % (fold)
\label{sub:51}
Sei $R$ ein kommutativer Ring mit $1$. Ein formaler Ausdruck
\[
	p = a_n X^n + a_{n-1} X^{n-1} + \ldots + a_1 X + a_0
\]
mit $a_n, \ldots , a_0 \in R$ heißt \Index{Polynom} mit Koeffizienten in $R$. Die Menge aller Polynome wird mit $R[X]$ bezeichnet. Der \Index{Grad} $d(p)$
von $p$ ist gegeben durch $d(p) = \max \set[i]{a_i \not=0}$. Das Element $a_{d(p)} \in R$ heißt \Index{Leitkoeffizient} von $p$. Für das Nullpolynom $0$ setzen wir
$d(0) := - \infty$ und $l(0)=0$. \\
Ein Polynom mit $a_i = 0$ für $i \ge 1$ heißt \bet{konstant}\index{Polynom!konstantes}. Ist $l(p)=1$, dann heißt $p$ \bet{normiert}\index{Polynom!normiertes}. 
% subsection 51 (end)

\subsection[Bemerkung: $X$ ist nur formale Variable!]{Bemerkung} % (fold)
\label{sub:52}
$X$ ist eine formale Variable. Es gilt $0 \cdot X = 0$, aber $X$ ist kein Element von $R$. Jedes $p \in R[X]$ definiert eine Abbildung $f_p : R \to R$, 
$\lambda  \mapsto p(\lambda )$.
% subsection 52 (end)

\subsection[Beispiele zum Unterschied zwischen einen Polynom und zugehöriger Abbildung]{Beispiele} % (fold)
\label{sub:53}
\begin{enumerate}[(i)]
	\item $p= X^2+X$ als Element von $\mathds{F}_2[X]$. In diesem Fall ist $f_p =0$, aber $p \not= 0$.
	\item $p=2X$ als Element von $\mathds{F}_2[X]$ ist das Nullpolynom.
	\item $p=2X$ als Element von $\mathds{Z}[X]$ ist nicht das Nullpolynom.
\end{enumerate}
% subsection 53 (end)

\subsection[Bemerkung: Addition und Multiplikation auf {$R[X]$}]{Bemerkung} % (fold)
\label{sub:54}
Auf $R[X]$ lassen sich eine Addition und eine Multiplikation erklären:
\begin{align*}
	&\enbrace*{\sum_{k=0}^{n} a_k \cdot X^k} + \enbrace*{\sum_{k=0}^{n} b_k \cdot X^k} := \sum_{k=0}^{n} (a_k + b_k) \cdot X^k  \\
	& \enbrace*{\sum_{k=0}^{m} a_k \cdot X^k} \cdot \enbrace*{\sum_{l=0}^{n} b_l \cdot X^l} :=  \sum_{j=0}^{m+n} \enbrace*{\sum_{k+l=j} a_k \cdot b_l } \cdot X^j 
\end{align*}
Das Konstante Polynom $1$ ist ein Einselement für $\cdot $. Das konstante Polynom $0$ ist ein Nullelement für $+$.
% subsection 54 (end)

\subsection[Definition: Nullteiler, Integritätsring]{Definition} % (fold)
\label{sub:55}
Eine Element $a \in R$ heißt \Index{Nullteiler}, falls $b \in R\setminus \{0\}$ existiert mit $a \cdot b = 0$ (oder $b \cdot a=0$). Falls $R$ nur $0$ als Nullteiler besitzt,
heißt $R$ \Index{nullteilerfrei} oder \Index{Integritätsring}.   
% subsection 55 (end)

\subsection[Beispiele zu Nullteilern]{Beispiele} % (fold)
\label{sub:56}
\begin{enumerate}[(i)]
	\item Jeder Körper ist nullteilerfrei.
	\item $2$ in $\nicefrac{\mathds{Z}}{4 \mathds{Z}}$ ist ein Nullteiler, denn $2 \cdot 2 =4$ und $4 \equiv 0 \mod 4$.
\end{enumerate}
% subsection 56 (end)

\subsection[Satz: Abschätzungen für den Grad eines Polynoms]{Satz} % (fold)
\label{sub:57}
Sei $R$ ein kommutativer Ring mit $1$. Seien $p,q \in R[X]$. Dann gilt:
\begin{enumerate}[(a)]
	\item $d(p+q) \le \max \set{d(p), d(q)} $
	\item $d(p \cdot q) \le d(p)+ d(q)$
\end{enumerate}
Ist $R$ nullteilerfrei, dann ist $d(p \cdot q) = d(p)+ d(q)$. Außerdem ist $l(p \cdot q)= l(p) \cdot l(q)$.
\minisec{Beweis}
Für $p=0$ oder $q=0$ ist die Aussage klar. \\ $p = \sum_i a_i X^i $, $q = \sum_j b_j X^j$. Dann ist $a_i + b_i = 0$ für $i > \max \set{d(p), d(q)} $. Hieraus folgt (a).
Außerdem ist $\sum_{k+l=j} a_k \cdot b_l= 0$ für $j\ge d(p)+ d(q)$. Für $j= d(p)+ d(q)$ erhalten wir
\[
	\sum_{k+l=j} a_k \cdot b_l = l(p) \cdot l(q)
\]
% subsection 57 (end)

\subsection{Division mit Rest} % (fold)
\label{sub:58}
Sei $K$ ein Körper und seien $p,q \in K[X]$ mit $q \not= 0$. Dann gibt es genau eine Zerlegung der Form $p= s \cdot q + r$ mit $d(r)< d(q)$.
\minisec{Beweis}
\begin{description}
	\item[Eindeutigkeit:] Sei $p = s_0 q + r_0 = s_1 q + r_1$ mit $d(r_i)< d(q)$ für $i \in \{0,1\}$.
	\[
		(s_o -s_1) q = r_1 -r_0 \enspace \Rightarrow \underbrace{d\big((s_0 -s_1) q\big)}_{d(s_0 -s_1) + d(q)} = d(r_1 -r_0) \le \max \set{d(r_0), d(r_1)}< d(q) 
	\]
	$\Rightarrow d(s_1 -s_1) < 0$ $\Rightarrow d(s_1 -s_1) = - \infty$ $\Rightarrow s_0 - s_1 = 0$, also auch $r_0 = r_1$.
	\item[Existenz] Für $d(p) < d(q)$ setzen wir $s=0$ und $r=p$. 
	
	Sei jetzt also $d(p) \ge d(q)$: Sei $n := d(q)$, sei $p= a X^{n+k}+ \ldots $ und $q = b \cdot X^n + \ldots $. Das Polynom $p-ab ^{-1} X^k q$ hat einen kleineren Grad als
	$p$. Nach Induktionsvorraussetzung gibt es $s,r  \in K[X]$ mit $d(r) < d(q)$ und 
	\[
		p - ab ^{-1} X^k q = s \cdot q +r \enspace\Longleftrightarrow \enspace p = (s+ ab ^{-1} X^k) q + r \bewende
	\]
\end{description}
% subsection 58 (end)

\subsection[Bemerkung: Division mit Rest funktioniert auch, wenn Leitkoeffizient Einheit ist]{Bemerkung} % (fold)
\label{sub:59}
Der Satz ist auch richtig für einen Polynomring\marginnote{$R$ kommutativer Ring mit 1} $R[X]$, falls der Leitkoeffizient $l(p) \in R$ eine Einheit ist, d.h. $\exists x \in R$ mit $x \cdot l(q) = 1$.

\uline{Beispiel:}
\marginnote{die Rechnung spare ich mir hier mal \ldots}
\begin{align*}
	(X^5 + 3X^4 + X^3 - 6X^2 + 2X +2) : (X^3 +2X^2 + X -1) = X^2 +X -2
\end{align*}
% subsection 59 (end)

\subsection[Korollar: Bei Division mit $(X-\alpha)$ ist der Rest konstant]{Korollar} % (fold)
\label{sub:510}
Ist $\alpha \in K$ eine Nullstelle von $p \in K[X]$ $(p(\alpha)= f_p(\alpha) =0)$. Dann gilt $p= q \cdot (X -\alpha)$ mit $q \in K[X]$.
\minisec{Beweis}
\emph{Übung}
% subsection 510 (end)

\subsection[Definition: Algebraisch abgeschlossener Körper]{Definition} % (fold)
\label{sub:511}
Ein Körper heißt \Index{algebraisch abgeschlossen}, falls jedes nichtkonstante Polynom $p \in K[X]$ eine Nullstelle in $K$ hat. Dann folgt, dass jedes Polynom $p \in K[X]$
in ein Produkt von Linearfaktoren zerfällt:
\[
	p = l(p) (X- \alpha_1) (X- \alpha_2) \ldots (X- \alpha_n)
\]
Beispiele: 
\begin{itemize}
	\item $\mathds{C}$ ist algebraisch abgeschlossen.
	\item $\mathds{R}$ ist nicht algebraisch abgeschlossen: Betrachte $X^2+1$ 
\end{itemize}
% subsection 511 (end)
% section 5 (end)
\newpage
\section{Ideale und Hauptidealringe} % (fold)
\label{sec:6}
In Abschnitt \ref{sec:6} seien alle Ringe kommutativ mit Eins
\subsection[Definition: Ideal]{Definition} % (fold)
\label{sub:61}
Sei $R$ ein Ring. Eine Teilmenge $I \subseteq R$ heißt \Index{Ideal}, falls gilt: 
\begin{enumerate}[(i)]
	\item $I \subseteq R$ ist eine Untergruppe bezüglich $+$
	\item $\forall a \in I : \forall r \in R : r \cdot a \in I$
\end{enumerate} 
% subsection 61 (end)

\subsection[Beispiele für Ideale]{Beispiel} % (fold)
\label{sub:62}
\begin{enumerate}[(i)]
	\item Für $N \in \mathds{N}$ ist $N \cdot \mathds{Z}$ ein Ideal in $\mathds{Z}$.
	\item $\set[{a_n X^n + a_{n-1}X^{n-1} + \ldots + a_1 X}]{a_n, \ldots , a_1 \in R} \subseteq R[X]$ ist ein Ideal.
	\item Sei $a \in R$. Dann ist $(a) := \set[r \cdot a]{r \in R} $ ein Ideal. Ideale dieser Form heißen \Index{Hauptideale}.
	\item Seien $a_1, \ldots , a_n \in R$. Dann ist $(a_1, \ldots ,a_n) := \set[r_1 a_1 + \ldots + r_n a_n]{r_1, \ldots , r_n \in R} $ ein Ideal.
	\item Sei $\psi : R \to S$ ein Ringhomomorphismus. Dann ist $\Kern \psi$ ein Ideal.
\end{enumerate}
% subsection 62 (end)

\subsection[Bemerkung: Die Faktorgruppe $\nicefrac{R}{I}$ ist ein Faktorring]{Bemerkung} % (fold)
\label{sub:63}
Sei $I \subseteq R$ ein Ideal. Dann wird die Faktorgruppe $\nicefrac{R}{I}$ (bezüglich $+$) zu einem Ring durch $(r+I)(s+I) := r \cdot s + I$. Dies ist wohldefiniert:
Seien $r+I = r' +I$ und $s+I = s' +I$. Dann $r-r', s-s' \in I$. 
\[
	r' s' +I = \enbrace*{r+(r'-r)} \enbrace*{s+(s'-s)}+I = r s + \underbrace{r(s'-s)}_{\in I}+ \underbrace{(r'-r)s  }_{\in I} + \underbrace{(r'-r)(s'-s)}_{\in I}+I 
	= r \cdot s +I
\]
$\nicefrac{R}{I}$ heißt der \Index{Faktorring} oder der \Index{Quotientenring}.
% subsection 63 (end)

\subsection[Bemerkung: Die Quotientenabbildung ist ein Ringhomomorphismus]{Bemerkung} % (fold)
\label{sub:64}
Sei $I \subseteq R$ ein Ideal. Dann ist $\pi : R \to \nicefrac{R}{I}$ mit $\pi(r) := r + I$ ein Ringhomomorphismus mit $\Kern \pi = I$.
% subsection 64 (end)

\subsection{Homomorphiesatz} % (fold)
\label{sub:65}
\begin{wrapfigure}[5]{R}{1cm}
	\centering{\begin{tikzcd}
		R \rar{\psi} \dar[swap]{\pi } & S \\
		\nicefrac{R}{I} \urar[dashed, swap]{\overline{\psi}} &
	\end{tikzcd}}
\end{wrapfigure}
Sei $\psi : R \to S$ ein Ringhomomorphismus und $I \subseteq R$ ein Ideal. Gilt $I \subseteq \Kern \psi$ so gibt es einen eindeutigen Ringhomomorphismus
$\overline{\psi} : \nicefrac{R}{I} \to S $ mit $\psi = \overline{\psi} \circ \pi  $.

Es gilt $\Kern \overline{\psi} = \pi ( \Kern \psi) $ und $\bild \overline{\psi} = \bild \psi $. Insbesondere ist $\overline{\psi} $ ein Isomorphismus, falls $\psi$ surjektiv
und $\Kern \psi = I$ ist. Es ist $ \overline{\psi}(r+ I) = \psi(r)$.
\minisec{Beweis}
Übung!
% subsection 65 (end)

\subsection[Korollar aus dem Homomorphiesatz]{Korollar} % (fold)
\label{sub:66}
Ist $\psi : R \to S$ ein Ringhomomorphismus, so gilt $\bild \psi \cong \nicefrac{R}{\Kern \psi}$. Ist $\psi$ surjektiv, so gilt $S \cong \nicefrac{R}{\Kern \psi}$.
% subsection 66 (end)

\subsection[Definition: Hauptidealring]{Definition} % (fold)
\label{sub:67}
Ein \Index{Hauptidealring} ist ein nullteilerfreier Ring, in dem jedes Ideal ein Hauptideal ist.
% subsection 67 (end)

\subsection[Beispiel eines Hauptidealringes]{Beispiel} % (fold)
\label{sub:68}
$\mathds{Z}$ ist ein Hauptidealring. Die Ideale sind die $n \mathds{Z}$.
% subsection 68 (end)

\subsection[Satz: Polynomringe über einen Körper sind Hauptidealringe]{Satz} % (fold)
\label{sub:69}
Sei $K$ ein Körper. Dann ist $K[X]$ ein Hauptidealring.
\minisec{Beweis}
Sei $I \subseteq K[X]$ ein Ideal. Ist $I=\{0\}$, so ist nichts zu zeigen. Sei also $I \not= \{0\}$. Wähle $f \in I \setminus \{0\}$ von minimalem Grad. Behauptung: $I= (f)$.
Sei $g \in I$. Division mit Rest liefert $g = q \cdot f +r$ wobei $r$ kleineren Grad als $f$ hat. Es ist aber $r = g - q \cdot f \in I$. Also folgt $r=0$, da $f$ minimal war.
Also $g = q \cdot f \in (f)$. \bewende
% subsection 69 (end)

\subsection[{Beispiel: $\mathds{Z}[X]$ ist kein Hauptidealring}]{Beispiel} % (fold)
\label{sub:610}
$\mathds{Z}[X]$ ist kein Hauptidealring: $(2X)$ ist kein Hauptideal in $\mathds{Z}[X]$.
% subsection 610 (end)

\subsection[Definition: Primideal und maximales Ideal]{Definition} % (fold)
\label{sub:611}
Sei $R$ ein Ring und $I \subseteq R$ ein Ideal.
\begin{enumerate}[(i)]
	\item $I$ heißt \Index{prim} oder ein \Index{Primideal}, wenn  $I \subsetneq R$ und für $r,s \in R$ gilt:
	\[
		r \cdot s \in I \Longrightarrow r \in I \text{ oder } s \in I
	\]
	\item $I$ heißt \bet{maximal}\index{Ideal!maximales}, wenn $I \subsetneq R$ und es kein Ideal $J$ gibt mit $I \subsetneq J \subsetneq R$. 
\end{enumerate}
% subsection 611 (end)

\subsection[Bemerkung, wann $(0)$ ein Primideal ist]{Bemerkung} % (fold)
\label{sub:612}
$(0)= \{0\}$ ist genau dann ein Primideal, wenn $R$ nullteilerfrei ist.
% subsection 612 (end)

\subsection[Lemma: $(0)$ ist in Körpern maximal]{Lemma} % (fold)
\label{sub:613}
$(0)= \{0\}$ ist genau dann maximal, wenn $R$ ein Körper ist.
\minisec{Beweis}
Ist $R$ ein Körper, so sind $\{0\}$ und $R$ die einzigen Ideale in $R$ (Ist $r \in I$ mit $r \not=0$, also $s= (s \cdot r ^{-1})r \in I$ für alle $s \in R$. Also $I=R$).
Sei umgekehrt $(0)$ ein maximales Ideal. Sei $r \in R \setminus \{0\}$. Zu zeigen: $r$ ist invertierbar in $R$. Betrachte das Ideal $(r)$. Es ist $(0) \subsetneq (r)$, da
$r \in (r)$. $(0)$ maximal $\Rightarrow (r) =R$. Also $1 \in (r)$. Also $\exists s \in R$ mit $1= s \cdot r$. Also $r ^{-1} = s \in R$. \bewende
% subsection 613 (end)

\subsection[Beispiel: Maximale und Primideale in $\mathds{Z}$]{Beispiel} % (fold)
\label{sub:614}
Sei $(0) \not= \underset{\not= \mathds{Z}}{I} = n \mathds{Z} \subseteq \mathds{Z}$ ein Ideal. Dann gilt $n \mathds{Z}$ ist maximal $\iff$ $n$ ist Primzahl $\iff$ ist ein 
Primideal.
\minisec{Beweis}
Für $n, m$ gilt
\[
	n  \text{ teilt } m \iff m Z \subseteq n \mathds{Z}
\]
Ist $n$ eine Primzahl und $n \mathds{Z} \subseteq k \mathds{Z}$. Dann folgt $k$ teilt $n$. Ist $n$ eine Primzahl so folgt $k \in \set{\pm n, \pm 1} $. Also 
$k \mathds{Z} \in \set{\mathds{Z}, n \mathds{Z}} $. Sei weiter $n$ eine Primzahl und $r \cdot s \in n \mathds{Z}$. Dann teilt $n$ das Produkt $r \cdot s$. Da $n$ eine 
Primzahl ist, muss $n$ $r$ oder $s$ teilen. Also $r\in n \mathds{Z}$ oder $s \in n \mathds{Z}$.\bewende
% subsection 614 (end)

\subsection[Satz: Zentrale Eigenschaften von maximalen und Primidealen]{Satz} % (fold)
\label{sub:615}
Sei $R$ ein Ring und $I \subsetneq R$ ein Ideal. Dann gilt:
\begin{enumerate}[(i)]
	\item $I$ ist prim $\iff \nicefrac{R}{I}$ ist nullteilerfrei
	\item $I$ ist maximal $\iff \nicefrac{R}{I}$ ist ein Körper.
\end{enumerate}
\minisec{Beweis}
\begin{enumerate}[(i)]
	\item Sei $I$ prim. Sei $(r+I)(s+I)= 0  \in \nicefrac{R}{I}$. Dann $r \cdot s + I = 0 \in \nicefrac{R}{I}$, also $r \cdot s \in I$. Mit $I$ prim folgt $r \in I$ oder 
	$s \in I$. $\Rightarrow r+I=0$ oder $s+I=0$.
	
	Sei $\nicefrac{R}{I}$ nullteilerfrei. Sei $r \cdot s  \in I$. Dann $(r+I)(s+I)= r \cdot s +I = 0 \in \nicefrac{R}{I}$. Da $\nicefrac{R}{I}$ nullteilerfrei ist, folgt:
	$r+I=0$ oder $s+I =0$ $\Rightarrow $ $r \in I$ oder $s \in I$. \bewende
	\item Sei $I$ maximal. Sei $a+I \in \nicefrac{R}{I}$, $a +I \not= 0 \in \nicefrac{R}{I}$, also $a \not\in I$. Gesucht ist ein multiplikatives Inverses zu $a+I$, also
	ein $b \in R$ mit $ab +I = 1+I$. Nun ist $(I,a) := \set[s+r \cdot a]{s \in I, r \in R} $ ein Ideal mit $I \subsetneq (I,a) \subseteq R$. Da $I$ maximal ist, folgt
	$(I,a)= R$. Da dann $1 \in (I,a)$ folgt $1= s+ b \cdot a$ mit $s \in I$ und $b \in R$. Also $a \cdot b +I = 1 +I$.
	
	Sei $\nicefrac{R}{I}$ ein Körper. Sei $J \subsetneq R$ ein Ideal mit $I \subsetneq J$. Dann ist $\overline{J} = \set[j+I]{j \in J}  $ ein Ideal in $\nicefrac{R}{I}$.
	Da $J \not= I$ ist $\overline{J} \not= \set{0+I}  $. Es folgt $\overline{J} = \nicefrac{R}{I}$ und damit $J=R$. \bewende
\end{enumerate}
% subsection 615 (end)

\subsection[Bemerkung: Mengenoperationen mit Idealen und koprime Ideale]{Bemerkung} % (fold)
\label{sub:616}
Seien $I,J \subseteq R$ Ideale
\begin{enumerate}[(i)]
	\item $I \cap J$ ist ein Ideal.
	\item $I \cup J$ ist nicht immer ein Ideal.
	\item $I+J := \set[a+b]{a \in I, b \in J}$ ist ein Ideal. Es ist das kleinste Ideal, das $I$ und $J$ enthält.
\end{enumerate}
Gilt $I+J = R$, so heißen $I$ und $J$ \Index{koprim}. 
% subsection 616 (end)

\subsection[Beispiel: Koprime Ideale in $\mathds{Z}$]{Beispiel} % (fold)
\label{sub:617}
Betrachte $\mathds{Z}$, dann gilt: $I := n \mathds{Z}, \, J := m \mathds{Z}$ sind genau dann koprim, wenn $n$ und $m$ teilerfremd sind. Allgemeiner gilt: 
\[
	n \mathds{Z} + m \mathds{Z} = \ggT(n,m)  \mathds{Z} \quad \text{ und } \quad n \mathds{Z} \cap m \mathds{Z} = \kgV(n,m) \mathds{Z}
\]
% subsection 617 (end)

\subsection{Chinesischer Restsatz} % (fold)
\label{sub:618}
Seien $I_1, \ldots , I_n \subseteq R$ Ideale, die paarweise koprim sind, also $I_i + I_j =R$ für $i\not= j$. Dann ist 
$\psi : R \to \nicefrac{R}{I_1} \times \ldots \times \nicefrac{R}{I_n}$ mit $\psi(x) := \enbrace*{x+ I_1, \ldots , x+ I_n}$ surjektiv und induziert einen Isomorphismus
\[
	\nicefrac{R}{I_1 \cap \ldots \cap I_n} = \nicefrac{R}{\Kern \psi} \cong \nicefrac{R}{I_1} \times \ldots \times \nicefrac{R}{I_n}
\]
\minisec{Beweis für $n=2$}
Da $I_1 + I_2 =R$ gibt es $a_1 \in I_1$ und $a_2 \in I_2$ mit $a_1 + a_2 = 1$. Seien $r_1 + I_1 \in \nicefrac{R}{I_1}$ und $r_2 + I_2 \in \nicefrac{R}{I_2}$ gegeben.
Gesucht ist $r \in R$ mit $r+ I_1 = r_1+I_1$, $r+ I_2 = r_2 + I_2$. Betrachte $r := r_1 a_2 + r_2 a_1$. Dann gilt 
\begin{align*}
	r + I_1 &= r_1 \cdot a_2 + r_2 \cdot a_1 + I_1 = r_1 \cdot a_2 +I_1 = r_1 \cdot a_2 + r_1 \cdot a_1 +I_1 = r_1(a_2 +a_1) +I_1 = r_1 + I_1 \\
	r + I_2 &= r_1 \cdot a_2 + r_2 \cdot a_1 + I_2 = r_2 \cdot a_1 + I_2 = r_2 \cdot a_1 + r_2 \cdot a_2 + I_2 = r_2(a_1+ a_2) + I_2 = r_2 + I_2
\end{align*}
Also ist $\psi$ surjektiv. Da $\Kern \psi = I_1 \cap \ldots \cap I_n$ folgt der Rest aus dem Homomorphiesatz. \bewende
\minisec{Beweis für allgemeines $n \in \mathds{N}$}
Nach Lemma \ref{sub:619} gibt es $d_j \in I_j$, $e_j \in J_j$ mit $d_j + e_j = 1$. Es folgt 
\[
	e_j + I_i = \begin{cases}
		0+ I_i = 0 \in \nicefrac{R}{I_i}, &\text{ falls }i \not= j\\
		e_j + d_j + I_j = 1+ I_j = 1 \in \nicefrac{R}{I_j}, &\text{ falls } i=j
	\end{cases}
\]
Sei $a_1, \ldots , a_n \in R$. Setze $x := a_1 \cdot e_1 + a_2 \cdot e_2 + \ldots + a_n \cdot e_n$. Dann gilt
\[
	x+I_j = (a_1 e_1 + I_j) + (a_2 e_2 + I_j) + \ldots + (a_n e_n + I_j) = a_j e_j + I_j = a_j + I_j.
\]
Also $\psi(x) = (a_1 + I_1, a_2 + I_2, \ldots , a_n +I_n)$. Damit ist $\psi$ surjektiv. \bewende
% subsection 618 (end)

\subsection[Lemma: Hilfslemma für den Chinesischen Restsatz]{Lemma} % (fold)
\label{sub:619}
Seien $I_1, \ldots , I_n \subseteq R$ paarweise koprime Ideale. Für $j \in \set{1, \ldots ,n} $ sei $J_j := \bigcap_{i\not= j} I_i$. Dann sind $I_j$ und $J_j$ koprim.
\minisec{Beweis}
Da $I_j $ koprim zu $I_i$ für $i \not= j$, gibt es $a_i \in I_i$, $a_i' \in I_j$ mit $1= a_i + a_i'$. Dann gilt $1 = \prod_{i \not= j}( a_1 +a_i') \in I_j + J_j$. Also
$R = I_j + J_j$ und $I_j$ und $J_j$ sind koprim.
% subsection 619 (end)

\subsection[Korollar zur Lösung von Kongruenzen]{Korollar} % (fold)
\label{sub:620}
Seien $a_1, \ldots , a_n \in \mathds{Z}$ paarweise teilerfremd. Seien $x_1, \ldots , x_n \in \mathds{Z}$ beliebig. Dann gibt es eine Lösung $x \in \mathds{Z}$ für die 
gemeinsamen Kongruenzen
\[
	x  \equiv x_i \mod a_i  \qquad {\color{light_gray} :\Leftrightarrow x + a_i \mathds{Z} = x_i + a_i \mathds{Z}} \tag{$\star$}
\]
Die Menge aller Lösungen von ($\star$) ist $x+ a_1 \cdot \ldots \cdot a_n \mathds{Z}$
\minisec{Beweis}
Betrachte $\psi : \mathds{Z} \to \nicefrac{\mathds{Z}}{a_1 \mathds{Z}}  \times \ldots \times \nicefrac{\mathds{Z}}{a_n \mathds{Z}} ,\, x \mapsto (x+a_1 \mathds{Z}, \ldots , x+ a_n \mathds{Z})$. Dann erfüllt $x$ die Kongruenzen $(\star)$ genau dann wenn
\[
	\psi(x) = (x_1 + a_1 \mathds{Z}, \ldots , x_n + a_n \mathds{Z})
\]
gilt. Da $\psi$ surjektiv ist, gibt es ein solches $x$. Alle weiteren Lösungen sind von der Form $x+y$ mit $y \in \Kern \psi$. \bewende
% subsection 620 (end)

\subsection[Beispiel: Bestimmung einer Lösung von Kongruenzen]{Beispiel} % (fold)
\label{sub:621}
Suche $x \in \mathds{Z}$ mit $x \equiv {\color{Cyan3} 1} \mod 3$, $x \equiv {\color{Firebrick2} 2} \mod 4$, $x \equiv {\color{Green2}1} \mod 5$. 
\begin{align*}
	1 &= \underset{\in 3 \mathds{Z}}{(21)} + \underset{\in 4 \mathds{Z}, 5 \mathds{Z}}{(-20)} = \underset{\in 4 \mathds{Z}}{(16)} +  \underset{\in 3 \mathds{Z}, 5 \mathds{Z}}{(-15)} = \underset{\in 5 \mathds{Z}}{(25)} + \underset{\in 3 \mathds{Z}, 4 \mathds{Z}}{(-24)} \\
	x &:= {\color{Cyan3} 1} \cdot (-20) + {\color{Firebrick2} 2} \cdot (-15) + {\color{Green2}1} \cdot (-24) = -20 -30-24 = -74
\end{align*}
% subsection 621 (end)
% section 6 (end)
\newpage


\section{Primfaktorzerlegung} % (fold)
\label{sec:7}
In Kapitel \ref{sec:7} sei $R$ immer ein kommutativer nullteilerfreier Ring.

\subsection[Definition: Gruppe der Einheiten]{Definition} % (fold)
\label{sub:71}
\[
	R^\times := \set[\varepsilon \in R]{\exists \delta \in R \text{ mit } \varepsilon \cdot \delta = 1} 
\]
heißt die Gruppe der \Index{Einheiten}. 
% subsection 71 (end)

\subsection[Definition: Teiler]{Definition} % (fold)
\label{sub:72}
Seien $a,b\in R$. Wir sagen $a$ teilt $b$, falls es ein $x \in R$ gibt mit $b = a \cdot x$. Wir schreiben dafür auch $a \mid b$.
% subsection 72 (end)

\subsection[Bemerkung: Multiplikation mit Einheiten hat keinen Einfluss auf Teilbarkeit]{Bemerkung} % (fold)
\label{sub:73}
Für $a,b \in R$, $\delta , \varepsilon \in R^\times$ gilt: $a \mid b \iff \varepsilon \cdot a \mid \delta \cdot b$.
% subsection 73 (end)

\subsection[Definition: irreduzible Elemente und Primelemente in Ringen]{Definition} % (fold)
\label{sub:74}
Sei $p \in R \setminus \enbrace*{R^\times \cup \set{0} } $
\begin{enumerate}[i)]
	\item $p$ heißt \Index{irreduzibel}, falls gilt $p= a \cdot b$ mit $a,b \in R$ $\Longrightarrow$ $a \in R^\times$ oder $b \in R^\times$
	\item $p$ heißt \Index{prim}  oder ein \Index{Primelement}, falls gilt:
	\[
		p \mid a \cdot b \text{ mit } a,b \in R \Longrightarrow p \mid a \text{ oder } p \mid b
	\] 
\end{enumerate}
% subsection 74 (end)

\subsection[Bemerkung: Primelemente in $\mathds{Z}$]{Bemerkung} % (fold)
\label{sub:75}
Für $n \in \mathds{Z} \setminus \set{-1,0,1} $. $n$ ist prim $\Leftrightarrow$ $n$ ist eine Primzahl $\iff$ $n$ ist irreduzibel.
% subsection 75 (end)

\subsection[Lemma: Primelemente und Primideale; prim $\Rightarrow $ irreduzibel]{Lemma} % (fold)
\label{sub:76}
Sei $p \in R \setminus \enbrace*{R^\times \cup \{0\}} $.
\begin{enumerate}[i)]
	\item $p$ ist prim $\iff$ $(p)$ ist Primideal.
	\item $p$ ist prim $\Longrightarrow$ $p$ ist irreduzibel.
\end{enumerate}
\minisec{Beweis}
\begin{enumerate}[i)]
	\item Ist klar, da $a \in (p) \iff p \mid a$.
	\item Sei $p$ prim und $p = a \cdot b$. Insbesondere $p \mid a \cdot b$ und es folgt $p \mid a$ oder $p \mid b$. O.B.d.A: $p \mid a$. Also $p \cdot r = a$ für ein 
	$r \in R$. Es folgt $p = a \cdot b = p \cdot r \cdot b$, also $1= r \cdot b$, also $b \in R^\times$. \bewende
\end{enumerate}
% subsection 76 (end)

\subsection[Bemerkung: In Hauptidealringen ist prim und irreduzibel äquivalent]{Bemerkung} % (fold)
\label{sub:77}
Ist $R$ ein Hauptidealring, so gilt auch "{}irreduzibel $\Rightarrow$ prim"{}. Im Allgemeinen ist das aber nicht richtig.
% subsection 77 (end)

\subsection[Definition: Faktorieller Ring]{Definition} % (fold)
\label{sub:78}
$R$ heißt \Index{faktoriell}, falls sich jedes $r \in R \setminus \enbrace*{R^\times \cup \{0\}} $ als Produkt von Primelementen schreiben lässt. 
% subsection 78 (end)

\subsection[Bemerkung: Hauptidealringe sind faktoriell]{Bemerkung} % (fold)
\label{sub:79}
Hauptidealringe sind faktoriell (LA2). Insbesondere ist $K[X]$ faktoriell, falls $K$ ein Körper ist.
% subsection 79 (end)

\subsection[Lemma: In faktoriellen Ringen ist prim äquivalent zu irreduzibel]{Lemma} % (fold)
\label{sub:710}
Sei $R$ ein faktorieller Ring und $p \in R\setminus \enbrace*{R^\times \cup \{0\}} $. Dann gilt: $p$ prim $\iff$ $p$ irreduzibel.
\minisec{Beweis}
\begin{description}
	\item["$\Rightarrow $":] ist sowieso richtig
	\item["$\Leftarrow$":] Sei $p$ irreduzibel. Aus $R$ faktoriell folgt $p = q_1 \cdot \ldots \cdot q_n$ mit $q_1, \ldots , q_n$ Primelementen. Insbesondere 
	$q_i \not\in R^\times$. Da $p$ irreduzibel ist, folgt $n=1$ und $p= q_1$ ist prim. \bewende
\end{description}
% subsection 710 (end)

\subsection{Lemma (Eindeutigkeit der Primfaktorzerlegung)} % (fold)
\label{sub:711}
Sei $R$ faktoriell. Sind $p_1, \ldots , p_n$ und $q_1, \ldots , q_m$ Primelemente mit $p_1 \cdot \ldots \cdot p_n = q_1 \cdot \ldots \cdot q_m$ so gilt $n=m$. Weiter gibt es Einheiten $\varepsilon_1, \ldots ,\varepsilon_n$ sodass nach Umordnung $p_i = \varepsilon_i \cdot q_i$ für $i=1,\ldots ,n$.
\minisec{Beweis}
Aus $p_1 \mid q_1 \cdots q_m$ folgt, möglicherweise nach Umnummerierung, $p_1 \mid q_1$. Da $q_1$ irreduzibel ist, folgt $q_1 = \varepsilon_1 \cdot p_1$. Weiter folgt
$p_2 \cdot \ldots \cdot p_n = (\varepsilon_1 \cdot q_2) \cdot \ldots \cdot q_m$. Per Induktion\marginnote{etwas –Zitat– "{}sloppy" \ldots :D} über die Anzahl der Faktoren folgt $n= m$ und $q_i =  \varepsilon_i p_i$
mit Einheiten $\varepsilon_i$. \bewende
% subsection 711 (end)

\subsection[Beispiele: Einheitengruppen, irreduzible Elemente]{Beispiele} % (fold)
\label{sub:712}
\begin{enumerate}[(i)]
	\item $\mathds{Z}^\times = \{-1,1\}$, $K$ Körper $K^\times = K \setminus \{0\}$.
	$K[X]^\times = K^\times$. $R$ Ring $R[X]^\times = R^\times$.
	\item In $K[X]$, $K$ Körper ist $(X- \alpha)$ für alle $\alpha \in K$ irreduzibel, also prim. \\
	(Beweis: $X- \alpha \mid f \in K[X] \iff f(\alpha)=0$)
	\item In $\mathds{R}[X]$ ist $X^2+1$ irreduzibel (also prim)
	\[
		\text{Was ist }\nicefrac{\mathds{R}[X]}{(X^2+1)} \text{ ?}
	\]
\end{enumerate}
% subsection 712 (end)

\subsection[Definition: $\ggT$ und $\kgV$]{Definition} % (fold)
\label{sub:713}
Sei $R$ ein Ring und $a_1, \ldots , a_n \in R$
\begin{enumerate}[i)]
	\item $d \in R$ heißt ein \Index{größter gemeinsamer Teiler ($\ggT$)} von $a_1, \ldots , a_n$ falls gilt:
	\begin{enumerate}[(a)]
		\item $d \mid a_i$ für $i=1,\ldots ,n$
		\item Ist $d' \in R$ mit $d' \mid a_i$ für $i=1,\ldots ,n$, so gilt $d' \mid d$.
	\end{enumerate} 
	\item $v \in R$ heißt ein \Index{kleinstes gemeinsames Vielfaches ($\kgV$)} von $a_1, \ldots , a_n$ falls gilt:
	\begin{enumerate}[(a)]
		\item $a_i \mid v$ für $i=1, \ldots ,n$
		\item Ist $v' \in R$ mit $a_i \mid v'$ für $i=1,\ldots ,n$, so gilt $v \mid v'$.
	\end{enumerate} 
\end{enumerate}
% subsection 713 (end)

\subsection[Bemerkung: Existenz von $\ggT$ und $\kgV$]{Bemerkung} % (fold)
\label{sub:714}
In beliebigen Ringen müssen weder $\kgV$ noch $\ggT$ existieren. Falls $\kgV (a_1, \ldots , a_n)$ existiert, so ist es auch eindeutig bis auf Multiplikation mit einer Einheit. Ebenso für $\ggT$.

In faktoriellen Ringen existieren $\kgV$ und $\ggT$ immer und können über die Primfaktorzerlegung bestimmt werden.
% subsection 714 (end)

\subsection[Satz: Zusammenhang zwischen Idealen und $\ggT$ und $\kgV$]{Satz} % (fold)
\label{sub:715}
Sei $R$ ein Hauptidealring und $a_1,\ldots ,a_n \in R$
\begin{enumerate}[i)]
	\item Wenn $(d) = (a_1, \ldots ,a_n)$, so ist $d$ ein $\ggT$ von $a_1, \ldots ,a_n$
	\item Wenn $(v) = (a_1) \cap (a_2) \cap \ldots \cap (a_n)$, so ist $v$ ein $\kgV$ von $a_1, \ldots , a_n$
\end{enumerate}
\minisec{Beweis}
\begin{enumerate}[i)]
	\item Da $a_i \in (a_1, \ldots ,a_n) = (d) =  \set[x \cdot d]{x \in R} $ gilt $d \mid a_i$ für alle $i$. Sei $d'$ ein weiterer gemeinsamer Teiler. Dann $a_i \in (d')$
	für $i=1,\ldots ,n$. Es folgt $(d) = (a_1, \ldots , a_n) \subseteq (d')$. Es folgt $d \in (d')$, also $d' \mid d$.
	\item Übung
\end{enumerate}
% subsection 715 (end)

\subsection[Bemerkung, die den Namen eigentlich nicht verdient hat]{Bemerkung} % (fold)
\label{sub:716}
$a \mid b \iff b \in (a)$.
% subsection 716 (end)

\subsection{Euklidischer Algorithmus} % (fold)
\label{sub:717}
Seien $f_1, f_2 \in K[X]$, $f_2 \not= 0$. Durch wiederholte Division mit Rest erhalten wir: 
\begin{align*}
	f_1 &= q_1 \cdot f_2 + f_3 \qquad d(f_3) < d(f_2) \\
	f_2 &= q_2 \cdot f_3 + f_4 \qquad d(f_4) < d(f_3) \\
	\vdots & \\
	f_n &= q_n \cdot f_{n+1}
\end{align*}
% subsection 717 (end)

\subsection[Lemma: Der euklidische Algorithmus berechnet den $\ggT$]{Lemma} % (fold)
\label{sub:718}
$f_{n+1}= \ggT(f_1,f_2)$
\minisec{Beweis}
Den Algorithmus aufsteigend sehen wir, dass $f_{n+1}$ alle $f_i$ teilt. Insbesondere ist $f_{n+1}$ ein gemeinsamer Teiler von $f_1$ und $f_2$. Ist $d'$ ein weiterer gemeinsamer
Teiler von $f_1$ und $f_2$, so sehen wir, indem wir dem Algorithmus absteigend folgen, dass $d'$ alle $f_i$ teilt. Insbesondere gilt $d' \mid f_{n+1}$. Damit ist $f_{n+1} = \ggT(f_1, f_2)$. \bewende
% subsection 718 (end)
% section 7 (end)
\newpage


\section{Satz von Gauß} % (fold)
\label{sec:8}
Sei $R$ in \ref{sec:8} immer nullteilerfrei, mit $1$ und kommutativ.

\subsection{Satz von Gauß} % (fold)
\label{sub:81}
Sei $R$ faktoriell. Dann ist auch $R[X]$ faktoriell.
% subsection 81 (end)

\subsection[Bemerkung: Hinweis zu Beweistrategie]{Bemerkung} % (fold)
\label{sub:82}
Im Beweis benutzen wir, dass $K[X]$ faktoriell ist, falls $K$ ein Körper ist.
% subsection 82 (end)

\subsection[Bemerkung: Werkzeuge für den Beweis des Satzes]{Bemerkung} % (fold)
\label{sub:83}
Sei $f \in \mathds{Q}[X]$. Dann gibt es $c \in \mathds{Z}$ mit $\tilde f := c \cdot f \in \mathds{Z}[X]$. Für jede Primzahl $p$ induziert 
$\mathds{Z} \to \nicefrac{\mathds{Z}}{p \mathds{Z}}$ einen
Ringhomomorphismus $\Phi : Z[X] \to \nicefrac{\mathds{Z}}{p \mathds{Z}}[X]$ mit $\sum_{n=1}^{N} a_n X^n \mapsto \sum_{n=1}^{N}(a_n + p \mathds{Z}) X^n$.
% subsection 83 (end)

\subsection{Konstruktion (Körper aus einem Ring)} % (fold)
\label{sub:84}
Sei $B = R \times R\setminus \{0\}$. Auf $B$ betrachten wir die Äquivalenzrelation 
\[
	(z,n) \sim (z',n') :\Leftrightarrow n'\cdot  z = n \cdot z'.
\]
Transitivität: $(z,n) \sim (z',n') \sim (z'',n'')$. Dann gilt
\[
	n'' \cdot z \cdot  n' = n'' \cdot z'\cdot  n = n' \cdot z'' \cdot n. 
\]
Da $R$ nullteilerfrei ist, folgt $n'' \cdot z = z'' \cdot n \Rightarrow (z,n) \sim (z'',n'')$. Wir schreiben nun $\frac{z}{n} $ für die Äquivalenzklasse von $(z,n)$
bezüglich dieser Äquivalenzrelation und $Q(R)$ für die Menge der Äquivalenzklassen.
% subsection 84 (end)

\subsection[Lemma: Körperstruktur auf $Q(R)$]{Lemma} % (fold)
\label{sub:85}
Durch $\frac{z}{n} + \frac{z'}{n'} := \frac{n' \cdot z+ n \cdot z'}{n \cdot n'}$ und $\frac{z}{n} \cdot \frac{z'}{n'} := \frac{z \cdot z'}{n \cdot n'}$ wird $Q(R)$ zu einem
Körper.
\minisec{Beweis}
Nachrechnen.
% subsection 85 (end)

\subsection[Definition: Quotientenkörper]{Definition} % (fold)
\label{sub:86}
$Q(R)$ heißt der \Index{Quotientenkörper} von $R$. \\
Beispiel: $Q(\mathds{Z}) = \mathds{Q}$ 
% subsection 86 (end)

\subsection[Bemerkung: Einbettung des Ringes in seinen Quotientenkörper]{Bemerkung} % (fold)
\label{sub:87}
Die Abbildung $i : R \to Q(R)$ mit $i(r) := \frac{r}{1} $ ist ein injektiver Ringhomomorphismus. Wir unterscheiden oft nicht zwischen $r$ und $i(r)= \frac{r}{1} $ und 
fassen $R$ als Unterring von $Q(R)$ auf.
% subsection 87 (end)

\subsection[Definition: Repräsentantensystem]{Definition} % (fold)
\label{sub:88}
Eine Menge $P \subseteq R$ von Primelementen heißt
ein \Index{Repräsentantensystem} der Primelemente, falls es zu jedem Primelement $q \in R$ ein eindeutiges $p \in P$ gibt
und $\varepsilon \in R^\times$ mit $q= \varepsilon \cdot p$. 
% subsection 88 (end)

\subsection[Beispiel für Repräsentantensysteme]{Beispiel} % (fold)
\label{sub:89}
\begin{enumerate}[i)]
	\item Die (positiven) Primzahlen $P=\set{2,3,5,7,11, \ldots} $ sind ein Repräsentantensystem für die Primzahlen in $\mathds{Z}$.
	\item Die Polynome $\set[(X- \alpha)]{\alpha \in \mathds{C}} $ sind ein Repräsentantensystem für die Primelemente in $\mathds{C}[X]$.
\end{enumerate}
% subsection 89 (end)

\subsection[Bemerkung: Darstellung beliebiger Ringelemente durch Repräsentantensystem]{Bemerkung} % (fold)
\label{sub:810}
Sei $P \subseteq R$ ein Repräsentantensystem der Primelemente. Sei $R$ faktoriell.
\begin{enumerate}[i)]
	\item Jedes $a \in R \setminus \{0\}$ lässt sich eindeutig schreiben als
	\[
		a = \varepsilon \cdot \prod_{p \in P} p^{v_p(a)}
	\]
	wobei $\varepsilon \in R^\times$ und $v_p(a) \in \mathds{N}_{\ge 0}$, $v_p(a)=0$ für fast alle $p \in P$.
	\item Jedes $x \in Q(R) \setminus \{0\}$ lässt sich eindeutig schreiben als
	\[
		x = \varepsilon \cdot \prod_{p \in P} p^{v_p(x)}
	\]
	wobei $\varepsilon \in R^\times$, $v_p(x) \in \mathds{Z}$ und $v_p(x)=0$ für fast alle $p \in P$. Wir setzen $v_p(0) := \infty$.
\end{enumerate}
% subsection 810 (end)

\subsection[Definition: Ordnung eines Polynoms bezüglich $p$]{Definition} % (fold)
\label{sub:811}
Sei $f = a_n X^n + a_{n-1} X^{n-1} + \ldots + a_1 X + a_0 \in Q(R)[X]$. $f \not= 0$, wobei $R$ faktoriell ist. Dann setzen wir für $p \in P$
\[
	v_p(f) := \min_{i=0,\ldots ,n} v_p(a_i) \quad , \quad v_p(0) := \infty
\]
% subsection 811 (end)

\subsection[Bemerkung: Ordnungsabbildung zum testen, ob Koeffizienten im Ring]{Bemerkung} % (fold)
\label{sub:812}
$f \in R[X] \iff v_p(f)\ge 0$ $\forall p \in P$.
% subsection 812 (end)

\subsection{Lemma von Gauß} % (fold)
\label{sub:813}
Sei $R$ faktoriell und $p \in R$ prim. Für $f,g \in Q(R)[X]$ gilt dann: 
\[
	v_p(f \cdot g) = v_p(f)+ v_p(g) \tag{$\star$}
\]
\minisec{Beweis}
Wir betrachten zunächst mehrere Spezialfälle:
\begin{enumerate}[(1)]
	\item $f =a, g=b \in Q(R)$. Dann ist $(\star)$ klar.
	\item $f= a \in Q(R), g= b_n X^n + \ldots + b_0 \in Q(R)[X]$. Dann gilt
	\[
		v_p(a \cdot g) = v_p(a b_n X^n + \ldots + a b_0) = \min_i v_p(a \cdot b_i) = \min_i (v_p(a)+ v_p(b_i)) = v_p(a) + v_p(g).
	\]
	\item $f=0$ oder $g=0$ ist auch klar.
	\item $f,g \in R[X]$ mit $v_p(f)=v_p(g)=0$. 
	
	Wir schreiben $R_p$ für das von $p$ erzeugte Hauptideal in $R$ und $R[X]_p$ für das von $p$ erzeugte Hauptideal in $R[X]$. Betrachte den durch $R \to \nicefrac{R}{R_p}$
	induzierten Ringhomomorphismus $\Phi : R[X] \to \nicefrac{R}{R_p}[X]$. Es gilt 
	\[
		\Kern \Phi = \set[{h \in R[X]}]{p \text{ teilt alle Koeffizienten von }h} \cup \{0\}  = \set[{h \in R[X]}]{v_p(h) > 0} \cup \{0\}. 
	\]
	Wegen $v_p(f) = v_p(g)=0$ folgt $\Phi(f) \not= 0 \not= \Phi (g)$. Es ist $\Phi(f \cdot g) = \Phi(f) \cdot \Phi(g)$ \hfill $(\#)$ \\
	$p$ ist prim in $R$ $\Rightarrow $ $R_p \subseteq R$ ist Primideal $\Rightarrow \nicefrac{R}{R_p}$ nullteilerfrei $\Rightarrow \nicefrac{R}{R_p}[X]$ ist nullteilerfrei.
	Also $(\#) \Rightarrow \Phi(f \cdot g) \not= 0$ da $\Phi(f) \not= 0 \not= \Phi(g)$. Es folgt also $f,g \not\in \Kern \Phi \Rightarrow v_p(f \cdot g)=0$.
\end{enumerate}
Seien nun $f = a_n X^n + \ldots + a_0$, $g = b_m X^m + \ldots + b_0 \in Q(R)[X]$, $f,g \not= 0$ aber sonst beliebig. Sei $A \in R$ das Produkt der Nenner der $a_i$ und 
$B \in R$ das Produkt der Nenner der $b_i$. Dann $A \cdot f, B \cdot g \in R[X]$. Setze $N := - v_p(A \cdot f), M := - v_p(B \cdot g)$. Mit (2) folgt:
\[
	v_p(p^N \cdot A \cdot f) = v_p(p^N)+ v_p(A \cdot f) = N + v_p(A \cdot f) = 0 \qquad v_p(p^M \cdot B \cdot g) = 0
\]
(4) liefert 
\[
	v_p\big((p^N \cdot A \cdot f)(p^M \cdot B \cdot g) \big) = v_p(p^N \cdot A \cdot f) + v_p(p^M \cdot B \cdot g) \tag{$\star\star$}
\]
mit (2) folgt
\begin{align*}
	v_p(p^N \cdot A \cdot f \cdot p^M \cdot B \cdot g) &= v_p(p^N) + v_p(A) + v_p(B) + v_p(p^M) + v_p(f \cdot g) \\
	v_p(p^N \cdot A \cdot f) &= v_p(p^N) + v_p(A) + v_p(f) \\
	v_p(p^M \cdot B \cdot g) &= v_p(B) + v_p(p^M) + v_p(g)
\end{align*}
$(\star\star)$ liefert $v_p(f \cdot g) = v_p(f) + v_p(g)$ \bewende
% subsection 813 (end)

\subsection[{Korollar: Normierte Faktoren von $h \in R[X]$ sind auch in $R[X]$}]{Korollar} % (fold)
\label{sub:814}
Sei $R$ faktoriell und $h=X^n + a_{n-1} X^{n-1} + \ldots + a_0 \in R[X]$ ein normiertes Polynom. Ist $h=f \cdot g$ mit $f,g \in Q(R)[X]$ beide normiert so gilt 
$f,g  \in R[X]$.
\minisec{Beweis}
Zu zeigen: $v_p(f), v_p(g) \ge 0$ für alle Primelemente $p \in R$. Da $f,g$ und $h$ normiert sind, gilt $v_p(f), v_p(g), v_p(h) \le 0$. Wegen $h \in R[X]$ folgt $v_p(h)=0$. Mit dem Gauß-Lemma folgt: 
\[
	0 = v_p(h)= v_p(f \cdot g) = v_p(f) + v_p(g) 
\]
Also $v_p(f)= v_p(g)=0$. \bewende
% subsection 814 (end)

\subsection[Definition: Primitives Polynom]{Definition} % (fold)
\label{sub:815}
Sei $R$ faktoriell. $f= a_n X^n + a_{n-1} X^{n-1} + \ldots + a_0 \in R[X]$ heißt \bet{primitiv}\index{Polynom!primitives}, wenn $\ggT(a_n, \ldots , a_0)=1$ ist.
% subsection 815 (end)

\subsection[Bemerkungen zu primitiven Polynomen]{Bemerkung} % (fold)
\label{sub:816}
\begin{enumerate}[(1)]
	\item Normierte Polynome sind primitiv.
	\item $f \in R[X]$ ist primitiv $\iff \forall  p \in R $ prim gilt $v_p(f)= 0$
	\item Jedes $f \in Q(R)[X]$, $f\not= 0$ lässt sich faktorisieren als $f= a \cdot \tilde f$ mit $a \in Q(R)$ und $\tilde f \in R[X]$ primitiv:
	\[
		a := \prod_{p \in P} p^{v_p(f)} \quad , \quad \tilde f = a ^{-1} \cdot f.
	\]
	Bsp: $\frac{3}{2} X^2 + \frac{3}{7}X + \frac{9}{2}  = \frac{1}{14} \enbrace*{21 X^2 +6 X + 63}  = \frac{3}{14} \enbrace*{7 X^2 +2 X + 21}$
\end{enumerate}
% subsection 816 (end)

\subsection[{Proposition: Primelemente in $R[X]$ und Faktorisierung in $R[X]$}]{Proposition} % (fold)
\label{sub:817}
Sei $R$ faktoriell
\begin{enumerate}[i)]
	\item Jedes Primelement $p \in R$ ist auch ein Primelement in $R[X]$.
	\item Sei $q \in R[X]$ primitiv und ein Primelement in $Q(R)[X]$. Dann ist $q$ auch prim in $R[X]$.
	\item Jedes $f \in R[X]\setminus (R^\times \cup \{0\})$ lässt sich faktorisieren als $f=p_1 \cdot \ldots \cdot p_n \cdot q_1 \cdot \ldots \cdot q_m$ mit 
	$p_1, \ldots , p_n \in R$ prim und $q_1, \ldots , q_m \in R[X]$ primitiv und prim in $Q(R)[X]$.
\end{enumerate}
\minisec{Beweis}
\begin{enumerate}[i)]
	\item Sei $\Phi : R[X] \to \nicefrac{R}{Rp}[X]$ der von $R \to \nicefrac{R}{Rp}$ induzierte Ringhomomorphismus. Dann liegt $f \in R[X]$ genau dann im $\Kern \Phi $
	wenn alle Koeffizienten von $f$ durch $p$ geteilt werden. Also $\Kern \Phi = R[X]p$. Mit dem Homomorphiesatz folgt:
	\[
		\nicefrac{R[X]}{R[X]p} \cong \nicefrac{R}{Rp}[X]
	\]
	$p$ prim $\Rightarrow Rp \subseteq P$ Primideal $\Rightarrow \nicefrac{R}{Rp}$ nullteilerfrei $\Rightarrow \nicefrac{R}{Rp}[X]$ nullteilerfrei.
	Daraus folgt $\nicefrac{R[X]}{R[X]p}$ nullteilerfrei, also ist $R[X]p \subseteq R[X]$ ein Primideal. Damit folgt $p \in R[X]$ ist Primelement.
	\item Seien $f,g \in R[X]$ mit $q \mid f \cdot g$ in $R[X]$. Damit teilt $q$ auch $f \cdot g$ in $Q(R)[X]$. Da $q$ prim ist in $Q(R)[X]$ folgt $q \mid f$ oder 
	$q \mid g$ in $Q(R)[X]$. O.B.d.A. $q \mid f$. Also $f = q \cdot h$ mit $h \in Q(R)[X]$. Sei $p \in R$ prim. Mit dem Lemma von Gauß folgt
	\[
		0 \le v_p(f) = v_p(q) + v_p(h) = v_p(h) \marginnote{da $q$ primitiv}
	\]
	Also ist $v_p(h) \ge 0$ für alle $p$. Damit $h \in R[X]$. Es folgt $q \mid f$ schon in $R[X]$.
	\item Sei $f= a \cdot \tilde f$ mit $a= \ggT$ der Koeffizienten von $f$. Dann ist $\tilde f \in R[X]$ primitiv. Da $R$ faktoriell ist, gibt es Primelemente 
	$p_1, \ldots , p_n \in R$ mit $a= p_1, \ldots , p_n$. Da $Q(R)[X]$ faktoriell ist,\marginnote{$Q(R)$ Körper $\Rightarrow $ $Q(R)[X]$ HIR} 
	gibt es $\tilde f_1 , \ldots , \tilde f_m \in Q(R)[X]$ prim mit
	\[
		\tilde f = \tilde f_1 \cdot \ldots \cdot \tilde f_m.
	\]
	Jedes $\tilde f_i$ lässt sich schreiben als $\tilde f_i = c_i \cdot q_i$ mit $c_i \in Q(R)$ und $q_i \in R[X]$ primitiv. Setze $c := c_1 \cdot \ldots \cdot c_m$.
	Es folgt $\tilde f = c \cdot q_1 \cdot \ldots \cdot q_m$. Mit $\tilde f_i$ ist auch $q_i$ prim $\in Q(R)[X]$. Es bleibt zu zeigen: $c \in R^\times$.
	
	Sei $p \in R$ prim. Dann 
	\[
		0 = v_p(\tilde f) = v_p(c) + v_p(q_1) + \ldots + v_p(q_m) = v_p(c)
	\]
	Also $v_p(c) = 0$ für alle $p \in R$ prim. Damit ist $c \in R^\times$. \bewende
\end{enumerate}
% subsection 817 (end)

\subsection{Beweis des Satz von Gauß} % (fold)
\label{sub:818}
Um zu zeigen, dass $R[X]$ faktoriell ist, müssen wir zeigen, dass sich jedes $f \in R[X] \setminus ( R^\times \cup \{0\})$ als Produkt von Primelementen schreiben lässt. Dies folgt sofort aus der Proposition \ref{sub:817}. \bewende
% subsection 818 (end)

\subsection[{Korollar: Äquivalenz zu \glqq prim in $R[X]$\grqq}]{Korollar} % (fold)
\label{sub:819}
Sei $R$ faktoriell. Dann ist $q \in R[X]$ genau dann prim wenn 
\begin{enumerate}[i)]
	\item $q \in R$ prim in $R$, \uline{oder}
	\item $q \in R[X]$ primitiv und prim in $Q(R)[X]$.
\end{enumerate}
\minisec{Beweis}
Nach i) und ii) in Proposition \ref{sub:817} sind die Elemente aus i) und ii) im Korollar prim in $R[X]$. Ist $q \in R[X]$ prim, so lässt sich wie in iii) der Proposition
faktorisieren. Da $q$ prim ist, besteht diese Faktorisierung nur aus einem Element. \bewende
% subsection 819 (end)
% section 8 (end)
\newpage


\section{Irreduzible Polynome} % (fold)
\label{sec:9}
Sei in Kapitel \ref{sec:9} $R$ immer faktoriell und $K:= Q(R)$
\subsection[Beispiele für irreduzible Polynome]{Beispiel} % (fold)
\label{sub:91}
\begin{enumerate}[i)]
	\item Sei $K$ ein Körper. Dann sind alle Polynome vom Grad $1$ irreduzibel.
	\item Ist $K$ algebraisch abgeschlossen, so sind die Polynome vom Grad $1$ genau die irreduziblen Polynome.
	\item Die irreduziblen Polynome über $\mathds{R}$ sind genau
	\begin{enumerate}[a)]
		\item alle Polynome vom Grad 1 und
		\item alle Polynome vom Grad 2, die keine Nullstelle in $\mathds{R}$ haben.
	\end{enumerate}
	\item Sei $K$ ein Körper und $p \in K[X]$ mit $2 \le d(p) \le 3$. Dann ist $p$ genau dann irreduzibel, wenn $p$ keine Nullstelle in $K$ hat.
	\item $(X^2+1)^2 \in \mathds{R}[X]$ ist nicht irreduzibel, aber besitzt keine Nullstelle in $\mathds{R}$.
\end{enumerate}
% subsection 91 (end)

\subsection[Bemerkung: Erkenntnisse zu primitiven Polynomen aus Kapitel \ref{sec:8}]{Bemerkung} % (fold)
\label{sub:92}
Sei $f \in K[X]\setminus \{0\}$. Dann gibt es $c \in K^\times$ und $\tilde f  \in R[X]$ primitiv mit $f = c \cdot \tilde f$. Es gilt Kapitel \ref{sec:8}.
$f$ ist irreduzibel in $K[X]$ $\iff$ $\tilde f $ ist irreduzibel in $K[X]$ $\iff \tilde f$ ist irreduzibel in $R[X]$.
% subsection 92 (end)

\subsection{Satz (Reduktionskriterium)} % (fold)
\label{sub:93}
Sei $p \in R$ prim. Sei $f = a_n X^n + \ldots + a_0 \in R[X]$, $\grad f = n \ge 1$. Sei $\Phi : R[X] \to \nicefrac{R}{Rp}[X]$ der von $R \to \nicefrac{R}{Rp}$ induzierte
Ringhomomorphismus. Weiter sei
\begin{enumerate}[(i)]
	\item $p \nmid a_n$
	\item $\Phi (f) $ irreduzibel in $\nicefrac{R}{Rp}[X]$
\end{enumerate}
Dann ist $f$ irreduzibel in $K[X]$.
\minisec{Beweis}
Wir nehmen zunächst an, dass $f$ primitiv ist. Dann genügt es zu zeigen, dass $f$ irreduzibel in $R[X]$ ist. Sei also $f=g \cdot h$ mit $g,h \in R[X]$. Dann ist $a_n$ das
Produkt der höchsten Koeffizienten von $g$ und $h$. Insbesondere teilt $p$ keinen dieser Koeffizienten und es folgt $\grad g = \grad \Phi (g)$ und 
$\grad h = \grad \Phi(h)$. Da $\Phi (f)= \Phi (g) \cdot \Phi(h)$ irreduzibel ist, ist $\Phi (g) \in \enbrace*{\nicefrac{R}{Rp}}^\times $ oder 
$\Phi (h) \in \enbrace*{\nicefrac{R}{Rp}}^\times$. O.B.d.A $\Phi (g) \in \enbrace*{\nicefrac{R}{Rp}}^\times $. Insbesondere $\grad (\Phi (g)) =0$. Also auch $\grad (g) =0$.
Also $g \in R$. Wegen $f=g \cdot h$ gilt $g \mid a_i \in R$ für $i=n, \ldots ,0$. Da $f$ primitiv ist, folgt $g \in R^\times$. 

Ist $f$ nicht primitiv, so gilt $f= \ggT(a_i) \cdot \tilde f$ mit $\tilde f \in R[X]$ primitiv. Es genügt zu zeigen, dass $\tilde f$ irreduzibel ist. Mit $\Phi (f)$ ist 
auch $\Phi(\tilde f)$ irreduzibel in $\nicefrac{R}{Rp}[X]$. Sei $\tilde a_n$ der höchste Koeffizient von $\tilde f$. Dann ist $a_n = \ggT(a_i) \cdot \tilde a_n$. Da
$p \nmid a_n$ folgt $p \nmid \tilde a_n$. Also können wir den ersten Teil des Beweises auf $\tilde f$ anwenden. Daher ist $\tilde f$ und damit auch $f$ irreduzibel. 
\bewende
% subsection 93 (end)

\subsection[Beispiel zur Anwendung des Reduktionskriteriums]{Beispiel} % (fold)
\label{sub:94}
$f = 19 X^3 + 17X + 15$ ist irreduzibel in $\mathds{Q}[X]$ (und in $\mathds{Z}[X]$): Das Bild von $f$ unter 
$\Phi : \mathds{Z}[X] \to \nicefrac{\mathds{Z}}{2 \mathds{Z}}[X]$ ist $\Phi(f) = X^3+ X + 1$ irreduzibel, da $\Phi(f)$ keine Nullstelle in 
$\nicefrac{\mathds{Z}}{2 \mathds{Z}}$ hat: $\Phi (f)(1) = 1+1+1 = 1 \not=0$ und $\Phi (f)(0) = 0 +0 +1 = 1 \not=0$.
% subsection 94 (end)

\subsection{Satz (Eisenstein)} % (fold)
\label{sub:95}
Sei $f = a_n X^n + \ldots + a_0 \in R[X]$ ein primitives Polynom von $\grad f = n \ge 1$. Weiter sei $p \in R$ prim mit 
\begin{enumerate}[(i)]
	\item $p \nmid a_n$
	\item $p \mid a_i$ für $i <n$
	\item $p^2 \nmid a_0$
\end{enumerate}
Dann ist $f$ irreduzibel in $R[X]$ und damit auch in $K[X]$.
\minisec{Beweis}
Sei $f = g \cdot h$ mit $g,h \in R[X]$. Sei dabei $g = c_k X^k + \ldots + c_0$, $h = d_l X^l + \ldots + d_0$ mit $\grad (g)= k$, $\grad(h)=l$. Sei 
$\Phi : R[X] \to \nicefrac{R}{R_p}[X]$ der von $R \to \nicefrac{R}{Rp}$ induzierte Ringhomomorphismus. Für $\alpha \in R$ setze 
$\overline{\alpha} := \alpha + R_p \in \nicefrac{R}{Rp} $. Es ist $\Phi (f) = \overline{a_n} X^n $ wegen ii). Da $\Phi (f) = \Phi (g) \cdot \Phi (h)$ folgt\footnote{ist nur richtig, wenn $\nicefrac{R}{Rp}[X]$ faktoriell, also wenn $\nicefrac{R}{Rp}$ faktoriell ist. Wenn nicht betrachte $K := Q(\nicefrac{R}{Rp})$ (nullteilerfrei,
da $Rp$ Primideal) und $\nicefrac{R}{Rp}[X] \subset K[X]$ und benutze die eindeutige Primfaktorzerlegung in $K[X]$.}
$\Phi (g) = \overline{c_k}X^k $ und $\Phi (h)= \overline{d_l}X^l $. 

Angenommen: $k \ge 1$ und $l \ge 1$. Dann folgt $\overline{c_o}=0 $ und $\overline{d_0}=0 $. Dann teilt $p$ sowohl $c_0$ als auch $d_0$. Aber 
$p^2 \nmid a_0 = c_0 \cdot d_0$. \light

Es folgt $k=0$ oder $l=0$. O.B.d.A. $k=0$. Also $g \in R$. Wegen $f= g \cdot h$ und $f$ primitiv folgt $g \in R^\times$. Damit ist $f$ irreduzibel. \bewende
% subsection 95 (end)

\subsection[{Beispiel: In $\mathds{Q}[X]$ existieren irreduzible Polynome von beliebigem Grad}]{Beispiel} % (fold)
\label{sub:96}
$X^n +p \in \mathds{Q}[X]$ für $p$ eine Primzahl ist irreduzibel. Insbesondere gibt es in $\mathds{Q}[X]$ irreduzible Polynome von beliebigem Grad $\ge 1$.
% subsection 96 (end)

\subsection[Beispiel: Körper der rationalen Funktionen]{Beispiel} % (fold)
\label{sub:97}
Sei $k$ ein Körper und $k(t) := Q(k[t])$ der \Index{Körper der rationalen Funktionen} in einer Variablen $t$ über $k$. Dann ist $f := X^n - t \in k(t)[X]$ irreduzibel.
Es ist nämlich $k[t]$ faktoriell und $t \in k[t]$ prim. Also können wir das Kriterium von Eisenstein anwenden.
% subsection 97 (end)

\subsection[Beispiel: Wenn $f$ irreduzibel, dann ist auch $f(X+1)$ irreduzibel]{Beispiel} % (fold)
\label{sub:98} 
Sei $p$ eine Primzahl. Behauptung: $f = X^{p-1}+ X^{p-2} + \ldots + X + 1 = \frac{X^p -1}{X-1} \in \mathds{Q}[X]$ ist irreduzibel. Wir zeigen 
$f(X+1) = \frac{(X+1)^p -1}{(X+1) -1} = X^{p-1} + \binom{p}{1} X^{p-2} + \ldots + \binom{p}{p-1}$ ist irreduzibel. Nun gilt $p \mid \binom{p}{r}$ für $r <p$ und 
$\binom{p}{1} =p$. Nach Eisenstein ist das irreduzibel.
% subsection 98 (end)

\subsection[Lemma: $f$ irreduzibel $\iff$ $f(X+1)$ irreduzibel]{Lemma} % (fold)
\label{sub:99}
Sei $f = a_n X^n + a_{n-1} X^{n-1} + \ldots + a_0 \in K[X]$. Dann ist $f$ genau dann irreduzibel, wenn 
\[
	f(X+1) = a_n (X+1)^n + a_{n-1} (X+1)^{n-1} + \ldots + a_1 (X+1) +a_0
\]
in $K[X]$ irreduzibel ist.
\minisec{Beweis}
Es ist $f \mapsto f(X+1)$ ein Ringisomorphismus von $K[X]$ (Das Inverse ist $f \mapsto f(X-1)$). Also ist $f$ genau dann irreduzibel, wenn $f(X+1)$ irreduzibel ist.
\bewende
% subsection 99 (end)
% section 9 (end)
\newpage


\section{Konstruktion mit Zirkel und Lineal} % (fold)
\label{sec:10}
\subsection[Beispiel: Mögliche Konstruktionen mit Zirkel und Lineal]{Beispiel} % (fold)
\label{sub:101}
\begin{enumerate}[1)]
	\item Sei en $P,Q \in \mathds{R}^2$. Dann lässt sich der Mittelpunkt $M := \frac{P+Q}{2} $ zwischen $P$ und $Q$ mit Zirkel und Lineal konstruieren:
	\[
		\begin{tikzpicture}[scale=0.7]
			%\draw[help lines] (0,0) grid (4,4);
			\coordinate (A) at (0,2);
			\coordinate (B) at (4,2);
			\draw[thick] (-0.3,2) -- (4.3,2);
			\draw[fill=Tomato2] (A) circle[radius=0.07];
			\node[color=Tomato2, below] at (A) {$P$};
			\draw[fill=SeaGreen3] (B) circle[radius=0.07];
			\node[color=SeaGreen4, below] at (B) {$Q$};
			\begin{scope}
			  \clip(1,0)--(3,0)--(3,4) -- (1,4);
			  \draw[color=Tomato2, ultra thick] (A) circle[radius=2.3];
			  \draw[color=SeaGreen3, ultra thick] (B) circle[radius=2.3];
			\end{scope}
			\draw[thick, color=SlateGray4] (2,0) -- (2,4);
		\end{tikzpicture}
	\]
	\item Winkelhalbierung
	\[
		\begin{tikzpicture}[scale=0.7]
			%\draw[help lines] (0,0) grid (4,4);
			\coordinate (O) at (0,0);
			\coordinate (A) at (15:3);
			\coordinate (B) at (67:3);
			\draw[color=SlateGray4, very thick] (O) -- (A) -- ++ (15:1.5);
			\draw[color=SlateGray4, very thick] (O) -- (B) -- ++ (67:1.5);
			\begin{scope}
			  \clip(0,0)--(4,0)--(4,4) -- (0,4);
			  \draw[color=Tomato2, ultra thick] (A) circle[radius=2];
			  \draw[color=SeaGreen3, ultra thick] (B) circle[radius=2];
			\end{scope}
			\draw[fill=Tomato2] (A) circle[radius=0.07cm];
			\draw[fill=SeaGreen3] (B) circle[radius=0.07cm];
			\draw[dashed, ultra thick] (0,0) -- (41:5);
		\end{tikzpicture}
	\]
	\item Drittelung einer Strecke:
	\[
		\begin{tikzpicture}
			%\draw[help lines] (0,0) grid (6,3);
			\coordinate (A) at (45:1);
			\coordinate (B) at (45:2);
			\coordinate (C) at (45:3);
			\draw[very thick] (0,0) -- (C);
			\draw[ultra thick, color=SeaGreen3] (0,0) -- (6,0);
			\draw[color=DodgerBlue2, very thick] (C) -- (6,0);
			\draw[color=DodgerBlue2, very thick] (B) -- (4,0);
			\draw[color=DodgerBlue2, very thick] (A) -- (2,0);
			\begin{scope}
			  \clip(45:0.5)-- ++ (1,0) -- ++ (0,1) -- ++ (-1,0);
			  \draw[color=Tomato2, ultra thick] (0,0) circle[radius=1];
			\end{scope}
			\begin{scope}
			  \clip(45:1.5)-- ++ (1,0) -- ++ (0,1) -- ++ (-1,0);
			  \draw[color=Tomato2, ultra thick] (A) circle[radius=1];
			\end{scope}
			\begin{scope}
			  \clip(45:2.5)-- ++ (1,0) -- ++ (0,1) -- ++ (-1,0);
			  \draw[color=Tomato2, ultra thick] (B) circle[radius=1];
			\end{scope}
			
		\end{tikzpicture}
	\]
\end{enumerate}
% subsection 101 (end)

\subsection{Konstruktionsprobleme} % (fold)
\label{sub:102}
\begin{enumerate}[(1)]
	\item Winkeldrittellung
	\item Delische Problem: Würfelverdoppellung
	\item Quadratur des Kreises: $\Vol (\bigcirc) = \Vol(\square)$
	\item Konstruktion des regelmäßigen $n$-Ecks: $\triangle, \square, $Pentagon
\end{enumerate}
% subsection 102 (end)

\subsection[Definition: Aus $M \subset \mathds{R}^2$ konstruierbar]{Definition} % (fold)
\label{sub:103}
Sei $M$ eine Menge von Punkten in $\mathds{R}^2$. 
\[
	\zirkel(M) := \set[P \in \mathds{R}^2]{ P \text{ ist mit Zirkel und Lineal aus } M \text{ konstruierbar}} 
\]
Genauer: Für $M \subseteq \mathds{R}^2$ sei   
\begin{align*}
	\Gr(M) &:= \set[g \subseteq \mathds{R}^2 \text{ Gerade}]{\abs*{g \cap M}\ge 2 } \\
	\Kr(M) &:= \set[k \subseteq \mathds{R}^2 \text{ ein Kreis}]{\text{Mittelpunkt }(k) \in M , \text{ Radius }(k) = \text{ Abstand }(P,Q) \text{ mit } P,Q \in M} 
\end{align*}
Setze nun
\[
	\zirkel (M) := \set[P \in \mathds{R}^2]{
	\vcenter{\hbox{\begin{minipage}[b]{6.5cm}
		\begin{align*}
			&\exists g \not= g \in \Gr(M) : P \in g \cap g' \\
					\text{oder } &\exists k \not= k' \in \Kr(M) \text{ mit } P \in k \cap k' \\
					\text{oder } &\exists k \in \Kr(M), g \in \Gr(M) \text{ mit } P \in k \cap g
		\end{align*}
	\end{minipage}}}} 
\]
und $\zirkel^{(0)}(M) = M$. $\zirkel^{(n)}(M) := \zirkel^{(1)} \enbrace*{\zirkel^{(n-1)}(M)} $ für $n \ge 1$. Dann ist $(M) = \bigcup_{n=0}^\infty \zirkel^{(n)}(M)	$. \\
\bet{Warnung:} $g \in \Gr(M), P \in g \not\Longrightarrow P \in \zirkel(M)$, $k \in \Kr(M), P \in k \not\Longrightarrow P \in \zirkel(M)$
% subsection 103 (end)

\subsection[Bemerkung: Übersetzung der Konstruktionsprobleme in Zahlen aus $\mathds{C}$]{Bemerkung} % (fold)
\label{sub:104}
Wir ersetzen $\mathds{R}^2$ durch den Körper $\mathds{C}$. Damit ergibt sich 
\begin{description}
	\item[(3)] Quadratur des Einheitskreises: $\sqrt{\pi }  \in \zirkel(\set{0,1})$
	\item[(2)]\label{delisches} Delisches Problem: Verdoppelung des Einheitswürfels: $\sqrt[3]{2} \in \zirkel (\set{0,1}) $
	\item[(4)] Konstruktion des regelmäßigen Einheits-$n$-Ecks: $e^{\frac{2 \pi }{n} i} \in \zirkel(\{0,1\})$ \marginnote{Zentriwinkel}
	\item[(1)] Drittellung des Winkels $\varphi \in [0, 2\pi ]$:. $e^{\frac{i \varphi}{3} } \in \zirkel(\set{0,1, e^{i \varphi}} )$
\end{description}
% subsection 104 (end)

\subsection[{Proposition 1: Elemente aus $M \subseteq \mathds{C}$, mit $\{0,1\} \subseteq M $}]{Proposition 1} % (fold)
\label{sub:105}
Sei $M \subseteq \mathds{C}$ mit $\set{0,1} \subseteq M$. Dann gilt
\begin{enumerate}[i)]
	\item $i \in \zirkel(M)$
	\item $z \in \zirkel(M) \Longrightarrow -z, \overline{z}, \re(z), \im(z), \abs*{z} \in \zirkel(M)$
	\item $z \in \zirkel(M), z \not= 0 \Longrightarrow \frac{1}{z} \in \zirkel(M) $
	\item $z_1, z_2 \in \zirkel(M) \Longrightarrow z_1 +z_2, z_1 \cdot z_2 \in \zirkel(M)$
	\item $z \in \mathds{C}, z^2 \in \zirkel(M) \Longrightarrow z \in \zirkel (M)$
\end{enumerate}
\minisec{Beweis}
Übung und Lorenz
% subsection 105 (end)

\subsection[Korollar: Die konstruierbaren Zahlen bilden einen Unterkörper von $\mathds{C}$]{Korollar} % (fold)
\label{sub:106}
\begin{enumerate}[i)]
	\item $\zirkel (\set{0,1} ) \subseteq \mathds{C}$ ist ein Unterkörper.
	\item $\mathds{Q} \subseteq \zirkel(\{0,1\})$, sogar $\mathds{Q}[i] \subseteq \zirkel(\set{0,1} )$
	\item $\mathds{Q}[i] \subsetneq \zirkel(\set{0,1} ) $
\end{enumerate}
\minisec{Beweis}
Proposition 1 (auch iii) mit $\sqrt{z} \in \zirkel(\set{0,1}), \sqrt{z} \not\in \mathds{Q}[i]$
% subsection 106 (end)

\subsection[Frage nach Unterkörpern von $\mathds{C}$]{Frage} % (fold)
\label{sub:107}
Was sind Beispiele von Unterkörpern von $\mathds{C}$?
\begin{multicols}{3}
	\begin{enumerate}[(i)]
		\item $\mathds{Q} \subseteq \mathds{C}$
		\item $\mathds{R} \subseteq \mathds{C}$
		\item $\mathds{C}$
	\end{enumerate}
\end{multicols}
% subsection 107 (end)

\subsection[{Beispiel: $\mathds{Q}[i]$ ist Unterkörper von $\mathds{C}$}]{Beispiel} % (fold)
\label{sub:108}
$\mathds{Q}[i] := \set[a+ b i]{a,b \in \mathds{Q}} $ ist ein Unterkörper von $\mathds{C}$
\[
	(a+bi)^{-1} = \frac{a-bi}{a^2 + b^2} = \frac{a}{a^2 +b^2}+ \frac{b}{a^2 + b^2} i \in \mathds{Q}[i]   
\]
% subsection 108 (end)

\subsection[{Lemma: Wenn $K$ Unterkörper von $\mathds{C}$, dann ist $K[\alpha]$ auch Unterkörper von $\mathds{C}$}]{Lemma} % (fold)
\label{sub:109}
Sei $K \subseteq \mathds{C}$ ein Unterkörper. Sei $\alpha \in \mathds{C} \setminus K$ mit $\alpha^2 \in K$. Dann ist $K[\alpha] := \set[a+ b \alpha]{a,b \in K} $ ein
Unterkörper von $\mathds{C}$.
\minisec{Beweis}
Zu zeigen: $\forall 0 \not= a + b \alpha \in K[\alpha]$ ist $(a+ b \alpha)^{-1} \in K[\alpha]$. Für $x,y \in K$ gilt $(x+y \alpha) = (a+ b \alpha) ^{-1}$ genau dann, wenn
gilt 
\begin{align*}
	ax + b \alpha^2 y &= 1 \\
	 bx + ay &= 0
\end{align*}
Es ist $\det \begin{pmatrix}
	a	& b \alpha^2 \\
	b & a
\end{pmatrix}= a^2 - b^2 \alpha^2$. Ist $b=0$, so ist $a\not=0$
und $\alpha^2 - b^2 \alpha^2 \not= 0$. ist $b \not= 0$, so ist auch $a^2 - b^2	\alpha^2 \not= 0$, denn sonst $\frac{a^2}{b^2} = \alpha^2 $ und damit 
$\alpha = \pm \frac{a}{b} \in K$, aber $\alpha \not\in K$.  

Damit ist das Gleichungssystem lösbar und $(a+ b \alpha)$ in $K[\alpha]$ invertierbar. \bewende
% subsection 109 (end)

\subsection[Definition: Quadratische Körpererweiterung]{Definition} % (fold)
\label{sub:1010}
Wir sagen dann: 
\[
	\nicefrac{K[\alpha]}{K} 
\]
ist eine \Index{quadratische Körpererweiterung}.
% subsection 1010 (end)

\subsection[Satz: Äquivalenz zu $z \in \mathds{C}$ ist konstruierbar aus $\{0,1\}$]{Satz} % (fold)
\label{sub:1011}
Für $z \in \mathds{C}$ sind äquivalent: 
\begin{enumerate}[i)]
	\item $z \in \zirkel (\set{0,1})$
	\item Es gibt eine Kette $\mathds{Q} = K_0 \subseteq K_1 \subseteq \ldots \subseteq K_n$ von Unterkörpern von $\mathds{C}$ mit $z \in K_n$ und $\nicefrac{K_i}{K_{i-1}}$
	ist eine quadratische Körpererweiterung, also gilt für jedes $i=1,\ldots ,n : \exists \omega_i \in K_i \setminus_{i-1}$ mit $\omega_i^2 \in K_{i-1}$ und 
	$K_i = K_{i-1}[\omega_i]$.
\end{enumerate}
\minisec{Beweis}
\begin{description}
	\item[ii) $\Rightarrow $ i)] Proposition 1 (\ref{sub:105})+ Induktion nach $n$
	\item[i) $\Rightarrow $ ii)] Sei $z \in \zirkel(\set{0,1})$. Dann gibt es $P_1, \ldots , P_n \in \mathds{C}$ mit $P_n = z$, so dass $P_{i+1}$ aus 
	$M_i = \set{0,1, P_1, \ldots , P_{i}} $ durch einen Konstruktionsschritt konstruiert werden kann. Also: Es gilt 
	\begin{enumerate}[(i)]
		\item $P_{i+1} \in g \cap g', g \not= g' \in \Gr(M_i)$ oder 
		\item $P_{i+1} \in k \cap k', k \not= k' \in \Kr(M_i)$ oder 
		\item $P_{i+1} \in k \cap g, k \in \Kr(M_i), g \in \Gr(M_1)$
	\end{enumerate}
	Sei $K_i$ der kleinste Unterkörper von $\mathds{C}$, der $M_i$ enthält. \\
	\uline{Behauptung (*):} Ist $K_{i+1} \supsetneq K_i$, so gibt es $\omega_{i+1} \in K_{i+1} \setminus K_i$ mit $\omega_{i+1}^2 \in K_i$ und $K_{i+1} = K_i[\omega_{i+1}]$
	
	Klar: (*) $\Rightarrow $ ii) Beweis modulo (*) \bewende
\end{description}
% subsection 1011 (end)

\subsection[Beispiel: Das regelmäßige 5-Eck ist konstruierbar]{Beispiel} % (fold)
\label{sub:1012}
\[
	\omega := e^{\frac{2 \pi i}{5} } 
\]
(Also $\omega \in \zirkel (\set{0,1} ) \iff$ regelmäßige 5-Eck ist konstruierbar)\\
Es ist $\omega^5=1$. Also ist $\omega$ Nullstelle von $X^5-1$. Es ist 
\[
	X^5-1 =  (X-1)(X^4+X^3+X^2+X+1)
\]
Insbesondere $(\omega-1)(\omega^4+ \omega^3+\omega^2+ \omega+1) =0$. Da $\omega \not= 1$, ist also $\omega^4+ \omega^3+\omega^2+\omega +1 =0$.\marginnote{$|\cdot \omega^{-2}$}
Sei $z:= \omega + \omega ^{-1}$. Dann folgt 
\[
	\omega^2 + \omega +1 + \omega ^{-1} + \omega ^{-2}=0 \tag{$\#$}
\]
Es ist $z^2 = \omega^2+ 2 + \omega ^{-2}$. Mit $(\#)$ folgt $z^2+z-1=0$. Also $z \in \set{- \frac{1}{2} \pm \frac{1}{2} \sqrt{5}}$. Daher $z \in \mathds{Q}[\sqrt{5}]$.
Weiter ist $\omega^2- z \omega = \omega^2- \omega^2-1 = -1$. Also 
\[
	\big(\omega- \frac{z}{2} \big)^2 = \omega^2 - z \omega+ \frac{z^2}{4} = -1 + \frac{z^2}{4} \in \mathds{Q}[\sqrt{5}]  
\]
Es ist also $\mathds{Q} \subseteq \mathds{Q}[\sqrt{5}] \subseteq \mathds{Q}[\sqrt{5}][\omega- \frac{z}{2} ]$ und $\omega \in \mathds{Q}[\sqrt{5}][\omega - \frac{z}{2}]$.\\
$\Longrightarrow$ Das regelmäßige $5$-Eck ist mit Zirkel und Lineal konstruierbar. \bewende
% subsection 1012 (end)
% section 10 (end)
\newpage

\section{Algebraische Körpererweiterungen} % (fold)
\label{sec:11}

\subsection[Definition: Körpererweiterung]{Definition} % (fold)
\label{sub:11.1}
Sei $K$ ein Unterkörper des Körpers $L$. Dann sagen wir $L$\marginnote{Die Notation gleicht leider der von Quotientenringen!} ist eine \Index{Körpererweiterung} (KE) von $K$ oder kurz $\nicefrac{L}{K}$ ist eine Körpererweiterung. 
% subsection 111 (end)

\subsection[Beispiele für Körpererweiterungen]{Beispiele} % (fold)
\label{sub:11.2}
$\nicefrac{\mathds{Q}[i]}{\mathds{Q}}$, $\nicefrac{\mathds{Q}[\sqrt{2}]}{\mathds{Q}}$, $\nicefrac{\mathds{C}}{\mathds{Q}}$, $\nicefrac{\mathds{R}}{\mathds{Q}}$,
$\nicefrac{\mathds{C}}{\mathds{R}}$, $\nicefrac{\mathds{R}}{\mathds{Q}[\sqrt{2} ]}$, \ldots 
% subsection 112 (end)

\subsection[Definition: Zwischenkörper]{Definition} % (fold)
\label{sub:11.3}
Ein \Index{Zwischenkörper} einer Körpererweiterung $\nicefrac{L}{K}$ ist ein Unterkörper $E$ von $L$ der $K$ enthält:
\[
	K \subseteq E \subseteq L
\] 
% subsection 113 (end)

\subsection[Bemerkung über den Schnitt von Zwischenkörpern]{Bemerkung} % (fold)
\label{sub:11.4}
Sind $E$ und $E'$ Zwischenkörper von $\nicefrac{L}{K}$, dann ist auch $E \cap E'$ ein Zwischenkörper von $\nicefrac{L}{K}$.
% subsection 11.4 (end)

\subsection[Definition: Von Teilmenge erzeugter kleinster Zwischenkörper]{Definition} % (fold)
\label{sub:11.5}
Sei $\nicefrac{L}{K}$ eine Körpererweiterung und $M \subseteq L$ eine Teilmenge. Dann gibt es einen kleinsten Zwischenkörper von $\nicefrac{L}{K}$, der $M$ enthält, nämlich
der Durchschnitt aller Zwischenkörper von $\nicefrac{L}{K}$, die $M$ enthalten. Wir bezeichnen diesen Zwischenkörper mit $K(M)$. $K(M)$ heißt von $M$ erzeugt. Wenn 
$M= \set{\alpha_1, \ldots ,\alpha_n} $, dann setzen wir $K(\alpha_1, \ldots , \alpha_2) := K(\set{\alpha_1, \ldots , \alpha_n} )$
% subsection 11.5 (end)

\subsection[Beispiel eines erzeugten Zwischenkörpers aus dem vorherigen Kapitel]{Beispiel} % (fold)
\label{sub:11.6}
Ist $\alpha \in \mathds{C}$ mit $\alpha^2 \in K$, so ist 
\[
	K(\alpha) = K[\alpha] = \set[a+ \alpha b]{a,b \in K} 
\]
% subsection 11.6 (end)

\subsection[Bemerkung: KEs lassen sich als Vektorraum über dem Basiskörper auffassen]{Bemerkung} % (fold)
\label{sub:11.7}
Ist $\nicefrac{L}{K}$ eine Körpererweiterung, so ist $L$ insbesondere ein $K$-Vektorraum. 
% subsection 11.7 (end)

\subsection[Definition: Grad einer Körpererweiterung]{Definition} % (fold)
\label{sub:11.8}
Sei $\nicefrac{L}{K}$ eine Körpererweiterung. $[L:K] := \dim_K L$ heißt der \bet{Grad von $\nicefrac{L}{K}$}\index{Körpererweiterung!Grad}. Ist $[L:K] < \infty$, so heißt
$\nicefrac{L}{K}$ endlich, sonst unendlich.
% subsection 11.8 (end)

\subsection[{Beispiel: Grad von $\nicefrac{\mathds{Q}[i]}{\mathds{Q}}$}]{Beispiel} % (fold)
\label{sub:11.9}
$\set{1,i}$ ist eine $\mathds{Q}$-Basis von $\mathds{Q}[i]$. Also $\big[\mathds{Q}[i] :Q\big] =2$. Jede quadratische Körpererweiterung hat Grad $2$. Die Umkehrung ist auch 
richtig (Übung).
% subsection 11.9 (end)

\subsection{Gradsatz} % (fold)
\label{sub:1110}
Sei $E$ ein Zwischenkörper von $\nicefrac{L}{K}$, dann gilt: 
\[
	[L : K] = [L: E] \cdot [E :K]
\]
\minisec{Beweis}
Sei $B$ eine $K$-Basis von $E$ und $A$ eine $E$-Basis von $L$. Es genügt zu zeigen, dass $A \cdot B := \set[a \cdot b]{a \in A, b \in B} $ eine $K$-Basis von $L$ ist.
\begin{description}
	\item[$A \cdot B$ ist $K$-EZS:] Sei $l \in L$. Da $A$ ein $E$-EZS von $L$ ist, gilt $l = \sum_{a \in A} \alpha_a \cdot a$ mit $\alpha_a \in E$, fast alle $\alpha_a=0$.
	Da $B$ ein $K$-EZS von $E$ ist, gilt
	\[
		\alpha_a = \sum_{b \in B} \beta_{b,a} \cdot b
	\]
	mit $\beta_{b,a} \in K$ (und fast alle gleich 0). Es folgt 
	\[
		l = \sum_{a \in A} \alpha_a \cdot a = \sum_{a \in A} \enbrace*{\sum_{b \in B} \beta_{b,a} \cdot b } \cdot a = \sum_{a \in A, b \in B}\beta_{b,a} \cdot b \cdot a
	\]
	\item[$A \cdot B$ ist linear unabhängig über $K$:] Sei also $\sum_{a \in A, b \in B}\alpha_{a,b} \cdot a \cdot b =0$ mit $\alpha_{a,b} \in K$, fast alle $=0$. Dann ist
	auch 
	\[
		\sum_{a \in A } \underbrace{\enbrace*{ \sum_{b \in B} \alpha_{a,b} \cdot b}}_{\in E} \cdot a = 0 
	\]
	Da $A$ linear unabhängig über $E$ folgt für alle $a \in A:$ $\sum \alpha_{a,b} \cdot b =0$. Da $B$ linear unabhängig über $K$ folgt $\alpha_{a,b}=0$ für alle 
	$a\in A, b \in B$. (Dies zeigt auch $a \cdot  b = a' \cdot b' \iff a=a' \wedge b=b'$) \bewende
\end{description}
% subsection 1110 (end)

\subsection[Beispiel: Grad der Körpererweiterung für die Konstruktion des $5$-Ecks]{Beispiel} % (fold)
\label{sub:1111}
Sei $\omega = e^{\frac{2 \pi  i}{5} }$. Am Donnerstag haben wir gezeigt: $(\omega - \frac{z}{2} )^2 \in \mathds{Q}[\sqrt{5}]$ mit $z \in \mathds{Q}[\sqrt{5}]$. Es ist 
$\mathds{Q} \subsetneq \mathds{Q}[\sqrt{5}]$, da $\sqrt{5} \not\in \mathds{Q}$. Außerdem $\mathds{Q}[\sqrt{5}] \subsetneq  \mathds{Q}[\sqrt{5}, \omega- \frac{z}{2}]$, da
$\mathds{Q}[\sqrt{5}] \subseteq \mathds{R}$ aber $\mathds{R} \not\ni \omega \in \mathds{Q}[\sqrt{5}, \omega- \frac{z}{2}]$. Es folgt 
\[
	\left[ \mathds{Q}\big[\sqrt{5}, \omega- \frac{z}{2} \big] : \mathds{Q}   \right] = 4
\]
(Es ist auch $[\mathds{Q}[\omega] : \mathds{Q}] = 4$)
% subsection 1111 (end)

\subsection[Korollar: Folgerung wenn der Grad prim ist]{Korollar} % (fold)
\label{sub:1112}
Sei $\nicefrac{L}{K}$ eine Körpererweiterung mit $[L:K]$ eine Primzahl. Dann sind $L$ und $K$ die einzigen Zwischenkörper von $\nicefrac{L}{K}$.
\minisec{Beweis}
Ist $E$ ein Zwischenkörper von $\nicefrac{L}{K}$, so gilt $p = [L :K] = [L : E] \cdot [E:K]$. Es gilt also $[L :E]=1$ oder $[E:K]=1$. Also gilt entweder $L=E$ oder $E=K$.
\bewende
% subsection 1112 (end)

\subsection[Definition: algebraisches Element und algebraische Körpererweiterung]{Definition} % (fold)
\label{sub:1113}
Sei $\nicefrac{L}{K}$ eine Körpererweiterung.
\begin{enumerate}[(i)]
	\item $\alpha \in L$ heißt \Index{algebraisch} über $K$, wenn es $f \in K[X]$, $f \not=0$ gibt mit $f(\alpha)=0$.
	\item $\nicefrac{L}{K}$ heißt \bet{algebraisch}, wenn alle $\alpha \in L$ algebraisch über $K$ sind. 
\end{enumerate}
% subsection 112 (end)

\subsection[Beispiele für algebraische Elemente aus $\nicefrac{\mathds{C}}{\mathds{Q}}$]{Beispiel} % (fold)
\label{sub:1114}
\begin{itemize}
	\item $\omega \in \mathds{C}$ ist algebraisch über $\mathds{Q}: \enspace(X^5-1)(\omega)= 0$
	\item Für alle $n,k \in \mathds{Z}$ ist $\sqrt[n]{k} \in \mathds{R} $ algebraisch über $\mathds{Q}$: $(X^n-k)(\sqrt[n]{k}  )=0$
\end{itemize}
% subsection 1114 (end)

\subsection[Fragen über algebraische Elemente]{Fragen} % (fold)
\label{sub:1115}
\begin{enumerate}[1)]
	\item Ist $\sqrt[n]{k} + \sqrt[n']{k'}$ algebraisch über $\mathds{Q}$ ?
	\item Ist die Summe und das Produkt von algebraischen Elementen wieder algebraisch?
\end{enumerate}
% subsection 1115 (end)

\subsection[Proposition: Minimalpolynom, Existenz und Eindeutigkeit]{Proposition} % (fold)
\label{sub:1116}
Sei $\nicefrac{L}{K}$ eine Körpererweiterung und $\alpha \in L$ algebraisch über $K$.
\begin{enumerate}[i)]
	\item Es gibt ein eindeutiges normiertes Polynom $p_\alpha \in K[X]$ mit $(p_\alpha) = \set[{f \in K[X]}]{f(\alpha)=0} $
	\item $p_\alpha$ ist irreduzibel und $(p_\alpha)$ ein Primideal in $K[X]$
\end{enumerate}
\minisec{Beweis}
$I := \set[{f \in K[X]}]{f(\alpha)=0} \subseteq K[X] $ ist ein Ideal ($f,g \in I: (f+g)(\alpha)= f(\alpha)+ g(\alpha)=0$ und $\varphi \in K[X], f \in I : (\varphi \cdot f)(\alpha) = \varphi(\alpha) \cdot f(\alpha)=0$). Da $K[X]$ ein Hauptidealring ist, gibt es $p_\alpha \in K[X]$ mit $I= (p_\alpha)$. $p_\alpha$ wird eindeutig durch
Normierung. $I$ ist sogar ein Primideal. Seien $f,g \in K[X]$ mit $f \cdot g \in I$. Dann $0=(f \cdot g)(\alpha) = f(\alpha) \cdot g(\alpha)$. Also $f(\alpha)=0$ oder 
$g(\alpha)=0$. Daraus folgt, dass $f \in I$ oder $g \in I$. Daher ist $p_\alpha$ prim und damit auch irreduzibel. \bewende
% subsection 1116 (end)

\subsection[Definition: Minimalpolynom]{Definition} % (fold)
\label{sub:1117}
$p_\alpha$ heißt das \Index{Minimalpolynom} von $\alpha$ über $K$. 
% subsection 1117 (end)

\subsection[Bemerkung: Das Minimalpolynom ist das kleinste Polynom mit $\alpha$ als Nullstelle]{Bemerkung} % (fold)
\label{sub:1118}
$p_\alpha$ ist das normierte Polynom von kleinstem Grad mit $p_\alpha(\alpha)=0$.
% subsection 1118 (end)

\subsection[Bemerkung: Irreduzible, normierte Polynome mit $\alpha$ als Nullstelle sind direkt $p_\alpha$]{Bemerkung} % (fold)
\label{sub:1119}
Sei $\nicefrac{L}{K}$ eine Körpererweiterung und $\alpha \in L$. Ist $p \in K[X]$ irreduzibel und normiert mit $p(\alpha)=0$, so ist $\alpha$ algebraisch über $K$ und 
$p_\alpha=p$. (Denn: $p_\alpha \mid p$, da $p_\alpha $ und $p$ irreduzibel und normiert sind, folgt $p_\alpha =p$).
% subsection 1119 (end)

\subsection[Definition: Einsetzungshomomorphismus]{Definition} % (fold)
\label{sub:1120}
Sei $\nicefrac{L}{K}$ eine Körpererweiterung und $\alpha \in L$. Dann heißt der Ringhomomorphismus $\phi_\alpha : K[X] \to L$ mit $\phi_\alpha (f) := f(\alpha)$
\Index{Einsetzungshomomorphismus}. Es ist $\Kern \phi_\alpha = \set[{f \in K[X]}]{f(\alpha)=0}$. Wir setzen 
\[
	K[\alpha] := \bild \phi_\alpha.
\]
% subsection 1120 (end)

\subsection[Bemerkung über $\Kern \Phi_\alpha$]{Bemerkung} % (fold)
\label{sub:1121}
\begin{enumerate}[(i)]
	\item $K[\alpha]$ ist Unterring von $L$, der $K$ enthält. Es gilt nach dem Homomorphiesatz (\ref{sub:65})
	\[
		K[\alpha] \cong \nicefrac{K[X]}{\Kern \phi_\alpha}
	\]
	\item $\alpha$ ist genau dann algebraisch, wenn $\Kern \phi_\alpha \not= \{0\}$. In diesem Fall ist $\Kern \phi_\alpha= (p_\alpha)$.
\end{enumerate}
% subsection 1121 (end)

\subsection[{Lemma: Zusammenhang zwischen dem Grad von $K[\alpha]$ und $p_\alpha$}]{Lemma} % (fold)
\label{sub:1122}
Sei $\alpha \in L$ algebraisch über $K$. Dann gilt $\dim_K K[\alpha] = d(p_\alpha)$
\minisec{Beweis}
Sei $p_\alpha = X^n + a_{n-1} X^{n-1} + \ldots + a_1 X + a_0 \in K[X]$. Wir zeigen: $1, \alpha, \alpha^2, \ldots , \alpha^{n-1}$ ist eine Basis von $K[\alpha]$.
\begin{description}
	\item[Lineare Unabhängigkeit:] Sei $b_{n-1} \cdot \alpha^{n-1} + b_{n-2}\cdot \alpha^{n-2}+ \ldots + b_1 \cdot \alpha + b_1 \cdot 1 =0$ mit $b_i \in K$.
	Setze $f := b_{n-1} X^{n-1}+ b_{n-2}X^{n-2} + \ldots + b_1 X + b_0 \in K[X]$. Es folgt $f(\alpha)=0$. Da $d(f) < d(p_\alpha) =n$ folgt $f=0$. Damit sind 
	$b_{n-1}=b_{n-2}= \ldots = b_1 =b_0=0$
	\item[Erzeugendensystem] Es ist $K[\alpha] = \langle \alpha^k \mid k=0,1,2, \ldots \rangle_K$.  Zu zeigen: Für $N\ge n$ ist 
	$\alpha^N \in \langle\alpha^0, \alpha^1, \ldots , \alpha^{N-1} \rangle$.
	Mit $p_\alpha(\alpha) =0$ folgt 
	\[
		\alpha^n = - a_{n-1}\cdot \alpha^{n-1} - a_{n-2}\cdot \alpha^{n-2}- \ldots  - a_1 \cdot \alpha - a_0 \cdot \alpha^0
	\]
	Daher $\alpha^N = - a_{n-1}\cdot \alpha^{N-1} - \ldots - a_1 \cdot \alpha^{N-n+1}- a_0 \cdot \alpha^{N-n} \in \langle \alpha^0, \ldots , \alpha^{N-1} \rangle$ \bewende
\end{description}
% subsection 1122 (end)

\subsection[Zusammenfassung des bisher gezeigten]{Zusammenfassung} % (fold)
\label{sub:1123}
\begin{align*}
	K(\alpha) &= \text{ der von $\alpha$ erzeugte Körper} \\
	K[\alpha] &= \text{ der von $\alpha$ erzeugte Ring} \\
	K[\alpha] &\cong \begin{cases}
		\nicefrac{K[X]}{(p_\alpha)}, &\text{ falls }\alpha \text{ algebraisch}\\
		K[X] &\text{ sonst}
	\end{cases} \\
	\dim_K K[\alpha] &= \begin{cases}
		d(p_\alpha), &\text{ falls }\alpha \text{ algebraisch}\\
		\infty ,&\text{ sonst}
	\end{cases}
\end{align*}
% subsection 1123 (end)

\subsection[{Satz: Wichtige Äquivalenzen zu \glqq\alpha$ ist algebraisch über $K$\grqq}]{Satz} % (fold)
\label{sub:1124}
Sei $\nicefrac{L}{K}$ eine Körpererweiterung und $\alpha \in L$. Dann sind äquivalent:
\begin{enumerate}[(i)]
	\item $\alpha$ ist algebraisch über $K$
	\item $\dim_K K[\alpha] < \infty$
	\item $K[\alpha] = K(\alpha)$
	\item $\nicefrac{K(\alpha)}{K}$ ist endlich.
\end{enumerate}
\minisec{Beweis}
Wegen $K[\alpha] \subseteq K(\alpha)$ gilt "{}iv) $\Rightarrow$ ii)"{}. Weiter gilt "{}ii)\& iii) $\Rightarrow $iv)"{}. Es genügt daher 
i) $\Rightarrow $ ii) $\Rightarrow $ iii) $\Rightarrow $ i) zu zeigen.
\begin{description}
	\item[i) $\Rightarrow $ ii)] folgt aus $\dim_K K[\alpha]= d(p_\alpha) < \infty$
	\item[ii) $\Rightarrow $ iii)] Es gilt immer $K[\alpha] \subseteq K(\alpha)$. Es genügt daher zu zeigen, dass $K[\alpha]$ ein Körper ist. Zu zeigen:
	\[
		\forall \beta \in K[\alpha] \setminus \{0\} \text{ gilt } \beta ^{-1} \in K[\alpha]
	\]
	Betrachte $M_\beta : K[\alpha] \to K[\alpha]$ mit $M_\beta(x)= \beta \cdot x$. Zu zeigen: $1 \in \bild M_\beta$. Behauptung: $M_\beta$ ist surjektiv. Da $\dim_K K[\alpha] < \infty$ und $M_\beta : K[\alpha] \to K[\alpha]$ $K$-linear ist, genügt es zu zeigen $\Kern M_\beta = \{0\}$. Sei $x \in \Kern M_\beta$, also 
	$M_\beta (x) = \beta \cdot x = 0$. Da $\beta \not= 0$ und $K[\alpha] $ nullteilerfrei ist (da $K[\alpha] \subseteq L$), folgt $x=0$.
	\item[iii) $\Rightarrow $i)] Angenommen $\alpha$ ist nicht algebraisch. Dann gilt $K[\alpha] \cong K[X]$. Nun ist aber $K[X]$ kein Körper. $\lightning$\bewende
\end{description}
% subsection 1124 (end)

\subsection[Korollar 1: Für $\alpha$ algebraisch ist der Grad von $K(\alpha)$ gleich dem Grad von $p_\alpha$]{Korollar 1} % (fold)
\label{sub:1125}
Sei $\nicefrac{L}{K}$ eine Körpererweiterung, $\alpha \in L$ algebraisch über $K$. Dann gilt $[K(\alpha) : K] = \deg(p_\alpha)$
\minisec{Beweis}
\[
	[K(\alpha) :K] = \dim_K K(\alpha) \stackrel{\text{\ref{sub:1124}}}{=}\dim_K K[\alpha] \stackrel{\text{\ref{sub:1122}}}{=} \deg p_\alpha \bewende
\]
% subsection 1125 (end)

\subsection[Korollar 2: Endliche Körpererweiterungen sind algebraisch]{Korollar 2} % (fold)
\label{sub:1126}
Endliche Körpererweiterungen sind algebraisch.
\minisec{Beweis}
Sei $\nicefrac{L}{K}$ endlich und $\alpha \in L$. Dann ist $K(\alpha) \subseteq L$. Daher ist $[K(\alpha) :K] \le [L :K] < \infty$. Mit Satz \ref{sub:1124} folgt:
$\alpha$ ist algebraisch über $K$. \bewende
% subsection 1126 (end)

\subsection[Korollar 3: Äquivalenz zu $\alpha$ algebraisch]{Korollar 3} % (fold)
\label{sub:1127}
Sei $\nicefrac{L}{K}$ eine Körpererweiterung und $\alpha \in L$. Dann gilt:
\[
	\alpha \text{ algebraisch über }K \iff \nicefrac{K(\alpha)}{K} \text{ ist algebraisch}
\]
\minisec{Beweis}
\begin{description}
	\item["$\Leftarrow$"{}:] ist klar, da $\alpha \in K(\alpha)$
	\item["$\Rightarrow $":] Sei $\alpha \in L$ algebraisch über $K$. Mit Satz \ref{sub:1124} folgt $\nicefrac{K(\alpha)}{K}$ ist endlich. Mit Korollar 2 (\ref{sub:1126}) 
	folgt $\nicefrac{K(\alpha)}{K}$ ist algebraisch. \bewende
\end{description}
% subsection 1127 (end)

\subsection[Korollar 4: Summen, Produkte und Inverse algebraische Elemente sind wieder algebraisch]{Korollar 4} % (fold)
\label{sub:1128}
Sei $\nicefrac{L}{K}$ eine Körpererweiterung. Seien $\alpha, \beta \in L$ algebraisch über $K$. Dann ist $\nicefrac{K(\alpha, \beta)}{K}$ endlich und damit algebraisch.
Insbesondere sind $-\alpha, \alpha+\beta, \alpha \cdot \beta$ algebraisch über $K$. Ist $\alpha\not= 0$, so ist auch $\alpha ^{-1}$ algebraisch über $K$.
\minisec{Beweis}
Mit der Gradformel folgt:
\[
	[K(\alpha, \beta) : K] = [K(\alpha, \beta) : K(\alpha)] \cdot [K(\alpha) :K]
\]
Da $\alpha$ algebraisch über $K$ ist, ist $[K(\alpha) :K] < \infty$. Es ist $K(\alpha, \beta)= K(\alpha)(\beta)$. Da $\beta$ algebraisch über $K$ ist, ist $\beta$ auch
algebraisch über $K(\alpha)$. Daher ist $\nicefrac{K(\alpha)(\beta)}{K(\alpha)}$ endlich. \bewende
% subsection 1128 (end)

\subsection[Bemerkung: Die Reihenfolge beim Erzeugen von Körpererweiterung ist irrelevant]{Bemerkung} % (fold)
\label{sub:1129}
Sei $\nicefrac{L}{K}$ eine Körpererweiterung und $A,B \subseteq L$. Dann gilt $K(A \cup B) = K(A)(B) = K(B)(A)$
\minisec{Beweis}
\begin{description}
	\item["$\supseteq$":] $K(A \cup B)$ ist ein Zwischenkörper von $\nicefrac{L}{K(A)}$ da $K(A) \subseteq K(A \cup B)$. Er enthält auch $B$. Also 
	$K(A)(B) \subseteq K(A \cup B)$, da $K(A)(B)$ der kleinster Zwischenkörper von $\nicefrac{L}{K(A)}$ ist, der $B$ enthält.
	\item["$\subseteq $":] Es ist $K(A \cup B)$ der kleinster Zwischenkörper von $\nicefrac{L}{K}$, der $A \cup B$ enthält. Da $A \cup B \subseteq K(A)(B)$ und
	$K(A)(B)$. Zu $K$ von $\nicefrac{L}{K}$ folgt "$\subseteq $"{}. 
\end{description}
% subsection 1129 (end)

\subsection[Korollar 5: Der algebraische Abschluss ist ein Zwischenkörper]{Korollar 5} % (fold)
\label{sub:1130}
Sei $\nicefrac{L}{K}$ eine Körpererweiterung. Dann ist 
\[
	\overline{K}^L := \set[\alpha \in L]{\alpha \text{ ist algebraisch über }K}  
\]
ein Zwischenkörper von $\nicefrac{L}{K}$.
\minisec{Beweis}
Korollar 4 (\ref{sub:1128}). 
% subsection 1130 (end)

\subsection[Definition: Algebraischer Abschluss]{Definition} % (fold)
\label{sub:1131}
$\overline{K}^L $ heißt der \bet{algebraische Abschluss}\index{algebraischer Abschluss} von $K$ in $L$.
% subsection 1131 (end)

\subsection[Bemerkung: Algebraischer Abschluss von $\mathds{Q}$]{Bemerkung} % (fold)
\label{sub:1132}
Sei $\overline{Q} := \overline{Q}^\mathds{C}  $. Dann ist $\nicefrac{\overline{\mathds{Q}} }{\mathds{Q}}$ algebraisch, aber nicht endlich.
% subsection 1132 (end)

\subsection[Beispiel: Begründung warum der algebraische Abschluss von $\mathds{Q}$ nicht endlich ist]{Beispiel} % (fold)
\label{sub:1133}
$[\mathds{Q}(\sqrt[n]{2}) :\mathds{Q}] = \deg p_{\sqrt[n]{2}} = n$ da $p_{\sqrt[n]{2}} = X^n -2$, was irreduzibel nach Eisenstein ist.
% subsection 11333 (end)

\subsection[Korollar: Das Delische Problem ist nicht mit Zirkel und Lineal lösbar]{Korollar} % (fold)
\label{sub:1134}
Das Delische Problem (\ref{sub:104}) ist nicht mit Zirkel und Lineal lösbar.
\minisec{Beweis}
Zu zeigen: $\sqrt[3]{2}  \not\in \zirkel (\{0,1\})$. Angenommen doch. Dann gibt es eine Folge von Körpern $\mathds{Q}=K_0 \subseteq K_1 \subseteq K_2 \subseteq  \ldots \subseteq K_n$ mit: $\sqrt[3]{2} \in K_n  $ und $\nicefrac{K_i}{K_{i-1}}$ ist eine quadratische Erweiterung.\\
Es ist $[K_i : K_{i-1}] =2$. Mit der Gradformel folgt: $[K_n :\mathds{Q}] = 2^n$. Andererseits ist $[\mathds{Q}(\sqrt[3]{2}  ) :\mathds{Q}]=3$. Es folgt 
\[
	2^n = [K_n : \mathds{Q}] = \left[K_n : \mathds{Q}(\sqrt[3]{2}  )\right] \cdot \left[\mathds{Q}(\sqrt[3]{2}) :\mathds{Q}\right] 
	= \left[K_n : \mathds{Q}(\sqrt[3]{2}  )\right] \cdot 3
\]
$\Rightarrow [K_n : \mathds{Q}(\sqrt[3]{2}  )] = \frac{2^n}{3} \lightning$  \bewende
% subsection 1134 (end)
% section 11 (end)
\newpage
\section{Auflösung von algebraischen Gleichungen über $\mathds{Q}$} % (fold)
\label{sec:12}
\subsection[Beispiel: $pq$-Formel]{Beispiel} % (fold)
\label{sub:12.1}
Sei $f= X^2 + pX+ q \in \mathds{Q}[X]$. Für $\alpha \in \mathds{C}$ gilt dann 
\begin{align*}
	f(\alpha)=0 \iff \alpha^2+ p \alpha= -q \iff \enbrace*{\alpha+ \frac{p}{2} }^2- \frac{p^2}{4} = -q &\iff \alpha+ \frac{p}{2} = \pm \sqrt{\frac{p^2}{4}-q }\\
	&\iff \alpha= - \frac{p}{2} \pm \sqrt{\frac{p^2}{4}-q }
\end{align*}
% subsection 12.1 (end)

\subsection[Frage, ob es entsprechende Formeln auch für Polynome höheren Grades gibt]{Frage} % (fold)
\label{sub:12.2}
Was passiert mit Polynomen von Grad $=3,4,5, \ldots $?
% subsection 12.2 (end)

\subsection[Bemerkung: Formel um Polynome auf die Form für die Cardon'sche Formel zu bringen]{Bemerkung} % (fold)
\label{sub:12.3}
Sei $f= X^3+ a X^2 + bX + c \in \mathds{Q}[X]$. Dann ist $f \enbrace*{X- \frac{a}{3} } = X^3+pX+q $ mit $p,q \in \mathds{Q}$.
% subsection 12.3 (end)

\subsection{Satz (Cardano'sche Formel)} % (fold)
\label{sub:12.4}
Sei $f= X^3+ pX+q \in \mathds{Q}[X]$. Sei $\xi := e^{2 \pi  i/3}$. Sei $\alpha \in \mathds{C}$ mit $\alpha^2= \enbrace*{\frac{p}{3} }^3+ \enbrace*{\frac{q}{2} }^2  $.
Seien $u,v \in \mathds{C}$ mit $u^3= - \frac{q}{2}+ \alpha $, $v^3= - \frac{q}{2}- \alpha $, $u \cdot v= - \frac{1}{3}p $. Dann sind
\[
	x_1 := u+v \qquad x_2 := \xi^2 u + \xi v \qquad x_3 := \xi u + \xi^2 v
\]
die Nullstellen von $f$.
\minisec{Beweis}
Bosch Seite 274 \bewende
% subsection 12.4 (end)

\subsection[Bemerkung: Es gibt auch Formeln für Polynome mit dem Grad 4]{Bemerkung} % (fold)
\label{sub:12.5}
Es gibt auch Formeln für Polynome vom Grad $4$. (Bosch Seite 278).
% subsection 12.5 (end)

\subsection[Frage, ob alle algebraischen Elemente durch Grundrechenarten+Wurzeln darstellbar]{Frage} % (fold)
\label{sub:12.6}
Sei $\alpha \in \mathds{C}$ algebraisch über $\mathds{Q}$. Lässt sich dann $\alpha$ mittels $+,-, \cdot , (\,)^{-1}, \sqrt[n]{(\,)}  $ durch Elemente aus $\mathds{Q}$ ausdrücken?\\
Beispiel: 
\[
	\frac{\sqrt[3]{7}- \sqrt[5]{\sqrt[2]{11} + \frac{5}{2}  }    }{\sqrt[p]{101}  } 
\]	

% subsection 12.6 (end)

\subsection[Definition: Radikalerweiterung]{Definition} % (fold)
\label{sub:12.7}
Eine Körpererweiterung $\nicefrac{L}{K}$ heißt eine \Index{Radikalerweiterung}, falls es $\alpha \in L$ und $n \in \mathds{N}^*$ gibt mit 
\begin{multicols}{2}
	\begin{enumerate}[a)]
		\item $\alpha^n \in K$
		\item $L= K(\alpha)$
	\end{enumerate}
\end{multicols}
% subsection 12.7 (end)

\subsection[Bemerkung für Körper mit $\Char K >0$]{Bemerkung} % (fold)
\label{sub:12.8}
Ist $\Char K >0$ (also $p \cdot 1 = 0$ für ein $p \ge 1$) so fast man die obige Definition üblicherweise etwas weiter.
% subsection 12.8 (end)

\subsection[Definition: Durch Radikale auflösbar]{Definition} % (fold)
\label{sub:12.9}
Sei $f \in \mathds{Q}[X]$. Dann heißt $f$ \Index{durch Radikale auflösbar}, falls es eine Kette von Zwischenkörpern  $\mathds{Q}= K_0 \subseteq K_1 \subseteq  \ldots  \subseteq K_n \subseteq \mathds{C}$ gibt mit: 
\begin{enumerate}[(i)]
	\item $f$ zerfällt über $K_n$ in Linearfaktoren ($\Leftrightarrow K_n$ enthält alle komplexen Nullstellen von $f$)
	\item Für $i=1, \ldots ,n$ ist $\nicefrac{K_i}{K_{i-1}}$ eine Radikalerweiterung.
\end{enumerate}
Ein Unterkörper $E \subseteq \mathds{C}$ heißt durch Radikale auflösbar, falls es eine Kette von Zwischenkörpern $\mathds{Q}= K_0 \subseteq K_1 \subseteq \ldots \subseteq K_n \subseteq \mathds{C}$ gibt mit 
\begin{multicols}{2}
	\begin{enumerate}[(i)]
		\item $E \subseteq K_n$
		\item wie vorher
	\end{enumerate}
\end{multicols}
% subsection 12.9 (end)

\subsection[Beispiele für durch Radikale auflösbare Polynome]{Beispiel} % (fold)
\label{sub:1210}
\begin{enumerate}[(i)]
	\item Jedes Polynom von Grad $2$ ist durch Radikale auflösbar: Ist $f=X^2+ pX +q$ so enthält $\mathds{Q}(\sqrt{p^2- 4q}  )$ alle Nullstellen von $f$
	\item Sei $f = X^3+ pX+q \in \mathds{Q}[X]$, so enthält $\mathds{Q}(\xi, \alpha, u)$ mit $\xi= e^{2 \pi i /3}$, 
	$\alpha^2= \enbrace*{\frac{p}{3} }^3+ \enbrace*{\frac{p}{2} }^2 $, $u^3= -\frac{p}{2}+ \alpha $ alle Nullstellen von $f$
	\[
		\mathds{Q}\subseteq \mathds{Q}(\xi) \subseteq  \mathds{Q}(\xi, \alpha) \subseteq \mathds{Q}(\xi,\alpha,u)
	\] 
	ist eine Kette von Radikalerweiterungen.
\end{enumerate}
Ausblick: Wir werden sehen, dass es Polynome in $\mathds{Q}[X]$ gibt, die nicht durch Radikale auflösbar sind.
% subsection 1210 (end)

\subsection[Definition: auflösbare Gruppe]{Definition} % (fold)
\label{sub:1211}
Eine Gruppe $G$ heißt \bet{auflösbar}\index{auflösbar!Gruppe}, wenn es eine Folge von von normalen Untergruppen\marginnote{Notation: $N \unlhd G$ falls $N \le G$ eine normale Untergruppe ist}
\[
	\{e\} = N_0 \unlhd N_1 \unlhd N_2 \unlhd \ldots \unlhd N_n = G
\]
gibt, so dass für $i=1, \ldots ,n$ $\nicefrac{N_i}{N_{i-1}}$ abelsch ist.
% subsection 1211 (end)

\subsection[Lemma: Eigenschaft von auflösbaren endlichen Gruppen]{Lemma} % (fold)
\label{sub:1212}
Sei $G$ eine endliche Gruppe. Ist $G$ auflösbar, so gibt es eine Folge 
\[
	\{e\} = N_0 \unlhd N_1 \unlhd \ldots \unlhd N_n = G
\]
von normalen Untergruppen, so dass für $i=1,\ldots ,n$ $\nicefrac{N_i}{N_{i-1}}$ zyklisch von Primzahlordnung ist.
\minisec{Beweis}
Übung.
% subsection 1212 (end)

\subsection[Lemma:Äquivalente Aussage zu $G$ ist auflösbar]{Lemma} % (fold)
\label{sub:1213}
Sei $N \unlhd G$ eine normale Untergruppe. Dann gilt
\[
	G \text{ auflösbar } \iff N \text{ und } \nicefrac{G}{N} \text{ sind auflösbar}
\]
\minisec{Beweis}
Übung.
% subsection 1231 (end)

\subsection[Korollar: $p$-Gruppen sind auflösbar]{Korollar} % (fold)
\label{sub:1214}
$p$-Gruppen sind auflösbar.
\minisec{Beweis}
Durch Induktion nach der Ordnung der Gruppe. Für die triviale Gruppe ist die Aussage sicher richtig. Sei $P$ eine $p$-Gruppe. Das Zentrum $Z$ von $P$ ist eine abelsche 
normale Untergruppe (siehe \ref{sub:317}). Da $P$ eine $p$-Gruppe ist, ist $Z\not= \{e\}$. Weiter ist $\abs*{\nicefrac{P}{Z}} < \abs*{P}  $. Nach Induktionsannahme ist
$\nicefrac{P}{Z}$ auflösbar. Da $Z$ als Zentrum abelsch und damit auflösbar ist, folgt mit dem Lemma \ref{sub:1213}: $P$ ist auflösbar. \bewende
% subsection 1214 (end)
% section 12 (end)
\newpage
\section{Primkörper} % (fold)
\label{sec:13}
\subsection[Bezeichnung: Natürliche Zahl in einem beliebigen Körper]{Bezeichnung} % (fold)
\label{sub:131}
Sei $K$ ein Körper. Dann setzen wir
\[
	n_K := \underbrace{1_K + \ldots + 1_K}_{n \text{-mal}} \tag{oft auch kurz $n=n_K$}
\]
Für eine Primzahl $p$ setzen wir $\mathds{F}_p := \nicefrac{\mathds{Z}}{p \mathds{Z}}$
% subsection 131 (end)

\subsection[Definition: Charakteristik eines Körpers]{Definition} % (fold)
\label{sub:132}
Sei $K$ ein Körper. Die kleinste positive Zahl $p \in \mathds{N}$, für die $p_K=0_K \in K$ ist, heißt die \Index{Charakteristik} $\Char K$ von $K$. Gibt es keine solche Zahl, so setzen wir $\Char K =0$ 
% subsection 132 (end)

\subsection[Beispiele für Charakteristiken]{Beispiel} % (fold)
\label{sub:133}
\begin{itemize}
	\item $\Char (\mathds{Q}) = \Char(\mathds{R})= \Char(\mathds{C}) =0$
	\item $\Char (\mathds{F}_p) = \Char Q(\mathds{F}_p[X])= p$
\end{itemize}
% subsection 133 (end)

\subsection[Bemerkung: Die Charakteristik ist immer 0 oder eine Primzahl]{Bemerkung} % (fold)
\label{sub:134}
Ist $(n \cdot m)_K=0$, so folgt $n_K=0$ oder $m_K=0$. Insbesondere ist $\Char K$ eine Primzahl oder gleich $0$.
% subsection 134 (end)

\subsection[Bemerkung: kleinster Unterkörper/Primkörper]{Bemerkung} % (fold)
\label{sub:135}
Sei $K$ ein Körper. Dann ist 
\[
	k := \set[\frac{n_K}{m_K} \in K]{n,m \in \mathds{Z}, m_k \not= 0} \qquad \text{ mit } (-n)_K = - n_K \text{ für } n \ge 0
\]
der kleinste Unterkörper von $K$. Es heißt der \Index{Primkörper} von $K$. Es gilt: 
\[
	k = \begin{cases}
		\mathds{Q}, &\text{ falls }\Char K =0\\
		\mathds{F}_p ,&\text{ falls } \Char K =p
	\end{cases}
\]
% subsection 135 (end)

\subsection[Bemerkung über Ringhomomorphismen zwischen Körpern]{Bemerkung} % (fold)
\label{sub:136}
Sei $\sigma : K \to L$ ein Ringhomomorphismus zwischen Körpern. Für $a \in K \setminus \{0\}$ gilt dann 
\[
	1 = \sigma(1) = \sigma(a \cdot a ^{-1}) = \sigma(a) \cdot \sigma(a ^{-1})
\]
\begin{enumerate}[1)]
	\item Es folgt $\sigma(a ^{-1})= \sigma(a) ^{-1}$. Insbesondere gibt es keinen Unterschied zwischen  "{}Ringhomomorphismus zwischen Körpern"{} und einem 
	"{}Körperhomomorphismus"{}.
	\item Es ist $\sigma(a) \not=0$. Insbesondere ist $\sigma$ injektiv.
\end{enumerate}
% subsection 136 (end)

\subsection[Lemma: Die Primkörper von homomorphen Körpern stimmen überein]{Lemma} % (fold)
\label{sub:137}
Sei $\sigma : K \to L$ ein Körperhomomorphismus. Dann stimmen die Primkörper von $K$ und $L$ überein und die Einschränkung von $\sigma$ auf die Primkörper ist die 
Identität.
\minisec{Beweis}
Es gilt $\sigma(n_K)= n_L$. Insbesondere ist $n_K=0 \iff n_L=0$. \bewende
% subsection 137 (end)
% section 13 (end)

\newpage
\section{Zerfällungskörper} % (fold)
\label{sec:14}
\subsection[Definition: Zerfällungskörper]{Definition} % (fold)
\label{sub:141}
Sei $K$ ein Körper und $f \in K[X]$. Ein Zerfällungskörper von $f$ ist ein Körper $L$ mit
\begin{enumerate}[(i)]
	\item $K \subseteq L$
	\item $f$ zerfällt über $L$ in Linearfaktoren. d.h. $f= \beta(X - \alpha_1) \cdots (X- \alpha_n)$ mit $\alpha_1, \ldots , \alpha_n  \in L, \beta \in K$
	\item $L= K(\alpha_1, \ldots , \alpha_n)$
\end{enumerate}
% subsection 141 (end)

\subsection[\protect{Beispiel: $\mathds{C}$ Zerfällungskörper von $X^2 +1 \in \mathds{R}[X]$}]{Beispiel} % (fold)
\label{sub:142}
$X^2+1 \in \mathds{R}[X]$ hat keine Nullstelle in $\mathds{R}$, zerfällt aber über $\mathds{C}$. $X^2+1= (X-i)(X+i)$. Da $\mathds{C}= \mathds{R}[i]$ ist $\mathds{C}$ ein
Zerfällungskörper für $X^2+1 \in \mathds{R}[X]$.
% subsection 142 (end)

\subsection{Satz von Kronecker} % (fold)
\label{sub:143}
Sei $K$ ein Körper und $g \in K[X]$ eine Polynom von Grad $\ge 1$. Dann gibt es eine endliche Körpererweiterung $\nicefrac{E}{K}$ in der $g$ eine Nullstelle hat.
\minisec{Beweis}
Sei $f$ ein irreduzibler Teiler von $g$. Dann ist $(f) \subseteq K[X]$ ein Primideal. Da $K[X]$ ein Hauptidealring ist, ist $(f)$ auch maximal. Daher ist 
$E:=\nicefrac{K[X]}{(f)}$ eine Körpererweiterung von $K$. Sei $\alpha := X+ (f) \in E$. Dann gilt:
\[
	g(\alpha) \stackrel{\text{\ref{sub:144}}}{=} g+ (f) = 0  \in E \tag{da $f \mid g$}
\]
Es ist $[E:K]= \deg(f) < \infty$. \bewende
% subsection 143 (end)

\subsection[Lemma: Hilfslemma für den Satz von Kronecker]{Lemma} % (fold)
\label{sub:144}
Seien $f,g \in K[X]$, $R:= \nicefrac{K[X]}{(f)}$. Dann gilt für $\alpha := X+ (f) \in R$ 
\[
	g(\alpha) = g+ (f)
\]
\minisec{Beweis}
Es ist $\alpha^i = (X + (f))^i = X^i + (f)$. Ist $g= \sum_{i=0}^{n} a_i X^i$, so folgt 
\[
	g(\alpha)= \sum_{i=1}^{n} a_i \alpha^i = \sum_{i=1}^{n} a_i X^i + (f)= g + (f) \bewende
\]
% subsection 144 (end)

\subsection{Korollar (Existenz von Zerfällungskörpern)} % (fold)
\label{sub:145}
Sei $f \in K[X]$. Dann gibt es einen Zerfällungskörper für $f$.
\minisec{Beweis}
Per Induktion nach $n := \deg f$. Ist $n=0$ so ist nichts zu zeigen.\\
\uline{Induktionsschritt: $n-1 \mapsto n$:} Sei $n= \deg f$. nach Kronecker gibt es eine Körpererweiterung $\nicefrac{L}{K}$ in der $f$ eine Nullstelle $\alpha \in L$ hat.
Also $f= \tilde f \cdot (X-\alpha) \in L[X]$ mit $\deg \tilde f = n-1$. Nach Induktionsannahme gibt es einen Zerfällungskörper $\nicefrac{\tilde L}{L}$ für $\tilde f$.
Dann zerfällt auch $f= (X- \alpha) \cdot \beta (X- \alpha_1) \cdot \ldots \cdot (X- \alpha_{n-1})$ über $\tilde L$. Nun ist 
$Z := K(\alpha, \alpha_1, \ldots , \alpha_n) \subseteq \tilde L$ ein Zerfällungskörper von $f$. \bewende
% subsection 145 (end)

\subsection[Definition: $K$-Homomorphismus]{Definition} % (fold)
\label{sub:146}
Seien $\nicefrac{L_1}{K}$ und $\nicefrac{L_2}{K}$ Körpererweiterungen. Ein \bet{$K$-Homomorphismus}\index{K-Homomorphismus@$K$-Homomorphismus} $\sigma : \nicefrac{L_1}{K} \to \nicefrac{L_2}{K}$ ist ein 
Ringhomomorphismus $\sigma : L_1 \to L_2$ mit $\sigma|_K = \id_K$. Gibt es einen bijektiven $K$-Homomorphismus $\sigma : \nicefrac{L_1}{K} \to \nicefrac{L_2}{K}$ so
heißen $\nicefrac{L_1}{K}$ und $\nicefrac{L_2}{K}$ \bet{isomorph}. Wir schreiben auch $\nicefrac{L_1}{K} \cong \nicefrac{L_2}{K}$. 
% subsection 146 (end)

\subsection[Bemerkung: Eigenschaften von $K$-Homomorphismen]{Bemerkung} % (fold)
\label{sub:147}
\begin{enumerate}[(i)]
	\item $K$-Homomorphismen sind $K$-linear: $\alpha \in K, x \in L_1$ dann: $\sigma(\alpha \cdot x) = \sigma(\alpha) \cdot \sigma(x) = \alpha \cdot \sigma(x)$
	\item $\nicefrac{L_1}{K} \cong \nicefrac{L_2}{K} \Longrightarrow [L_1 :K ] = [L_2 : K]$.
\end{enumerate}
% subsection 147 (end)

\subsection[{Bemerkung: $\varphi :K \to L$ induziert Ringhomomorphismus $\Phi :L[X] \to L[X]$}]{Bemerkung} % (fold)
\label{sub:148}
Sei $\varphi : K \to L$ ein Ringhomomorphismus zwischen Körpern. Sei $\Phi : K[X] \to L[X]$ der durch $\varphi$ induzierte Ringhomomorphismus.
\[
	\Phi \enbrace*{\sum_{i=0}^{n} a_i X^i} = \sum_{i=0}^{n} \varphi(a_i) X^i 
\]
Wir schreiben $\Phi (f) = f^\varphi$. Für $f \in K[X]$ und $\alpha \in K$ gilt dann $\varphi\big(f(\alpha)\big) = f^\varphi \big(\varphi(\alpha)\big)$. Insbesondere
\[
	f(\alpha) =0 \iff f^\varphi \big(\varphi(\alpha)\big)=0
\]
% subsection 148 (end)

\subsection{Fortsetzungssatz} % (fold)
\label{sub:149}
\begin{wrapfigure}[6]{R}{4cm}
	\begin{tikzcd}
		L_1 & L_2 \\
		K_1(\alpha_1)\uar[hook] \arrow[dashed]{r}{\hat \sigma} & K_2(\alpha_2) \uar[hook]\\
		K_1 \rar{\sigma} \uar[hook]& K_2 \arrow[hook]{u}
	\end{tikzcd}
\end{wrapfigure}
Seien
\begin{enumerate}[(i)]
	\item $\sigma : K_1 \to K_2$ ein Isomorphismus von Körpern
	\item $\nicefrac{L_1}{K_1}$ und $\nicefrac{L_2}{K_2}$ Körpererweiterungen
	\item $f \in K_1[X]$ irreduzibel
	\item $\alpha_1 \in L_1$ eine Nullstelle von $f$, $\alpha_2 \in L_2$ eine Nullstelle von $f^\sigma$
\end{enumerate}
Behauptung: Es existiert genau eine Fortsetzung $\hat \sigma: K_1(\alpha_1) \xrightarrow{\cong} K_2(\alpha_2) $ von $\sigma$ mit $\hat \sigma(\alpha_1)= \alpha_2$. Weiter
gilt: $\hat \sigma \big(g(\alpha_1)\big) = g^\sigma(\alpha_2)$
\minisec{Beweis}
Sei o.B.d.A $f$ normiert. Dann ist auch $f^\sigma$ normiert. Da mit $\sigma$ auch die induzierte Abbildung $\psi : K_1[X] \to K_2[X], g \mapsto g^\sigma$ ein Isomorphismus
ist, ist mit $f$ auch $f^\sigma = \psi(f)$ irreduzibel. Es folgt 
\[
	p_{\alpha_1} = f \enspace \text{ und } \enspace p_{\alpha_2}= f^\sigma 
\]
Damit folgt auch
\begin{align*}
	(f) &= \Kern \enbrace*{\Phi_{\alpha_1} : K_1[X] \twoheadrightarrow K_1[\alpha_1] = K_1(\alpha_1) }\\
	(f^\sigma) &= \Kern \enbrace*{\Phi_{\alpha_2} : K_2[X] \twoheadrightarrow K_2[\alpha_2] = K_2(\alpha_2)}
\end{align*}
Wir erhalten induzierte Isomorphismen nach dem Homomorphiesatz:
\begin{align*}
	\overline{\Phi_{\alpha_1}} : \nicefrac{K_1[X]}{(f)} \xrightarrow{\cong} K_1(\alpha_1) \quad &, \quad g+(f) \mapsto g(\alpha_1) \\
	\overline{\Phi_{\alpha_2}} : \nicefrac{K_2[X]}{(f^\sigma)} \xrightarrow{\cong} K_2{\alpha_2} \quad &, \quad g+ (f^\sigma) \mapsto g(\alpha_2)     
\end{align*}
Da $\psi\big((f)\big)= \big(\psi(f)\big) = (f^\sigma)$ induziert $\psi$ einen Isomorphismus
\[
	\overline{\psi} : \nicefrac{K_1[X]}{(f)} \xrightarrow{\cong} \nicefrac{K_2[X]}{(f^\sigma)} \quad , \quad g+ (f) \mapsto g^\sigma + (f^\sigma)  
\]
Insgesamt ist 
\[
	\begin{array}{rcccccl}
		\hat \sigma : K_1(\alpha_1) &\xrightarrow{\enbrace*{\overline{\Phi_{\alpha_1} } }^{-1}} & \nicefrac{K_1[X]}{(f)} & \xrightarrow{\overline{\psi}} & \nicefrac{K_2[X]}{(f^\sigma)} & \xrightarrow{\overline{\Phi_{\alpha_2}} }& K_2(\alpha_2) \\
		   g(\alpha_1) &\longmapsto & g+ (f) &\longmapsto &g^\sigma + (f^\sigma) &\longmapsto &g^\sigma(\alpha_2)
	\end{array}
\]
Eindeutigkeit ist klar, da $K_1(\alpha_1)=K_1[\alpha_1]$. \bewende
% subsection 149 (end)

\subsection[Lemma: Anzahl der möglichen Fortsetzungen eines Isomorphismus von Körpern]{Lemma} % (fold)
\label{sub:1410}
Sei $\sigma : K_1 \to K_2$ ein Körperisomorphismus und $\nicefrac{L_1}{K_1}$, $\nicefrac{L_2}{K_2}$ Körpererweiterungen. Sei $\alpha_1 \in L_1$ algebraisch über $K_1$.
Dann gibt es höchstens $[K_1(\alpha_1) :K_1] = \deg p_{\alpha_1}$ Fortsetzungen von $\sigma$ zu Körperhomomorphismen
\[
	\tau : K_1(\alpha_1) \to L_2
\] 
\minisec{Beweis}
Ist $\tau : K_1(\alpha_1) \to L_2$ eine solche Fortsetzung, so ist $\tau(\alpha_1)$ eine Nullstelle von $p_{\alpha_1}^\sigma$. Nun ist nach dem Fortsetzungssatz $\tau$
durch $\tau(\alpha_1)$ schon eindeutig festgelegt. Da $p_{\alpha_1}^\sigma$ höchstens $\deg p_{\alpha_1}^\sigma = \deg p_{\alpha_1}$-viele Nullstellen in $L_2$ hat, gibt
es höchstens $\deg p_{\alpha_1}$ viele $\tau$`s. \bewende
% subsection 1410 (end)

\subsection[Proposition: $K$-Endomorphismen algebraischer KE sind Automorphismen]{Proposition} % (fold)
\label{sub:1411}
Sei $\sigma : \nicefrac{L}{K} \to \nicefrac{L}{K}$ ein $K$-Endomorphismus. Ist $\nicefrac{L}{K}$ algebraisch, so ist $\sigma$ bijektiv, also ein 
$K$-Automorphismus.
\minisec{Beweis}
Jeder $K$-Homomorphismus ist injektiv. Es bleibt die Surjektivität zu zeigen. Sei $\alpha \in L$. Sei 
\[
	N= \set[\beta \in L]{p_\alpha(\beta) =0} \ni \alpha 
\]
Da $\sigma$ ein $K$-Homomorphismus ist, gilt $\sigma(N) \subseteq N$. Da $\sigma$ injektiv und $N$ endlich ist, gilt $\sigma(N)=N$. Da $\alpha \in N$ gibt es $\beta\in N$
mit $\sigma(\beta)=\alpha$. \bewende
% subsection 1411 (end)

\subsection{Satz (Eindeutigkeit des Zerfällungskörpers)} % (fold)
\label{sub:1412}
Sei $\sigma : K_1 \to K_2$ ein Körperisomorphismus und $f \in K_1[X]$. Sei $\nicefrac{L_1}{K_1}$ ein Zerfällungskörper von $f$ und $\nicefrac{L_2}{K_2}$ ein Zerfällungskörper von $f^\sigma$. Dann lässt sich $\sigma : K_1 \to K_2$ zu einem Körperisomorphismus $\tau : L_1 \to L_2$ fortsetzen.
\minisec{Beweis}
Per Induktion nach $\deg f$. Für $\deg f \le 0$ ist nichts zu zeigen. Sei die Eindeutigkeit für alle Polynome von Grad $<n$ gezeigt. \\
Sei $\nicefrac{L_1}{K_1}$ Zerfällungskörper von $f$, $\nicefrac{L_2}{K_2}$ Zerfällungskörper von $f^\sigma$. Sei $g$ ein irreduzibler Teiler von $f$, der keine Nullstelle in $K_1$ hat (Gibt es keinen solchen Teiler, so zerfällt $f$ über $K_1$ und $f^\sigma$ über $K_2$, also $L_1=K_1, L_2=K_2$). Sei $\alpha_1 \in L_1$ eine Nullstelle von $g$
und $\alpha_2 \in L_2$ eine Nullstelle von $g^\sigma$. Nach dem Fortsetzungssatz (\ref{sub:149}) setzt sich $\sigma :K_1 \to K_2$ zu einen Körperisomorphismus
$\tau_0 :K_1(\alpha_1) \to K_2(\alpha_2)$ mit $\tau_0(\alpha_1)=\alpha_2$ fort. Sei nun $\tilde f \in K_1(\alpha_1)[X]$ mit $f= \tilde f \cdot (X-\alpha_1)$. Dann ist
$\nicefrac{L_1}{K_1(\alpha_1)}$ ein Zerfällungskörper von $\tilde f$ und $\nicefrac{L_2}{K_2(\alpha_2)}$ ein Zerfällungskörper von $\tilde f^{\tau_0}$.

Nach der Induktionsannahme gibt es eine Fortsetzung von $\tau_0 : K_1(\alpha_1) \xrightarrow{\cong} K_2(\alpha_2)$ zu $\tau : L_1 \xrightarrow{\cong} L_2$. Da $\tau$
eine Fortsetzung von $\tau_0$ ist und $\tau_0$ eine Fortsetzung von $\sigma$ ist, ist $\tau$ eine Fortsetzung von $\sigma$. \bewende
% subsection 1412 (end)

\subsection{Korollar (Fortsetzungssatz für Zerfällungskörper)} % (fold)
\label{sub:1413}
Sei $f \in K[X]$ und $L$ der Zerfällungskörper von $f$. Sei $E$ ein Zwischenkörper von $\nicefrac{L}{K}$ und $\sigma : \nicefrac{E}{K} \to \nicefrac{L}{K}$ ein
$K$-Homomorphismus. Dann gibt es eine Fortsetzung von $\sigma$ zu einen $K$-Automorphismus von $L$.
\minisec{Beweis}
Betrachte $\nicefrac{L}{E}$ und $\nicefrac{L}{\sigma(E)}$ als Zerfällungskörper von $f \in E[X]$ und $f^\sigma \in \sigma(E)[X]$. Die Eindeutigkeit liefert eine
Fortsetzung von $\sigma$ zu einem $K$-Automorphismus. \bewende
% subsection 1413 (end)

\subsection[Frage: Sind Zwischenkörper invariant unter Körperautomorphismen?]{Frage} % (fold)
\label{sub:1414}
Gilt immer $\sigma(E) \subseteq E$ ?

\uline{Beispiel:} $f=X^4-2 \in \mathds{Q}[X]$. Der Zerfällungskörper ist $\mathds{Q}(i, \sqrt[4]{2}  )$. Sei $E := \mathds{Q}(\sqrt[4]{2}  )$. Was für $\mathds{Q}$-Homomorphismen $\nicefrac{E}{\mathds{Q}} \to \nicefrac{\mathds{Q}(\sqrt[4]{2},i)}{\mathds{Q}}$ gibt es? \\
Die Nullstellen von $f$ sind $\sqrt[4]{2}  ,  - \sqrt[4]{2}  ,  i \sqrt[4]{2} ,  -i \sqrt[4]{2}$.	
Da $f$ irreduzibel, gibt es nach dem Fortsetzungssatz genau 4 $\mathds{Q}$-Homomorphismen $\sigma_i : \mathds{Q}(\sqrt[4]{2}  ) \to L$. Diese sind bestimmt durch
\[
	\id = \sigma_1(\sqrt[4]{2}) = \sqrt[4]{2} ,\quad  \sigma_2(\sqrt[4]{2}) = - \sqrt[4]{2} ,\quad   \sigma_3(\sqrt[4]{2})= i \sqrt[4]{2},\quad  \sigma_4(\sqrt[4]{2})= -i \sqrt[4]{2}
\]
Es ist
\[
\begin{array}{ll}
	\sigma_1\big(\mathds{Q}(\sqrt[4]{2})\big) = \mathds{Q}(\sqrt[4]{2})\qquad \qquad  &
	\sigma_2\big(\mathds{Q}(\sqrt[4]{2})\big) = \mathds{Q}(\sqrt[4]{2}) \\
	\sigma_3\big(\mathds{Q}(\sqrt[4]{2})\big) \nsubseteq \mathds{Q}(\sqrt[4]{2})\qquad  &
	\sigma_4\big(\mathds{Q}(\sqrt[4]{2})\big) \nsubseteq \mathds{Q}(\sqrt[4]{2})
\end{array}
\]
% subsection 1414 (end)

\subsection{Proposition (Invarianz von Zerfällungskörpern unter $K$-Homomorphismen)} % (fold)
\label{sub:1415}
\begin{wrapfigure}[10]{R}[1cm]{3cm}
	\begin{tikzcd}
		E \rar{\sigma}& E \\
		L \uar[hook] \rar{\sigma|_L} & L \uar[hook] \\
		K \uar[hook] \rar{\cong} &K \uar[hook]
	\end{tikzcd}
\end{wrapfigure}
Sei $f \in K[X]$, $L$ eine Zerfällungskörper von $f$. Sei $\nicefrac{E}{L}$ eine Körpererweiterung und $\sigma : \nicefrac{E}{K} \to \nicefrac{E}{K}$ ein 
$K$-Homomorphismus. Dann gilt $\sigma(L)=L$
\minisec{Beweis}
Seien $\alpha_1, \ldots , \alpha_n$ die Nullstellen von $f$ in $L$. Es gilt $L=K(\alpha_1, \ldots , \alpha_n)$ da $L$ Zerfällungskörper von $f$ ist. Da $\sigma$ ein
$K$-Homomorphismus ist, gilt für $\beta \in L:$ 
\[
	f(\beta)=0 \iff f^\sigma\big(\sigma(\beta)\big)= f\big(\sigma(\beta)\big)=0.
\]
Also gilt $\sigma(\alpha_i) \in \set{\alpha_1, \ldots , \alpha_n} \subseteq L $ für alle $i$. Es folgt $\sigma(L) \subseteq L$. Da $\nicefrac{L}{K}$
endlich ist, folgt $\sigma(L)=L$. \bewende
% subsection 1415 (end)

\subsection[Proposition: Anzahl der Fortsetzungen bei endlichen Körpererweiterungen]{Proposition} % (fold)
\label{sub:1416}
Sei $\sigma : K_1 \to K_2$ ein Körperisomorphismus,\marginnote{Erinnerung: Do. $L_1=K(\alpha_1)$}
 $\nicefrac{L_1}{K_1}$ eine endliche Körpererweiterung, $\nicefrac{L_2}{K_2}$ eine beliebige Körpererweiterung. Dann gibt
es höchstens $n:= [L_1 : K_1]$ viele Fortsetzungen $\tau_1, \ldots ,  \tau_n : L_1 \to L_2$ von $\sigma :K_1 \to K_2$.
\minisec{Beweis}
Per Induktion nach $n$. Für $n=1$ ist nichts zu zeigen. Sei also $n= [L_1 :K_1] >1$ und $\alpha \in L_1 \setminus K_1$. Wir wissen schon, dass es höchstens 
$m=[K_1(\alpha_1): K_1]$ viele verschiedene Fortsetzungen $\sigma_1, \ldots , \sigma_m : K_1(\alpha) \to L_1$ von $\sigma$ gibt.

Nach Induktionsannahme $([L_1:K_1(\alpha)] < [L_1:K_1])$ gibt es zu jedem $\sigma_i$ höchstens $\frac{n}{m} = [L_1:K_1(\alpha)]$-viele Fortsetzungen $\tau_{i,j} : L_1 \to L_2$
von $\sigma_i :K_1(\alpha) \to L_2$.  Nun sind die $\tau_{i,j}$ aber alle möglichen Fortsetzungen von $\sigma$. Da es höchstens $m \cdot \frac{n}{m}=n $ viele
$\tau_{i,j}$ gibt, folgt die Behauptung. \bewende
% subsection 1416 (end)

\subsection[Beispiel: Anzahl von von Fortsetzungen]{Beispiel} % (fold)
\label{sub:1417}
$K_1 = K_2 = \mathds{Q}$, $L_1 = \mathds{Q}\big(\sqrt[4]{2}\big)$, $L_2 = \mathds{Q}\big(\sqrt[4]{2},i\big)$. Es gibt $4$ Fortsetzungen $\sigma_1, \ldots , \sigma_4 : \mathds{Q}\big(\sqrt[4]{2}\big) \to \mathds{Q}\big(\sqrt[4]{2},i\big)$ von $\id : K_1 \xrightarrow{\cong} K_2 $. Es ist $[L_1 :K_1] = \left[\mathds{Q}(\sqrt[4]{2}) : \mathds{Q}\right] = 4.$
% subsection 1417 (end)
% section 14 (end)
\newpage
\section{Normale Körpererweiterungen} % (fold)
\label{sec:15}

\subsection[Definition: normale Körpererweiterung]{Definition} % (fold)
\label{sub:151}
Eine Körpererweiterung $\nicefrac{L}{K}$ heißt \bet{normal}\index{Körpererweiterung!normale}, wenn folgendes gilt: Sei $f \in K[X]$ irreduzibel. Besitzt $f$ eine Nullstelle
in $L$, so zerfällt $f$ über $L$ in Linearfaktoren.

(Kurz: "{}eine Nullstelle $\Longrightarrow$ alle Nullstellen"{})
% subsection 151 (end)

\subsection[Bemerkung: alternative Charakterisierung von normalen Körpererweiterungen]{Bemerkung} % (fold)
\label{sub:152}
$\nicefrac{L}{K}$ ist genau dann normal, wenn für alle algebraischen $\alpha \in L$ das Minimalpolynom $p_\alpha \in K[X]$ über $L$ in Linearfaktoren zerfällt.
% subsection 152 (end)

\subsection[Beispiel für eine normale Körpererweiterung]{Beispiel} % (fold)
\label{sub:153}
$\nicefrac{\mathds{Q}(\sqrt[4]{2})}{\mathds{Q}}$ ist nicht normal. $\nicefrac{\mathds{Q}(\sqrt[4]{2},i)}{\mathds{Q}}$ ist normal nach dem nächsten Satz.
% subsection 153 (end)

\subsection[\protect{Satz über äquivalente Aussagen zu $\nicefrac{L}{K}$ ist normal}]{Satz} % (fold)
\label{sub:154}
Sei $\nicefrac{L}{K}$ eine endliche Körpererweiterung. Dann sind äquivalent:
\begin{enumerate}[(i)]
	\item $\nicefrac{L}{K}$ ist normal
	\item $L$ ist ein Zerfällungskörper eines Polynoms $f \in K[X]$. 
	\item Sei $\nicefrac{E}{L}$ eine Körpererweiterung und $\tau : \nicefrac{E}{K} \to \nicefrac{E}{K}$ ein $K$-Homomorphismus. Dann gilt $\tau(L)=L$.
\end{enumerate}
\minisec{Beweis}
\begin{description}
	\item[(i)$\Rightarrow $(ii):] Da $\nicefrac{L}{K}$ endlich ist, gibt es $\alpha_1, \ldots , \alpha_n \in L$ mit $L=K(\alpha_1, \ldots ,\alpha_n)$. Da 
	$\nicefrac{L}{K}$ normal ist, ist $L$ der Zerfällungskörper von $p_{\alpha_1} \cdot p_{\alpha_2} \cdot \ldots \cdot p_{\alpha_n}$
	\item[(ii)$\Rightarrow $(iii):] Heute um 10:30 :D \ldots  oder auch \ref{sub:1415}
	\item[(iii)$\Rightarrow $(i):] Sei $f \in K[X]$ irreduzibel und $\alpha \in L$ eine Nullstelle von $f$. Sei $E$ ein Zerfällungskörper von $f$ über $L$. Sei 
	$\beta \in E$ eine Nullstelle von $f$. Zu zeigen: $\beta \in L$. Nach Fortsetzungssatz gibt es einen $K$-Homomorphismus $\sigma : \nicefrac{K(\alpha)}{K} \to \nicefrac{E}{K}$ mit $\sigma(\alpha)=\beta$. Nach dem Fortsetzungssatz für Zerfällungskörper besitzt $\sigma$ eine Fortsetzung zu einen $K$-Automorphismus
	$\tau : \nicefrac{E}{K} \to \nicefrac{E}{K}$. Nach Vorraussetzung folgt $\tau(L)=L$ also $\beta= \sigma(\alpha) =\tau (\alpha) \in L$. \bewende
\end{description}
% subsection 154 (end)

\subsection[Definition: Gruppe der $K$-Automorphismen]{Definition} % (fold)
\label{sub:155}
Sei $\nicefrac{L}{K}$ eine Körpererweiterung. Mit $\Aut(\nicefrac{L}{K})$\index{K-Homomorphismus@$K$-Homomorphismus!Aut(L/K)@$\Aut(\nicefrac{L}{K})$}
 bezeichnen wir die Gruppe der $K$-Automorphismen von $\nicefrac{L}{K}$.
% subsection 155 (end)

\subsection[Bemerkung: $\Aut(\nicefrac{L}{E})$ ist eine Untergruppe von $\Aut(\nicefrac{L}{K})$]{Bemerkung} % (fold)
\label{sub:156}
Sei $E$ ein Zwischenkörper von $\nicefrac{L}{K}$. Da jeder $E$-Automorphismus von $\nicefrac{L}{E}$ auch ein $K$-Automorphismus von $\nicefrac{L}{K}$ ist, ist $\Aut(\nicefrac{L}{E})$ eine Untergruppe von $\Aut(\nicefrac{L}{K})$.
% subsection 156 (end)

\subsection[Proposition: Die Automorphismengruppe normaler Zwischenkörper ist normal]{Proposition} % (fold)
\label{sub:157}
Sei $E$ ein Zwischenkörper von $\nicefrac{L}{K}$. Sei weiter $\nicefrac{E}{K}$ normal. Dann ist $\Aut(\nicefrac{L}{E})$ eine normale Untergruppe von $\Aut(\nicefrac{L}{K})$
\minisec{Beweis}
$\Aut(\nicefrac{L}{E}) = \Kern \enbrace*{\Aut(\nicefrac{L}{K}) \xrightarrow{\sigma \mapsto \sigma|_E } \Aut(\nicefrac{E}{K}) } $. Der Kern eines Gruppenhomomorphismus ist immer 
eine normale Untergruppe. \bewende
% subsection 157 (end)
% section 15 (end)
\newpage
\section{Separable Körpererweiterungen} % (fold)
\label{sec:16}

\subsection[Definition: Separable Polynome]{Definition} % (fold)
\label{sub:161}
Sei $f \in K[X]$ ein Polynom von Grad $n$ heißt \bet{separabel}\index{Polynom!separabel}, wenn $f$ im Zerfällungskörper von $f$ $n$ verschiedene Nullstellen hat.
% subsection 161 (end)

\subsection[Beispiele für Separable und nicht separable Polynome]{Beispiel} % (fold)
\label{sub:162}
\begin{enumerate}[(i)]
	\item $X^2+1 \in \mathds{Q}[X]$ ist separabel: $X^2+1= (X-i)(X+i)$
	\item $X^2+1 \in \mathds{F}_2[X]$ ist nicht separabel: $X^2+1=(X+1)^2$
\end{enumerate}
% subsection 162 (end)

\subsection[Definition: Formale Ableitung]{Definition} % (fold)
\label{sub:163}
Die $K$-lineare Abbildung 
\[
	D : K[X] \to K[X] \quad , \quad D \enbrace*{\sum_{i=0}^{n} a_i X^i} = \sum_{i=0}^{n} a_i \cdot i X^{i-1} 
\]
heißt die \Index{formale Ableitung}.\\
\uline{Warnung:} Für $f=X^2 \in \mathds{F}_2[X]$ ist $D f = 2 \cdot X =0$
% subsection 163 (end)

\subsection[Bemerkung: Leibnizregel]{Bemerkung} % (fold)
\label{sub:164}
Es gilt die \Index{Leibnizregel}:\marginnote{wurde auch in einer Übungsaufgabe im letzten Semester gezeigt}
\[
	\D(f \cdot g) = \D (f)\cdot  g + f \cdot \D(g)
\]
% subsection 164 (end)

\subsection[Lemma: Vielfachheit von Nullstellen bestimmen mittels der formalen Ableitung]{Lemma} % (fold)
\label{sub:165}
Sei $f \in K[X]$ und $\alpha \in K$. Dann sind äquivalent:
\begin{enumerate}[(1)]
	\item $\alpha$ ist mehrfache Nullstelle von $f$, also $(X-\alpha)^2 \mid f$
	\item $f(\alpha) = f'(\alpha)=0$
\end{enumerate}
\minisec{Beweis}
Schreibe $f=g(X-\alpha)^r$ mit $g(\alpha)\not= 0$. Dann ist 
\[
	\D (f) = \D (g) \cdot (X-\alpha)^r + g \cdot \D((X-\alpha)^r) = \D (g) \cdot  (X-\alpha)^r + g \cdot r(X-\alpha)^{r-1}
\]
Es folgt $f(\alpha)= f'(\alpha)=0 \iff r\ge 2$. \bewende
% subsection 165 (end)

\subsection[Korollar: Irreduzible Polynome sind separabel, wenn $\Char K=0$]{Korollar} % (fold)
\label{sub:166}
Sei $\Char K =0$ und $f$ irreduzibel. Dann hat $f$ keine mehrfachen Nullstellen und ist separabel.
\minisec{Beweis}
Da $\grad f \ge 1$ ist $f' \not=0$. Wäre $\alpha$ (im Zerfällungskörper von $f$) eine Nullstelle von $f$ und $f'$, so wäre $\alpha$ auch eine Nullstelle vom $\ggT$ von
$f$ und $f'$ (in $K[X]$). Da $\grad f' < \grad f$ und $f$ irreduzibel ist, ist dieser $\ggT$ aber $1$ und hat keine Nullstelle. \bewende\\
(Ist $h= \ggT (f,f') \in K[X]$, so gibt es $k_1, k_2 \in K[X]$ mit $h=k_1 f+ k_2 f'$)

% subsection 166 (end)

\subsection[Definition: Separable Körpererweiterung]{Definition} % (fold)
\label{sub:167}
Eine algebraische Körpererweiterung $\nicefrac{L}{K}$ heißt \bet{separabel}\index{Körpererweiterung!separable}, falls jedes irreduzible Polynom $f \in K[X]$, das eine
Nullstelle hat, separabel ist.

Kurz: "{}eine Nullstelle $\Rightarrow n$ verschiedene Nullstellen"{}
% subsection 167 (end)

\subsection[Bemerkung: Charakterisierung von separabel über die Minimalpolynome]{Bemerkung} % (fold)
\label{sub:168}
Eine algebraische Körpererweiterung $\nicefrac{L}{K}$ ist genau dann separabel, wenn alle $p_\alpha$, $ \alpha \in L$, separabel sind.
% subsection 168 (end)

\subsection[Bemerkung: Algebraische Körpererweiterungen sind separabel, wenn $\Char K =0$]{Bemerkung} % (fold)
\label{sub:169}
$\nicefrac{L}{K}$ algebraisch, $\Char K =0$ $\Longrightarrow \nicefrac{L}{K}$ ist separabel. 
% subsection 169 (end)

\subsection[Bemerkung: $\nicefrac{L}{K}$ separabel $\Rightarrow $ $\nicefrac{L}{E}, \nicefrac{E}{K}$ separabel]{Bemerkung} % (fold)
\label{sub:1610}
$\nicefrac{L}{K}$ separabel, und $E$ ist ein Zwischenkörper von $\nicefrac{L}{K}$, dann sind auch $\nicefrac{L}{E}$ und $\nicefrac{E}{K}$ separabel.
% subsection 1610 (end)

\subsection[Bemerkung: Endliche KE lassen sich zu normalen Körpererweiterungen ausbauen]{Bemerkung} % (fold)
\label{sub:1611}
Ist $\nicefrac{L}{K}$ eine endliche Körpererweiterung, so gibt es eine Körpererweiterung $\nicefrac{N}{L}$, so dass $\nicefrac{N}{K}$ normal ist. $\nicefrac{L}{K}$ endlich
$\Rightarrow \alpha_1, \ldots , \alpha_n \in L$ mit $L=K(\alpha_1, \ldots , \alpha_n)$. Sei nun $N$ der Zerfällungskörper von 
\[
	f := p_{\alpha_1} \cdot \ldots \cdot p_{\alpha_n} \in K[X] \subseteq L[X]
\]
über $L$. Dann ist $\nicefrac{N}{K}$ der Zerfällungskörper von $f$ über $K$.
% subsection 1611 (end)

\subsection[Satz: Charakterisierung von separabel über Anzahl der $K$-Homo. in normalem Oberkörper]{Satz} % (fold)
\label{sub:1612}
Sei $\nicefrac{L}{K}$ eine endliche Körpererweiterung und $\nicefrac{N}{L}$ eine Körpererweiterung so dass $\nicefrac{N}{K}$ normal ist. Dann sind äquivalent:
\begin{enumerate}[(i)]
	\item $\nicefrac{L}{K}$ separabel
	\item Es gibt $n= [L:K]$-viele verschiedene $K$-Homomorphismen $\sigma_1, \ldots , \sigma_n : \nicefrac{L}{K} \to \nicefrac{N}{K}$.
\end{enumerate}
\minisec{Beweis}
\begin{description}
	\item[(i)$\Rightarrow $(ii):] folgt aus dem Lemma \ref{sub:1613}
	\item[(ii)$\Rightarrow $(i):] Sei $\alpha \in L$. Sei $\tau : \nicefrac{L}{K} \to \nicefrac{N}{K}$ ein $K$-Homomorphismus. Dann ist auch 
	$\tau|_{K(\alpha)} : \nicefrac{K(\alpha)}{K} \to \nicefrac{N}{K}$ ein $K$-Homomorphismus. Wir wissen (Fortsetzungssatz), dass es höchstens 
	$m=[K(\alpha):K] = \deg p_\alpha$-viele $K$-Homomorphismen $\sigma_i : \nicefrac{K(\alpha)}{K} \to \nicefrac{L}{K}$ gibt. Weiter gibt es zu jedem $\sigma_i$ höchstens 
	$\frac{n}{m}=[L:K(\alpha)] $-viele Fortsetzungen $\tau_{ij} : \nicefrac{L}{K} \to \nicefrac{N}{K}$ (\ref{sub:1416}). 
	
	Da es nach (ii) $n$ $K$-Homomorphismen
	$\nicefrac{L}{K} \to \nicefrac{N}{K}$ gibt, muss es also auch $m$ verschiedene $K$-Homomorphismen 
	$\sigma_1, \ldots , \sigma_m : \nicefrac{K(\alpha)}{K} \to \nicefrac{L}{K}$ geben. Dann sind die $\sigma_i(\alpha)$ paarweise verschieden (denn 
	$\sigma_i(\alpha) = \sigma_j(\alpha) \Rightarrow  \sigma_i = \sigma_j$). Andererseits ist jedes $\sigma_i(\alpha)$ eine Nullstelle von $p^{\sigma_i}_\alpha = p_\alpha$.
	Die $\sigma_i(\alpha)$ sind also die gesuchten $m$-vielen Nullstellen von $p_\alpha$. \bewende
\end{description}
% subsection 1612 (end)

\subsection[Lemma: Hilfslemma über die Anzahl der Fortsetzungen eines Körperisomorphismus]{Lemma} % (fold)
\label{sub:1613}
Sei $\sigma : K_1 \to K_2$ eine Körperisomorphismus. Seien $\nicefrac{L_1}{K_1}$ und $\nicefrac{L_2}{K_2}$ Körpererweiterungen. Dabei gelte
\begin{enumerate}[1)]
	\item $\nicefrac{L_1}{K_1}$ ist separabel und endlich
	\item Für alle $\alpha \in L_1$ zerfällt $p^\sigma_\alpha \in K_2[X]$ über $L_2$ in Linearfaktoren.
\end{enumerate}
Dann gibt es $n=[L_1 :K_1]$-viele Fortsetzungen von $\sigma : K_1 \to K_2$ zu $\tau :L_1 \to L_2$.
\minisec{Beweis}
Per Induktion nach $n$. Für $n=1$ ist nichts zu zeigen. Sei $\alpha \in L_1 \setminus K_1$. Da $p_\alpha$ separabel ist, hat $p_\alpha$ nach Vorraussetzung 
$m= \grad p_\alpha= [K_1(\alpha):K_1]$-viele verschiedene Nullstellen $\beta_1, \ldots , \beta_m$ in $L_2$. Nach dem Fortsetzungssatz (\ref{sub:149}) gibt es zu jeder
dieser Nullstellen eine Fortsetzung $\tau_i : K_1(\alpha) \to L_2$ von $\sigma : K_1 \to K_2$ mit $\tau_i(\alpha)= \beta_i$. Für jedes $\tau_i$ gibt es nach 
Induktionsannahme $\frac{n}{m}=[L_1:K_1(\alpha)] $-viele Fortsetzungen $\tau_{ij} : L_1 \to L_2$ von $\tau_i$. Insgesamt gibt es $m \cdot \frac{n}{m} $ viele Fortsetzungen 
von $\sigma$. \bewende
% subsection 1613 (end)

\subsection[Korollar: Wenn $p_\alpha$ separabel ist, ist auch $\nicefrac{K(\alpha)}{K}$ separabel]{Korollar} % (fold)
\label{sub:1614}
Sei $\nicefrac{L}{K}$ eine Körpererweiterung. Sei $\alpha \in L$ algebraisch über $K$. Ist $p_\alpha \in K[X]$ separabel, so ist $\nicefrac{K(\alpha)}{K}$ separabel.
\minisec{Beweis}
Sei $\nicefrac{N}{K(\alpha)}$ eine Körpererweiterung, so dass $\nicefrac{N}{K}$ normal ist. Da $p_\alpha$ separabel ist, gibt es nach Fortsetzungssatz (\ref{sub:149})
$n=[K(\alpha):K]= \grad p_\alpha$-viele $K$-Homomorphismen $\nicefrac{K(\alpha)}{K} \to \nicefrac{N}{K}$. nach Satz \ref{sub:1612} ist $\nicefrac{K(\alpha)}{K}$ separabel.
\bewende
% subsection 1614 (end)

\subsection{Satz vom primitiven Element} % (fold)
\label{sub:1615}
Sei $\nicefrac{L}{K}$ eine endliche und separable Körpererweiterung. Dann gibt es $\alpha \in L$ mit $L=K(\alpha)$.
\minisec{Beweis}
Da $\nicefrac{L}{K}$ endlich ist, gibt es $\alpha_1, \ldots , \alpha_n \in L$ mit $L=K(\alpha_1, \ldots , \alpha_n)$. Die Behauptung folgt per Induktion aus folgendem
Lemma \ref{sub:1616} \bewende
% subsection 1615 (end)

\subsection[Lemma: Induktionsschritt für den Satz vom primitiven Element]{Lemma} % (fold)
\label{sub:1616}
Sei $\nicefrac{L}{K}$ eine separable Körpererweiterung. Seien $\alpha, \beta \in L$. Dann gibt es $\gamma \in L$ mit $K(\alpha, \beta) = K(\gamma)$.
\minisec{Beweis (für den Fall $\abs*{K}=\infty$)}
Mit $\nicefrac{L}{K}$ ist auch $\nicefrac{K(\alpha, \beta)}{K}$ separabel. Sei $\nicefrac{N}{K(\alpha, \beta)}$ eine Körpererweiterung, so dass $\nicefrac{N}{K}$ normal
ist. Dann gibt es $n=[K(\alpha,\beta):K]$-viele $K$-Homomorphismen $\sigma_1, \ldots , \sigma_n : K(\alpha,\beta) \to \nicefrac{N}{K}$. Da die $\sigma_i$ verschieden sind,
gilt für $i \not= j: $ $\sigma_i(\alpha) \not= \sigma_j(\alpha)$ oder $\sigma_i(\beta) \not= \sigma_j(\beta)$. Für $i \not= j$ gilt also
\[
	0 \not= \big( \sigma_i(\alpha)- \sigma_j(\alpha)\big) + \big(\sigma_i(\beta)- \sigma_j(\beta)\big) X = (\sigma_i(\alpha) + \sigma_i(\beta) X) - (\sigma_j(\alpha)+ \sigma_j(\beta)X)
	\in N[X]
\]
Betrachte 
\[
	f= \prod_{i\not= j} \big((\sigma_i(\alpha)+ \sigma_i(\beta)X) - (\sigma_j(\alpha)+ \sigma_j(\beta)X)\big) \not= 0 
\]
Da $\abs*{K}=\infty $ und $f\not= 0$, gibt es $x \in K$ mit $f(x) \not= 0$. Es folgt, dass die Elemente $\sigma_i(\alpha)+ \sigma_i(\beta)x = \sigma_i(\alpha+ \beta x)$
paarweise verschieden sind. Sei $\gamma := \alpha+ \beta x \in L$. Da die $\sigma_i(\gamma)$ alle Nullstellen von $p_\gamma$ sind, folgt 
$[K(\gamma):K] \ge n = [K(\alpha, \beta):K]$. Da $\gamma \in K(\alpha, \beta)$ folgt $K(\gamma) \subseteq K(\alpha,\beta)$ und damit insgesamt 
$K(\gamma)= K(\alpha, \beta)$. \bewende
% subsection 1616 (end)
% section 16 (end)
\newpage

\section{Galois-Theorie} % (fold)
\label{sec:17}

\subsection[Definition: Galois-Erweiterung und Galois-Gruppe]{Definition} % (fold)
\label{sub:171}
Eine algebraische Körpererweiterung $\nicefrac{L}{K}$ heißt eine \Index{Galois-Erweiterung} (oder galoissch), wenn sie normal und separabel ist. Dann heißt 
$\Gal \enbrace*{\nicefrac{L}{K}} := \Aut \enbrace*{\nicefrac{L}{K}}  $ die \Index{Galois-Gruppe} von $\nicefrac{L}{K}$. 
% subsection 171 (end)

\subsection[Bemerkung: separabel, normal und galoissch im Bezug auf Zwischenkörper]{Bemerkung} % (fold)
\label{sub:172}
Sei $E$ ein Zwischenkörper von $\nicefrac{L}{K}$. 
\begin{enumerate}[(i)]
	\item Ist $\nicefrac{L}{K}$ separabel, so ist auch $\nicefrac{E}{K}$ und $\nicefrac{L}{E}$ separabel. 
	
	Beweis für Letzteres: Sei $\alpha \in L$ und $p_\alpha^E \in E[X]$
	das Minimalpolynom von $\alpha$ über $E$ und $p_\alpha^K$ das Minimalpolynom von $\alpha$ über $K$. Dann gilt $p_\alpha^E \mid p_\alpha^K$ in $E[X]$. Da $p_\alpha^K$
	separabel ist, ist auch $p_\alpha^E$ separabel.
	\item Ist $\nicefrac{L}{K}$ normal, so ist auch $\nicefrac{L}{E}$ normal. Beweis: Gleiches Argument wie eben.
	\item Ist $\nicefrac{L}{K}$ galoissch, so ist auch $\nicefrac{L}{E}$ galoisch.
\end{enumerate}
% subsection 172 (end)

\subsection[{Proposition 1: $\abs*{\Gal (\nicefrac{L}{K})}= [L :K]$, falls $[L:K] <\infty$ und galoissch}]{Proposition 1} % (fold)
\label{sub:173}
Sei $\nicefrac{L}{K}$ eine endliche Galoiserweiterung. Dann gilt $\abs*{\Gal (\nicefrac{L}{K})}= [L :K]$.
\minisec{Beweis}
Nach dem Satz vom primitiven Element (\ref{sub:1615}) gibt es $\alpha \in L$ mit $L=K(\alpha)$. Da $\nicefrac{L}{K}$ normal ist, zerfällt $p_\alpha$ über $L$ in 
Linearfaktoren.  Da $\nicefrac{L}{K}$ separabel ist, besitzt $p_\alpha$ $n=[L:K]= \grad p_\alpha$-viele Nullstellen in $L$. Nach dem Fortsetzungssatz gibt es genau
$n$-viele $K$-Homomorphismen $\sigma_i :\nicefrac{L}{K} \to \nicefrac{L}{K}$ (Diese wenden durch $\sigma_i(\alpha)$ bestimmt; die $\sigma_i(\alpha)$ sind genau die
Nullstellen von $p_\alpha$).

Da $\nicefrac{L}{K}$ algebraisch (also sogar endlich) ist, sind die $\sigma_i$ Automorphismen. Es folgt $\Gal(\nicefrac{L}{K})= \set{\sigma_1, \ldots, \sigma_n} $. Insbesondere
\[
	\abs*{\Gal( \nicefrac{L}{K})} = n = [L :K] \bewende 
\]
% subsection 173 (end)

\subsection[Definition: Fixkörper]{Definition} % (fold)
\label{sub:174}
Sei $G$ eine Untergruppe von $\Aut(\nicefrac{L}{K})$. Dann heißt $L^G := \set[\alpha \in L]{\sigma(\alpha)=\alpha \text{ für alle } \sigma \in G} $ der \Index{Fixkörper}
von $G$. 
% subsection 174 (end)

\subsection[Bemerkung: Begründung warum der Fixkörper ein Körper ist]{Bemerkung} % (fold)
\label{sub:175}
Für $\alpha, \beta \in L^G$ und $\sigma \in G$ gilt: 
\begin{itemize}
	\item $\sigma(\alpha + \beta)= \sigma(\alpha)+ \sigma(\beta)= \alpha + \beta$, also $\alpha+ \beta \in L^G$.
	\item $\sigma(\alpha \cdot \beta) = \sigma(\alpha) \cdot \sigma(\beta) = \alpha \cdot  \beta$, also $\alpha \cdot \beta \in L^G$.
	\item $\sigma(-\alpha)= - \sigma(\alpha) = - \alpha$, also $- \alpha \in L^G$.
	\item Falls $\alpha \not= 0$: $ \sigma(\alpha ^{-1}) = \sigma(\alpha) ^{-1} = \alpha ^{-1}$, also $\alpha ^{-1} \in L^G$.
	\item Für alle $\lambda  \in K$ ist $\sigma(\lambda )=\lambda$, also $K \subseteq L^G$.
\end{itemize}
$L^G$ ist also ein Zwischenkörper von $\nicefrac{L}{K}$.
% subsection 175 (end)

\subsection[Proposition 2: Die Galois-Gruppe zu $L^G$ ist gleich der Untergruppe $G \le \Aut(\nicefrac{L}{K})$]{Proposition 2} % (fold)
\label{sub:176}
Sei $\nicefrac{L}{K}$ eine Körpererweiterung. Sei $G$ eine endliche Untergruppe von $\Aut(\nicefrac{L}{K})$. Sei $E:= L^G$. Dann ist $\nicefrac{L}{E}$ eine
endliche Galois-Erweiterung mit $\Gal(\nicefrac{L}{E}) = G$. Insbesondere ist $[L:E]= \abs*{G} $.
\minisec{Beweis}
Zu $\alpha \in L$ sei $G_\alpha = \set[\sigma \in G]{\sigma(\alpha)=\alpha}$ die Standgruppe von $\alpha$ und 
\[
	f_\alpha := \prod_{\sigma G_\alpha \in \nicefrac{G}{G_\alpha}} \big(X-\sigma(\alpha)\big)
\]
Für $\tau \in G$ gilt 
\[
	f^\tau_\alpha = \prod_{\sigma G_\alpha \in \nicefrac{G}{G_\alpha}} \big(X-\tau\sigma(\alpha)\big) = \prod_{\sigma G_\alpha \in \nicefrac{G}{G_\alpha}} \big(X-\sigma(\alpha)\big)= f_\alpha
\]
Es folgt $f_\alpha \in E[X]$. Für $\sigma, \tau \in G$ gilt 
\begin{align*}
	\sigma(\alpha)=\tau(\alpha) \iff \tau ^{-1} \sigma(\alpha)= \alpha \iff \tau ^{-1} \sigma \in G_\alpha 
	\iff \tau G_\alpha = \sigma G_\alpha
\end{align*}
Daher sind die Nullstellen von $f_\alpha$ paarweise verschieden. Sei $p_\alpha \in E[X]$ das Minimalpolynom von $\alpha$. Da $f_\alpha(\alpha)=0$ gilt 
$p_\alpha \mid f_\alpha$. Damit hat $p_\alpha$ keine doppelten Nullstellen und zerfällt über $L$ in Linearfaktoren. Daher ist $\nicefrac{L}{E}$ separabel und normal.
Wegen $f_\alpha(\alpha)=0$ ist $\nicefrac{L}{E}$ auch algebraisch. Insgesamt ist $\nicefrac{L}{E}$ eine Galois-Erweiterung.

Wir zeigen als nächstes $[L:L^G] \le \abs*{G} $. Angenommen: $[L:L^G] \gneq \abs*{G} $. Dann gibt es einen Zwischenkörper $E$ von $\nicefrac{L}{L^G}$ mit $\nicefrac{E}{L^G}$
endlich und $[E:L^G] \gneq \abs*{G} $. Da mit $\nicefrac{L}{L^G}$ auch $\nicefrac{E}{L^G}$ separabel ist, gibt es nach dem Satz vom primitiven Element (\ref{sub:1615}) 
$\alpha \in E$ mit $E=L^G(\alpha)$. Es ist 
\[
	[L^G(\alpha): L^G]= \grad p_\alpha \stackrel{f_\alpha(\alpha)=0}{\le} \grad f_\alpha \le \abs*{G} \lightning 
\]
Es gilt also $[L:L^G] \le \abs*{G} $. Nach Proposition 1 (\ref{sub:173}) ist $[L:L^G]= \abs*{\Gal(\nicefrac{L}{L^G})}$. Es ist aber $G$ eine Untergruppe von $\Gal(\nicefrac{L}{L^G})$. Es folgt $G=\Gal(\nicefrac{L}{L^G})$ (da $\abs*{\Gal(\nicefrac{L}{L^G})} \le \abs*{G}  $). \bewende
% subsection 176 (end)

\subsection[Bemerkung zur Übersetzung zwischen Zwischenkörpern und Untergruppen von $\Aut(\nicefrac{L}{K})$]{Bemerkung} % (fold)
\label{sub:177}
Sei $\nicefrac{L}{K}$ eine Körpererweiterung. Seien 
\begin{align*}
	\mathcal{Z}(\nicefrac{L}{K}) &= \text{ die Menge aller Zwischenkörper von } \nicefrac{L}{K} \\
	\mathcal{U}(\nicefrac{L}{K}) &= \text{ die Menge aller Untergruppen von } \Aut(\nicefrac{L}{K}) \\
	\intertext{Wir haben Abbildungen:}
	\mathrm{ug} &: \mathcal{Z}(\nicefrac{L}{K}) \to \mathcal{U}(\nicefrac{L}{K}) \quad \mathrm{ug}(E) := \Aut(\nicefrac{L}{E}) \\
	\mathrm{zw} &: \mathcal{U}(\nicefrac{L}{K}) \to \mathcal{Z}(\nicefrac{L}{K}) \quad \mathrm{zw}(G) := L^G \\
	\intertext{Diese erfüllen}
	&H \le G \Longrightarrow \mathrm{zw}(H) \supseteq \mathrm{zw}(G) \\
	&E \subseteq F \Longrightarrow \mathrm{ug}\,(E) \ge \mathrm{ug}(F)
\end{align*}
% subsection 177 (end)

\subsection{Hauptsatz der Galois-Theorie (für endliche Galois-Erweiterungen)} % (fold)
\label{sub:178}
Sei $\nicefrac{L}{K}$ eine endliche Galois-Erweiterung. Dann gelten: 
\begin{enumerate}[(1)]
	\item $\mathrm{ug}$ und $\mathrm{zw}$ sind zueinander inverse Abbildungen zwischen der Menge der Untergruppen von $\Gal(\nicefrac{L}{K})$ der Menge der Zwischenkörper
	von $\nicefrac{L}{K}$.
	\item Für $E \in \mathcal{Z}(\nicefrac{L}{K})$ gilt $[L:E]= \abs*{\Gal(\nicefrac{L}{E})}$.
	\item Sei $E$ ein Zwischenkörper von $\nicefrac{L}{K}$. Dann ist $\nicefrac{E}{K}$ genau dann normal (und damit galoissch), falls 
	$H:=\mathrm{ug}(E)= \Gal(\nicefrac{L}{E})$ ein Normalteiler von $\Gal(\nicefrac{L}{K})$ ist. In diesem Fall ist $\Gal(\nicefrac{L}{E})$ der Kern des surjektiven
	Gruppenhomomorphismus $\Gal(\nicefrac{L}{K}) \to \Gal(\nicefrac{E}{K}), \sigma \mapsto \sigma\big|_E$. Insbesondere ist 
	\[
		\Gal(\nicefrac{E}{K}) \cong \nicefrac{\Gal(\nicefrac{L}{K})}{\Gal(\nicefrac{L}{E})}.
	\]
\end{enumerate}
\minisec{Beweis}
\begin{enumerate}[(1)]
	\item Sei $H \le \Gal(\nicefrac{L}{K})$. Dann ist nach Proposition 2 (\ref{sub:176}) $H= \Gal(\nicefrac{L}{L^H})$. Also $\mathrm{ug} \circ \mathrm{zw}(H)=H$. Sei
	$E \in \mathcal{Z}(\nicefrac{L}{K})$. Dann ist sicher $E  \subseteq L^{\Gal(\nicefrac{L}{E})}$. Nach Proposition 1\& 2 gilt 
	\[
		[L:L^{\Gal(\nicefrac{L}{E})}] \stackrel{\text{Prop. 1 \& 2}}{=} \abs*{\Gal(\nicefrac{L}{E})} \stackrel{\text{Prop. 1}}{=} [L :E]
	\]
	Es folgt $E= L^{\Gal(\nicefrac{L}{E})}$. Also ist $\mathrm{zw} \circ \mathrm{ug}(E)=E$.
	\item Haben wir schon in Proposition 1 (\ref{sub:173}) gesehen.
	\item Sei $E \in \mathcal{Z}(\nicefrac{L}{K})$ und $\nicefrac{E}{K}$ normal. In Kapitel \ref{sec:14} haben wir gesehen, dass sich jedes $\sigma \in \Gal(\nicefrac{L}{K})$
	zu einem Element $\sigma\big|_E \in \Gal(\nicefrac{E}{K})$ einschränkt (Da $\sigma(E) \subseteq E$). Wir erhalten einen Gruppenhomomorphismus 
	\[
		\Gal(\nicefrac{L}{K}) \xrightarrow{\varphi} \Gal(\nicefrac{E}{K}) \text{ mit } \varphi(\sigma) := \sigma\big|_E 
	\]
	und es ist $\Kern \varphi = \Gal(\nicefrac{L}{E})$. Insbesondere ist $\Gal(\nicefrac{L}{E})$ normal in $\Gal(\nicefrac{L}{K})$. Weiter gilt
	\[
		\abs*{\im \varphi} = \frac{\abs*{\Gal(\nicefrac{L}{K})} }{\abs*{\Gal(\nicefrac{L}{E})} }  = \frac{[L:K]}{[L:E]} = [E:K] = \abs*{\Gal(\nicefrac{E}{K})} 
	\]
	Also ist $\varphi$ surjektiv. 
	
	Sei umgekehrt $H \unlhd \Gal(\nicefrac{L}{K})$ eine normale Untergruppe. Sei $E=L^H$. Sei $\alpha \in E$ und $\sigma \in \Gal(\nicefrac{L}{K})$. Für 
	$\tau\in H= \Gal(\nicefrac{L}{E})$ gilt
	\[
		\tau\big(\sigma(\alpha)\big) = \sigma \underbrace{\sigma ^{-1} \tau \sigma}_{\in H, \text{ da } H \text{ normal}} (\alpha) \stackrel{\alpha \in E=L^H}{=} \sigma(\alpha)
	\]
	Es folgt $\sigma(\alpha) \in E=L^H$. Also $\sigma(E) =E$. Sei $G:= \set[\sigma\big|_E]{\sigma \in \Gal(\nicefrac{L}{K})} \le \Aut(\nicefrac{E}{K})$. Dann ist 
	\[
		K \subseteq E^G \subseteq L^{\Gal(\nicefrac{L}{K})} = K
	\]
	Also $E^G=K$. Nach Prop. 2 (\ref{sub:176}) ist $\nicefrac{E}{K} = \nicefrac{E}{E^G}$ dann galoissch und insbesondere normal. \bewende
\end{enumerate}
% subsection 178 (end)

\subsection[Definition: Galois-Gruppe eines Polynoms]{Definition} % (fold)
\label{sub:179}
Sei $f \in K[X]$ separabel. Sei $L$ der Zerfällungskörper von $f$ (Dann ist $\nicefrac{L}{K}$ galoissch.). Dann heißt $\Gal_K(f) := \Gal(\nicefrac{L}{K})$ die
\bet{Galois-Gruppe}\index{Galois-Gruppe!eines Polynoms} von $f$ über $K$\index{Galois-Gruppe!eines Polynoms}.
% subsection 179 (end)

\subsection[Beispiel einer Galois-Gruppe eines Polynoms]{Beispiel} % (fold)
\label{sub:1710}
Sei $f=X^3-2 \in \mathds{Q}[X]$. Der Zerfällungskörper von $f$ ist  $\mathds{Q}(\sqrt[3]{2}, \zeta \sqrt[3]{2}, \zeta^2 \sqrt[3]{2}) =\mathds{Q}(\sqrt[3]{2}, \zeta  )$ 
mit $\zeta= e^{2 \pi  i /3} $. \\
\uline{Behauptung:} $\Gal_\mathds{Q}(f) \cong S_3$.
\minisec{Beweis}
Sei $Z := \set{\sqrt[3]{2}, \zeta \sqrt[3]{2}, \zeta^2 \sqrt[3]{2} } $ die Menge der Nullstellen von $f$. Für $\sigma \in \Gal_\mathds{Q}(f)$ gilt $\sigma(Z)=Z$. Wir 
erhalten also eine Wirkung der Gruppe $\Gal_\mathds{Q}(f)$ auf $Z$. In anderen Worten erhalten wir einen Gruppenhomomorphismus 
$\varphi : \Gal_\mathds{Q}(f) \to S_Z \cong S_3$. Da $\mathds{Q}(\sqrt[3]{2}, \zeta  ) = \mathds{Q}(Z)$ ist $\varphi$ injektiv. Da 
$\abs*{\Gal_\mathds{Q}(f)} =  [\mathds{Q}(\sqrt[3]{2}, \zeta  ): \mathds{Q}]=6 = \abs*{S_3} $ folgt $\Gal_Q(f)  \cong S_3$. \bewende
% subsection 1710 (end)

\subsection[Bemerkung: Die Galois-Gruppe wirkt auf der Menge der Nullstellen]{Bemerkung} % (fold)
\label{sub:1711}
Ist $f \in K[X]$ separabel und $n= \deg f$, so hat $f$ im Zerfällungskörper $L$ von $f$ $n$ verschiedene Nullstellen. Wir erhalten eine Wirkung von $\Gal_K(f)$ auf der 
Menge der Nullstellen und einen injektiven Gruppenhomomorphismus $\Gal_K(f) \to S_n$ (Injektivität: $L=K(M)$). Insbesondere 
\[
	[L:K] = \abs*{\Gal_K(f)} \le \abs*{S_n} = n!  
\]
% subsection 1711 (end)

\subsection[Beispiel zur Ungleichung in \ref{sub:1711}]{Beispiel} % (fold)
\label{sub:1712}
Der Zerfällungskörper von $X^4-2 \in \mathds{Q}[X]$ ist $\mathds{Q}(\sqrt[4]{2},i)$. Es ist 
\[
	\left[\mathds{Q}(\sqrt[4]{2},i) :\mathds{Q}\right] = \left[\mathds{Q}(\sqrt[4]{2},i) : \mathds{Q}(\sqrt[4]{2})\right] \cdot \left[\mathds{Q}(\sqrt[4]{2}) : \mathds{Q}\right] = 2 \cdot 4 = 8
\]
Insbesondere $\abs*{\Gal_\mathds{Q}(X^4-1)} = 8 $, aber $\abs*{S_4} = 24 $.
% subsection 1712 (end)
% section 17 (end)
\newpage

\section{Der Fundamentalsatz der Algebra} % (fold)
\label{sec:18}

\subsection[Definition: Algebraisch abgeschlossen]{Definition} % (fold)
\label{sub:181}
Ein Körper $K$ heißt \Index{algebraisch abgeschlossen}, wenn jedes Polynom $f \in K[X]$ von Grad $\ge 1$ eine Nullstelle in $K$ besitzt.
% subsection 181 (end)

\subsection[Bemerkung: In algebraisch abgeschlossenen Körper zerfällt jedes Polynom]{Bemerkung} % (fold)
\label{sub:182}
Ist $K$ algebraisch abgeschlossen, so zerfällt jedes Polynom über $K$ in Linearfaktoren.
% subsection 182 (end)

\subsection[{Bemerkung: Algebraisch abgeschlossen $\iff$ $\not\exists L$ mit $[L:K]=2$}]{Bemerkung} % (fold)
\label{sub:183}
$K$ ist genau dann algebraisch abgeschlossen, wenn es keine endliche Körpererweiterung $\nicefrac{L}{K}$ gibt mit $[L:K]\ge 2$.\hfill (Übung)
% subsection 183 (end)

\subsection{Fundamentalsatz der Algebra} % (fold)
\label{sub:184}
Der Körper $\mathds{C}$ ist algebraisch abgeschlossen.
% subsection 184 (end)

\subsection[Bemerkung zu benutzten Erkenntnissen aus der Analysis]{Bemerkung} % (fold)
\label{sub:185}
Wir wissen schon (teilweise mit Mitteln der Analysis):
\begin{enumerate}[a)]
	\item Jedes Polynom $f \in \mathds{C}[X]$ vom Grad 2 hat eine Nullstelle in $\mathds{C}$.\hfill (quadratische Erweiterung)
	\item Jedes Polynom $f \in \mathds{R}[X]$ von ungeradem Grad besitzt eine Nullstelle in $\mathds{R}$.\hfill (ZWS)
\end{enumerate}
% subsection 185 (end)

\subsection[Lemma: Grundlegende Eigenschaften des Grades von KE  $\nicefrac{L}{\mathds{C}}$ und $\nicefrac{L}{\mathds{R}}$]{Lemma} % (fold)
\label{sub:186}
\begin{enumerate}[a)]
	\item Es gibt keine Körpererweiterung $\nicefrac{L}{\mathds{C}}$ mit $[L : \mathds{C}] =2$.
	\item Sei $\nicefrac{L}{\mathds{R}}$ eine Körpererweiterung $\nicefrac{L}{\mathds{R}}$ von ungeradem Grad $[L:\mathds{R}]\ge 3$. Dann gilt $L=\mathds{R}$.
\end{enumerate}
\minisec{Beweis}
\begin{enumerate}[a)]
	\item Angenommen doch. Sei $\alpha \in L \setminus \mathds{C}$. Dann ist $p_\alpha \in \mathds{C}[X]$ ein irreduzibles Polynom vom Grad$=2$. \light
	\item Sei $\alpha \in L$. Dann ist auch $\nicefrac{\mathds{R}(\alpha)}{\mathds{R}}$ eine Körpererweiterung von ungeradem Grad. Damit ist $p_\alpha \in \mathds{R}[X]$ ein
	irreduzibles Polynom von ungeradem Grad. Es folgt $\deg p_\alpha =1$ und $\alpha \in \mathds{R}$. \bewende
\end{enumerate}
% subsection 186 (end)

\subsection[Bemerkung zur Beweisstrategie]{Bemerkung} % (fold)
\label{sub:187}
Ist $\nicefrac{L}{\mathds{C}}$ eine endliche Körpererweiterung. Da $\Char \mathds{C} =0$ ist $\nicefrac{L}{\mathds{C}}$ separabel. Es gibt  eine endliche Körpererweiterung $\nicefrac{N}{L}$, so dass $\nicefrac{N}{\mathds{C}}$ normal und damit galoissch ist. Um den Fundamentalsatz zu 
beweisen, genügt es daher zu zeigen, dass es keine (endlichen) Galoiserweiterungen von $\mathds{C}$ gibt.
% subsection 187 (end)

\subsection[Lemma 2: Endliche KE von $\mathds{R}$ haben Grad $2^k$]{Lemma 2} % (fold)
\label{sub:188}
Sei $\nicefrac{L}{\mathds{R}}$ eine endliche Galoiserweiterung. Dann gilt $[L:\mathds{R}]= 2^k$ mit $k \in \mathds{N}$.
\minisec{Beweis}
Sei $H$ eine $2$-Sylowuntergruppe von $G := \Gal(\nicefrac{L}{\mathds{R}})$. Dann gilt $\abs*{H}= 2^k $ und $[G:H]$ ist ungerade. Sei $E := L^H$. Mit dem Hauptsatz der 
Galois-Theorie (\ref{sub:178}) folgt: $\Gal(\nicefrac{L}{E}) =H$ und $[L:E]= \abs*{H} = 2^k $. Ebenfalls mit dem Hauptsatz der Galois-Theorie folgt
\[
	[L:\mathds{R}] = \abs*{G} = \abs*{H} \cdot [G:H].  
\]
Es folgt, dass $[E:\mathds{R}] = [G :H]$ ungerade ist. Damit folgt $E=\mathds{R}$. Also $[L:\mathds{R}] = [L:E] = 2^k$. \bewende
% subsection 188 (end)

\subsection[Proposition: Ketten von Normalteilern in einer $p$-Gruppe]{Proposition} % (fold)
\label{sub:189}
Sei $p$ eine Primzahl und $P$ eine $p$-Gruppe. Dann gibt es eine Folge $\{e\} = N_0 \unlhd N_1 \unlhd N_2 \unlhd \ldots \unlhd N_k =P$ von Normalteilern mit
$\abs*{N_i} =p^i $.
\minisec{Beweis}
Sei $\abs*{P}= p^k$. Induktion nach $k$. Für $k=0$ ist nicht zu zeigen. \\
\uline{Induktionsschritt: $k -1 \mapsto k$} Sei $\abs*{P}=p^k $. \\
Da $P$ eine $p$-Gruppe ist, ist das Zentrum von $P$ nichttrivial (Kapitel \ref{sec:3}). Sei $z \in P$ ein nicht-triviales zentrales Element. Da $\ord(z)$ die Ordnung von 
$P$ teilt, gilt $\ord (z) = p^l$ für ein $l$. Dann ist $g:= z^{p^{l-1}}$ ein nicht-triviales zentrales Element der Ordnung $p$. Dann ist $N_1 := \langle g \rangle$ ein
Normalteiler mit $\abs*{N_1} =p $.

Nun wenden wir die Induktionsannahme auf $Q := \nicefrac{P}{N_1}$ an (Es ist $\abs*{Q} = \frac{\abs*{P} }{\abs*{N_1} } = p^{k-1}$). Wir erhalten Normalteiler: 
\[
	\{e\} = M_0 \unlhd M_1 \unlhd \ldots \unlhd M_{k-1} =Q 
\]
mit $\abs*{M_i}=p^i$. Sei $\pi  : P \to Q$ die Quotientenabbildung. Dann sind die $N_i := \pi ^{-1}(M_{i-1})$ für $i=1, \ldots ,k$ die gesuchten Normalteiler. \bewende
% subsection 189 (end)

\subsection[Korollar zu Galoiserweiterungen mit Grad $p^k$]{Korollar} % (fold)
\label{sub:1810}
Sei $p$ eine Primzahl und $\nicefrac{L}{K}$ eine Galoiserweiterung von Grad $p^k$. Dann gibt es eine Folge von Zwischenkörpern 
\[
	K=E_0 \subseteq E_1 \subseteq \ldots \subseteq E_k =L
\]
mit $[E_i :K] = p^i$ und $[E_i : E_{i-1}]=p$.
\minisec{Beweis}
Für $P := \Gal(\nicefrac{L}{K})$ gilt nach dem Hauptsatz $\abs*{P} = [L :K] =p^k $. Mit Proposition \ref{sub:189} folgt:
\[
	\exists \{e\} = N_0 \unlhd N_1 \unlhd \ldots \unlhd N_k =P \quad \text{mit } \abs*{N_i}=p^i  
\]
Setze nun $E_i := L^{N_{k-i}}$. Dann $K=E_0  \subseteq E_1 \subseteq \ldots \subseteq E_k=L$ und nach dem Hauptsatz folgt
\[
	[E_i :K] = \frac{[L:K]}{[L:E_i]} = \frac{p^k}{p^{k-i}} = p^{i}. \bewende  
\]
% subsection 1810 (end)

\subsection{Beweis des Fundamentalsatzes} % (fold)
\label{sub:1811}
Sei $\nicefrac{L}{\mathds{C}}$ eine endliche Körpererweiterung. Zu zeigen ist $L=\mathds{C}$. Es ist auch $\nicefrac{L}{\mathds{R}}$ eine endliche Körpererweiterung. Sei
$\nicefrac{N}{L}$ eine Körpererweiterung für die $\nicefrac{N}{\mathds{R}}$ normal und damit galoissch ist. Wegen Lemma 2 (\ref{sub:188}) ist der Grad von 
$\nicefrac{N}{\mathds{R}}$ und damit der auch der Grad von $\nicefrac{N}{\mathds{C}}$ eine 2er-Potenz. Da $\nicefrac{N}{\mathds{R}}$ galoissch ist, ist auch 
$\nicefrac{N}{\mathds{C}}$ galoissch und die Galoisgruppe von $\nicefrac{N}{\mathds{C}}$ eine $2$-Gruppe, $\Gal(\nicefrac{N}{\mathds{C}})= 2^k$. Ist $k=0$ so folgt mit dem
Hauptsatz der Galoistheorie $N=\mathds{C}$ und damit $L=\mathds{C}$.

Angenommen $k \ge 1$. Nach dem Korollar gibt es dann einen Zwischenkörper $E$ von $\nicefrac{N}{\mathds{C}}$ mit $[E:\mathds{C}] =2$ im Widerspruch zu 
Lemma 1 (\ref{sub:186}) a). \bewende
% subsection 1811 (end)
% section 18 (end)
\newpage

\section{Einheitswurzeln} % (fold)
\label{sec:19}

\subsection[Definition: Einheitswurzel]{Definition} % (fold)
\label{sub:191}
Sei $K$ ein Körper. $\zeta \in K$ heißt eine \bet{$n$-te Einheitswurzel}\index{n-te Einheitswurzel@$n$-te Einheitswurzel}, falls $\zeta^n =1$. 
Die $n$-ten Einheitswurzeln bilden
eine Gruppe $E_n(K)$ bezüglich Multiplikation in $K$. $\zeta$ heißt eine 
\bet{primitive $n$-te Einheitswurzel}\index{n-te Einheitswurzel@$n$-te Einheitswurzel!primitive} falls
$\zeta^n =1$ aber $\zeta^k\not= 1$ für $0 < k \le n-1$. Die Menge der primitiven Einheitswurzeln bezeichnen wir mit $PE_n(K)$.
% subsection 191 (end)

\subsection[Bemerkung: Primitive Einheitswurzeln erzeugen die Gruppe der Einheitswurzeln]{Bemerkung} % (fold)
\label{sub:192}
$\zeta \in E_n(K)$ ist genau dann eine primitive Einheitswurzel, wenn $\zeta$ ein Erzeuger von $E_n(K)$ ist, also $\langle \zeta \rangle = E_n(K)$.
% subsection 192 (end)

\subsection[Beispiel: Einheitswurzeln in $\mathds{C}$]{Beispiel} % (fold)
\label{sub:193}
\begin{align*}
	E_n(\mathds{C}) &= \set[e^{2 \pi i k /n}]{0 \le k \le n-1} \cong \nicefrac{\mathds{Z}}{n \mathds{Z}} \tag{$e^{\frac{2 \pi i k}{n} } \mapsto k+ n \mathds{Z}$} \\
	PE_n(\mathds{C}) &= \set[e^{2 \pi i k /n}]{1 \le k \le n-1 , \ggT(k,n)=1}  
\end{align*}
% subsection 193 (end)

\subsection[Bemerkung: Vergleich der Einheiten in $\nicefrac{\mathds{Z}}{n \mathds{Z}}$ mit $PE_n(\mathds{C})$]{Bemerkung} % (fold)
\label{sub:194}
Für die Einheitengruppe des Rings $\nicefrac{\mathds{Z}}{n \mathds{Z}}$ gilt
\[
	\enbrace*{\nicefrac{\mathds{Z}}{n \mathds{Z}}}^\times = \set[k+ n \mathds{Z}]{1 \le k \le n-1, \ggT(k,n)=1}.  
\]
Insbesondere ist $\abs*{PE_n(\mathds{C})} = \abs*{(\nicefrac{\mathds{Z}}{n \mathds{Z}})^\times}$.
% subsection 194 (end)

\subsection[Definition: Eulersche $\varphi$-Funktion]{Definition} % (fold)
\label{sub:195}
Für $n \in \mathds{N}_{>0}$ setzen wir $\varphi(n) := \abs*{ (\nicefrac{\mathds{Z}}{n \mathds{Z}})^\times} $. Die so erklärte Funktion 
$\varphi : \mathds{N}_{>0} \to \mathds{N}_{>0}$ heißt die \Index{Eulersche $\varphi$-Funktion}. Es ist $\varphi(n) = \abs*{PE_n(\mathds{C})} $.
% subsection 195 (end)

\subsection[Lemma über die Funktionswerte von $\varphi$]{Lemma} % (fold)
\label{sub:196}
\begin{enumerate}[(i)]
	\item Ist $p$ eine Primzahl so gilt $\varphi(p)= p-1$ und $\varphi(p^\nu) =(p-1) \cdot p^{\nu -1}$.
	\item Ist $\ggT(n,m)=1$ so gilt $\varphi(n \cdot m)= \varphi(n) \cdot \varphi(m)$.
	\item Ist $n= p_1^{\nu_1} \cdot \ldots \cdot p_r^{\nu_r}$ die Primfaktorzerlegung von $n$ mit paarweise verschiedenen Primzahlen $p_i$ so gilt 
	\[
		\varphi(n) = \prod_{i=1}^r (p_i -1) \cdot p_i^{\nu_i-1}
	\]
\end{enumerate}
\minisec{Beweis}
\begin{enumerate}[(i)]
	\item Es ist $\ggT(p^\nu,k)>1$ genau dann wenn $p \mid k$. Es folgt 
	\begin{align*}
		\varphi(p^\nu) = \abs*{ \set[k]{1 \le k \le p^\nu -1, \rule{0cm}{0.4cm} p \nmid k} } &= (p^\nu -1) - \abs*{ \set[k]{1 \le k \le p^\nu -1 , \rule{0cm}{0.4cm}p \mid k} }  \\
		&= (p^\nu -1) - \abs*{\set{p, 2 \cdot p, \ldots , (p^{\nu -1} -1)\cdot p} } \\
		&= (p^\nu -1) -(p^{\nu -1} -1) = p^\nu - p^{\nu-1} = (p-1) p^{\nu -1}.
	\end{align*}
	\item Nach dem Chinesischen Restsatz (\ref{sub:618}) gibt es einen Ringisomorphismus 
	$\nicefrac{\mathds{Z}}{n \cdot m \mathds{Z}} \xrightarrow{\cong}  \nicefrac{\mathds{Z}}{n \mathds{Z}} \times \nicefrac{\mathds{Z}}{m \mathds{Z}}$,
	$k+ nm \mathds{Z} \mapsto (k+ n \mathds{Z}, k+ m \mathds{Z})$, da $n$ und $m$ teilerfremd
	sind. Wegen
	\[
		\enbrace*{\nicefrac{\mathds{Z}}{n \mathds{Z}} \times \nicefrac{\mathds{Z}}{m \mathds{Z}}}^\times = \enbrace*{\nicefrac{\mathds{Z}}{n \mathds{Z}}}^\times \times 
		\enbrace*{\nicefrac{\mathds{Z}}{m \mathds{Z}}  }^\times
	\]
	folgt die Behauptung.
	\item folgt aus (i) und (ii) \bewende
\end{enumerate}
% subsection 196 (end)

\subsection[Korollar: $\abs*{PE_n(\mathds{C})} = \varphi(n)$]{Korollar} % (fold)
\label{sub:197}
\[
	\abs*{PE_n(\mathds{C})} = \varphi(n) 
\]
% subsection 197 (end)

\subsection[Definition: $n$-ter Kreisteilungskörper]{Definition} % (fold)
\label{sub:198}
Sei $\zeta \in PE_n(\mathds{C})$. Dann heißt $\mathds{Q}_n := \mathds{Q}(\zeta)$ der \bet{$n$-te Kreisteilungskörper}\index{n-ter Kreisteilungskörper@$n$-ter Kreisteilungskörper}.
% subsection 198 (end)

\subsection[Bemerkung: Eigenschaften des $n$-ten Kreisteilungskörpers]{Bemerkung} % (fold)
\label{sub:199}
Mit $\zeta$ liegen alle $n$-ten Einheitswurzeln in $\mathds{Q}_n$.
$\mathds{Q}_n$ ist der Zerfällungskörper von $X^n-1 \in \mathds{Q}[X]$. Insbesondere ist $\nicefrac{\mathds{Q}_n}{\mathds{Q}}$ normal. Da algebraische Erweiterungen von 
$\mathds{Q}$ ($\Char \mathds{Q}=0$) immer separabel sind, ist $\nicefrac{\mathds{Q}_n}{\mathds{Q}}$ eine Galois-Erweiterung.
% subsection 199 (end)

\subsection[Lemma 1: Es ex. ein injektiver Gruppenhomomorphismus $\Gal(\nicefrac{\mathds{Q}_n}{\mathds{Q}})\to\enbrace*{\nicefrac{\mathds{Z}}{n\mathds{Z}}}^\times$]{Lemma 1} % (fold)
\label{sub:1910}
Es gibt einen injektiven Gruppenhomomorphismus
\[
	\psi : \Gal( \nicefrac{\mathds{Q}_n}{\mathds{Q}}) \to \enbrace*{\nicefrac{\mathds{Z}}{n \mathds{Z}}}^\times 
\]
\minisec{Beweis}
Sei $\sigma \in \Gal \enbrace*{\nicefrac{\mathds{Q}_n}{\mathds{Q}}} $. Dann permutiert $\sigma$ die Nullstellen von $X^n-1$ und induziert einen Gruppenhomomorphismus 
$\sigma|_{E_n(\mathds{C})} : E_n(\mathds{C}) \to E_n(\mathds{C})$. Insbesondere gilt $\sigma(PE_n)  \subseteq PE_n$. Sei $\zeta := e^{2 \pi  i/n} \in PE_n$. Definiere nun 
$\psi (\sigma) := k+ n \mathds{Z}$ durch
$\sigma(\zeta) = \zeta^k$. Da $\mathds{Q}_n = \mathds{Q}(\zeta)$ gilt
\[
	\sigma = \id \iff  \sigma(\zeta) = \zeta \iff k+ n \mathds{Z}  = 1_{\nicefrac{\mathds{Z}}{n \mathds{Z}}} \tag{$\star$}
\]
Für $\sigma, \tau \in \Gal(\nicefrac{\mathds{Q}_n}{\mathds{Q}})$ mit $\sigma(\zeta) = \zeta^k$ und $\tau (\zeta) = \zeta^l$ gilt
\[
	\sigma \circ \tau(\zeta) = \sigma (\zeta^l) = \sigma(\zeta)^l = \zeta^{k \cdot l}
\]
Daher ist $\psi$ ein Gruppenhomomorphismus und somit injektiv wegen ($\star$). \bewende
% subsection 1910 (end)

\subsection[Definition: $n$-tes Kreisteilungspolynom]{Definition} % (fold)
\label{sub:1911}
\[
	\Phi_n := \prod_{\zeta \in PE_n(\mathds{C})} (X- \zeta)
\]
heißt das \bet{$n$-te Kreisteilungspolynom}\index{n-tes Kreisteilungspolynom@$n$-tes Kreisteilungspolynom}. Es ist $\deg \Phi_n = \varphi(n)$ und die primitiven
Einheitswurzeln sind gerade die Nullstellen von $\Phi_n$.
% subsection 1911 (end)

\subsection[Lemma: Folgerung aus dem Gauß-Lemma]{Lemma} % (fold)
\label{sub:1912}
Seien $f \in \mathds{Z}[X], g \in \mathds{C}[X]$ normiert mit $f \cdot g \in \mathds{Z}[X]$. Dann gilt $g \in \mathds{Z}[X]$.
\minisec{Beweis}
Nach einem Korollar zum Gauß-Lemma (\ref{sub:814}) genügt es zu zeigen $g \in\mathds{Q}[X]$. Sei $f = a_n X^n + a_{n-1} X^{n-1} + \ldots + a_0$ und $g = b_m X^m + b_{m-1} X^{m-1} + \ldots + b_0$.
O.B.d.A. können wir $a_0 \not= 0$ annehmen. (Sonst ersetzen wir $f$ durch $f \cdot X^{-l}$ für geeignetes $l$.)

Mit $f \cdot g \in \mathds{Z}[X]$ folgt $\sum_{i+j=k} a_i \cdot b_j \in \mathds{Z}$ für alle $k$. Insbesondere ist 
\[
	b_k - \frac{1}{a_0} \sum_{i=1}^{k} a_i \cdot b_{k-i} \in \mathds{Q}. 
\]
Per Induktion folgt $b_k \in \mathds{Q}$ für alle $k$. Also $g \in \mathds{Q}[X]$. \bewende
% subsection 1912 (end)

\subsection[{Lemma 2: Das Kreisteilungspolynom liegt in $\mathds{Z}[X]$}]{Lemma 2} % (fold)
\label{sub:1913}
\[
	\Phi_n \in \mathds{Z}[X]
\]
\minisec{Beweis}
Es ist 
\begin{align*}
	X^n-1 = \prod_{\zeta \in E_n} (X- \zeta) = \prod_{d \mid n} \prod_{\substack{\zeta \in E_n \\ \ord (\zeta) =d}} (X- \zeta) = \prod_{d \mid n} \prod_{\zeta \in PE_d} (X- \zeta) 
	= \prod_{d \mid n} \Phi_d.
\end{align*}
Wir zeigen nun $\Phi_n \in \mathds{Z}[X]$ per Induktion nach $n$. Für $n=1$ ist $\Phi_1 = X-1 \in \mathds{Z}[X]$. Für $n>1$ ist
\[
	\Phi_n \cdot \underbrace{\prod_{d \mid n , \, d \not= n} \Phi_d}_{\text{I.A.: } \in \mathds{Z}[X]} = \underbrace{X^n-1 }_{\in \mathds{Z}[X]}
\]
Die Behauptung folgt aus dem vorigen Lemma. \bewende
% subsection 1913 (end)

\subsection[{Satz: $\Phi_n$ ist irreduzibel in $\mathds{Q}[X]$, $[\mathds{Q}_n :\mathds{Q}] = \varphi(n)$}]{Satz} % (fold)
\label{sub:1914}
\begin{enumerate}[a)]
	\item $\Phi_n$ ist irreduzibel in $\mathds{Q}[X]$.
	\item $[\mathds{Q}_n : \mathds{Q}] = \varphi(n)$.
	\item $\Gal(\nicefrac{\mathds{Q}_n}{\mathds{Q}}) \cong (\nicefrac{\mathds{Z}}{n \mathds{Z}})^\times$.
\end{enumerate}
\minisec{Beweis}
Gilt a), so ist $\Phi_n$ das Minimalpolynom für (jedes) $\zeta \in PE_n$ und b) folgt. Wegen Lemma 1 (\ref{sub:1910}) und 
\[
	\abs*{\Gal(\nicefrac{\mathds{Q}_n}{\mathds{Q}})} = [\mathds{Q}_n : \mathds{Q}] \stackrel{\text{b)}}{=} \varphi(n) = \abs*{ (\nicefrac{\mathds{Z}}{n \mathds{Z}})^\times}  
\]
ist c) eine Folgerung aus b). Es bleibt a) zu zeigen. 

Sei $\zeta \in PE_n$. Wegen $\Phi_n(\zeta) =0$ gilt $p_\zeta \mid \Phi_n$. Da $\deg \Phi_n = \abs*{PE_n} = \abs*{\set[\zeta^k]{\ggT(k,n)=1} }  $ genügt es zu zeigen
$p_\zeta (\zeta^k)=0$ für alle $k$ mit $\ggT(k,n)=1$. Wir betrachten zunächst den Fall, dass $k=p$ für eine zu $n$ teilerfremde Primzahl $p$ ist.

Sei $g \in \mathds{Q}[X]$ mit $p_\zeta \cdot g = X^n-1 =: h$. Nach Gauß gilt $g , p_\zeta \in \mathds{Z}[X]$. Angenommen $p_\zeta (\zeta^p) \not= 0$. Dann folgt
$g(\zeta^p)=0$ und $\zeta$ ist eine Nullstelle von $g(X^p)$. Wieder mit Gauß folgt $p_\zeta \mid g(X^p)$ in $\mathds{Z}[X]$. Wir werden nun folgende Beobachtung benutzen:
In $\mathds{F}_p$ gilt 
\[
	(a+b)^p = a^p + b^p,
\]
da alle anderen Binomialkoeffzienten Vielfache von $p$ sind. Daher gilt $f(X^p) = f(X)^p$ für $f \in \mathds{F}_p[X]$. Sei
\[
	\mathds{Z}[X] \to \nicefrac{\mathds{Z}}{p \mathds{Z}}[X] \enspace, \enspace f \mapsto \overline{f} 
\]
der von $\mathds{Z} \to \nicefrac{\mathds{Z}}{p \mathds{Z}}$ induzierte Ringhomomorphismus. Wegen $p_\zeta \mid g(X^p)$ folgt $\overline{p}_\zeta \mid \overline{g}^p  $.
Folglich besitzen $\overline{p}_\zeta $ und $\overline{g}$ eine gemeinsame Nullstelle. Aber wegen $\overline{h}' = n \cdot X^{n-1} $ hat 
$\overline{h} = \overline{p}_\zeta \cdot \overline{g}$ keine doppelte Nullstelle. \light \\
Also $p_\zeta( \zeta^p)=0$.

\bigskip Sei nun $k$ mit $\ggT(n,k)=1$ beliebig. Dann ist $k=p_1 \cdot \ldots  \cdot p_r$ wobei $p_1, \ldots , p_r$ Primzahlen sind, die $n$ nicht teilen. Wegen
$p_\zeta(\zeta^{p_1}) =0$ ist $p_\zeta$  auch das Minimalpolynom von $\zeta^{p_1}$ und es folgt
\[
	p_\zeta ( \zeta^{p_1 \cdot p_2}) = p_{\zeta^{p_1}} \Big( (\zeta^{p_1})^{p_2}\Big) = 0
\]
Induktiv folgt $p_\zeta (\zeta^k) =0$. \bewende
% subsection 1914 (end)
% section 19 (end)
\newpage

\section{$n$-Teilung des Kreises} % (fold)
\label{sec:20}

\subsection{Satz (Charakterisierung von Konstruierbarkeit)} % (fold)
\label{sub:201}
Sei $\alpha \in \mathds{C}$ algebraisch über $\mathds{Q}$. Sei $L$ der Zerfällungskörper von $p_\alpha \in \mathds{Q}[X]$. Dann sind äquivalent
	\begin{enumerate}[1)]
		\item $\alpha \in \zirkel (\{0,1\})$
		\item $[L:\mathds{Q}]$ ist eine Potenz von $2$. 
	\end{enumerate}
\minisec{Beweis}
\begin{description}
	\item[2) $\Rightarrow $ 1):] Als Zerfällungskörper ist $\nicefrac{L}{\mathds{Q}}$ normal. Wegen $\Char \mathds{Q} =0$ ist $\nicefrac{L}{\mathds{Q}}$ auch separabel und
	damit galoissch. In Kapitel \ref{sec:18} haben wir gesehen, dass es dann eine Folge von Zwischenkörpern $\mathds{Q}= K_0 \subseteq  K \subseteq K_1 \subseteq  \ldots \subseteq K_n=L$
	gibt mit $[K_i :\mathds{Q}]=2^i$. Insbesondere $[K_i : K_{i-1}]=2$ und $\nicefrac{K_i}{K_{i-1}}$ eine quadratische Erweiterung. Damit folgt 
	$\alpha \in \zirkel(\{0,1\})$ nach \ref{sub:1011}.
	\item[1) $\Rightarrow $ 2:] Sei $\alpha \in \zirkel( \{0,1\})$. Nach \ref{sub:1011} existiert eine Folge $\mathds{Q}= K_0 \subseteq K_1 \subseteq \ldots \subseteq K_k$
	mit $\alpha \in K_k$, $K_i = K_i(\omega_i)$ mit $\omega_i^2 \in K_{i-1}$. Sei $p_i$ das Minimalpolynom von $\omega_i$ über $\mathds{Q}$. Sei
	$f_i = p_1 \cdot \ldots \cdot p_i$ und $Z_i$ der Zerfällungskörper von $f_i$. Dann $K_i \subseteq Z_i$ und $\mathds{Q} \subseteq Z_1  \subseteq  \ldots \subseteq Z_k$.\\	
	\uline{Behauptung:} $[Z_i : Z_{i-1}]$ ist eine Potenz von $2$. Dann ist auch $[Z_k : \mathds{Q}]$ eine Potenz von $2$. Da $\nicefrac{Z_k}{\mathds{Q}}$ normal ist und
	$\alpha \in K_k \subseteq Z_k$ enthält $Z_k$ den Zerfällungskörper $L$ von $p_\alpha$. Damit ist auch $[L:\mathds{Q}]$ eine Potenz von $2$ (da 
	$[L:\mathds{Q}] \mid [Z_k :\mathds{Q}]$). Es bleibt die Behauptung zu beweisen.
	
	Seien $\omega_i = \alpha_1, \ldots , \alpha_s$ die Nullstellen von $p_i$. Nach dem Fortsetzungssatz für normale Körpererweiterungen (\ref{sub:1416}) gibt es 
	$\sigma_j  \in \Gal (\nicefrac{Z_i}{\mathds{Q}})$ mit $\sigma_j(\alpha_1) = \alpha_j$. Da $\nicefrac{Z_{i-1}}{\mathds{Q}}$ normal ist, gilt 
	$\sigma_j(Z_{i-1}) = Z_{i-1}$ nach \ref{sub:154}. Es ist $\alpha_1^2 = \omega_i^2 \in K_{i-1} \subseteq Z_{i-1}$. Daher 
	$\alpha_j^2 = \sigma_j(\alpha_1^2) \in \sigma_j(Z_{i-1}) = Z_{i-1}$. Damit auch
	\[
		Z_{i-1} \subseteq Z_{i-1}(\alpha_1) \subseteq Z_{i-2}(\alpha_1, \alpha_2) \subseteq  \ldots \subseteq Z_{i-1} (\alpha_1, \ldots , \alpha_s) = Z_i
	\]
	eine Folge von quadratischen Körpererweiterungen. Daher ist $[Z_i : Z_{i-1}]$ eine Potenz von 2. \bewende
\end{description}
% subsection 201 (end)

\subsection{Satz (Gauß)} % (fold)
\label{sub:202}
Das regelmäßige $n$-Eck ist genau dann mit Zirkel und Lineal konstruierbar, wenn $n = 2^e \cdot p_1 \cdot \ldots \cdot p_r$ mit beliebigem $e$ und paarweise verschiedenen
Primzahlen $p_1, \ldots , p_r$ der Form $p_i = 1+ 2^{k_i}$.
\minisec{Beweis}
Sei $\zeta := e^{2 \pi  i / n}$. Es ist $\mathds{Q}_n := \mathds{Q}(\zeta)$ der Zerfällungskörper von $p_\zeta$ und es gilt nach Kapitel \ref{sec:19}
$[\mathds{Q}_n :\mathds{Q}] = \varphi(n)$. Ist $n=2^e \cdot p_1^{e_1} \cdot \ldots \cdot  p_r^{e_r}$ die Primfaktorzerlegung von $n$, so gilt
\[
	\varphi(n) = 2^{e-1} \cdot p_1^{e_1 -1} \cdot (p_1 -1) \cdot \ldots \cdot p_r^{e_r -1} \cdot (p_r-1) \tag{siehe \ref{sub:196}}
\]
Es ist also $\varphi(n)$ genau dann eine Potenz von $2$, wenn $e_1= \ldots = e_r =1$ und $(p_i-1)$ für $i=1, \ldots ,r$ eine Potenz von $2$ ist, also $p_i= 2^{k_i}+1$.
Mit dem vorigen Satz (\ref{sub:201}) folgt die Behauptung. \bewende
% subsection 202 (end)

\subsection[Lemma über Primzahlen der Form $1+2^m$]{Lemma} % (fold)
\label{sub:203}
Sei $m \in \mathds{N}$ und $1+ 2^m$ eine Primzahl. Dann ist $m$ eine Potenz von $2$.
\minisec{Beweis}
Ist $m$ keine Potenz von $2$, so können wir $m=k \cdot p$ mit $p>2$ ungerade schreiben. Dann ist\footnote{besser bei \url{http://de.wikipedia.org/wiki/Fermat-Zahl}}
\[
	(1+2^m) = 1- (-2^k)^p = (1+2^k) \Big(1+ (-2)^k + (-2)^{2 \cdot k} + \ldots + (-2)^{(p-1)\cdot k}\Big) \bewende
\] 
% subsection 203 (end)

\subsection[Bemerkung zu Fermatschen Zahlen bzw. Fermatschen Primzahlen]{Bemerkung} % (fold)
\label{sub:204}
$F_k = 1+2^{(2^k)}$ heißt die $k$-te \Index{Fermatsche Zahl}. Es sind $F_0 = 3, F_1 = 5, F_2=17, F_3 = 257, F_4 = 65537$ Primzahlen. $F_5$ ist keine Primzahl mehr.
Es ist nicht bekannt, ob es weitere Fermatsche Primzahlen gibt. 
% subsection 204 (end)
% section 20 (end)
\newpage

\section{Auflösen von algebraischen Gleichungen über $\mathds{Q}$, II.} % (fold)
\label{sec:21}
In Kapitel \ref{sec:21} seien alle Körper Unterkörper von $\mathds{C}$.

\subsection[Definition: Auflösbare Körpererweiterung]{Definition} % (fold)
\label{sub:21.1}
Eine endliche Körpererweiterung $\nicefrac{L}{K}$ heißt \Index{auflösbar}, wenn es eine endliche Körpererweiterung $\nicefrac{E}{L}$ gibt, so dass $\nicefrac{E}{K}$
galoissch ist mit auflösbarer\footnote{$G$ heißt auflösbar, wenn es eine Folge von Normalteilern $\set{e} = N_0 \unlhd N_1 \unlhd \ldots \unlhd N_k =G$ gibt,
so dass für $i=1, \ldots ,k$ $\nicefrac{N_i}{N_{i-1}}$ abelsch ist; siehe auch \ref{sub:1211}} Galois-Gruppe $\Gal(\nicefrac{E}{K})$. 
% subsection 211 (end)

\subsection[Bemerkung: Untergruppen einer auflösbaren Gruppe sind auflösbar]{Bemerkung} % (fold)
\label{sub:21.2}
$G$ auflösbar, $H \le G$, so ist auch $H$ auflösbar. Ist $N \unlhd G$, so gilt: $G$ auflösbar $\iff N, \nicefrac{G}{N}$ auflösbar.
% subsection 21.2 (end)

\subsection[Wiederholung: Radikalerweiterungen und durch Radikale auflösbar]{Wiederholung} % (fold)
\label{sub:21.3}
\begin{itemize}
	\item $\nicefrac{L}{K}$ heißt Radikalerweiterung $:\Leftrightarrow \alpha \in L : L=K(\alpha), \alpha^d \in K$, $d \in \mathds{N}$
	\item $\nicefrac{L}{K}$ heißt durch Radikale auflösbar, wenn es eine Kette von Körpern gibt mit $K=K_0 \subseteq K_1 \subseteq \ldots \subseteq K_r$ mit 
	$L \subseteq K_r$ und $\nicefrac{K_i}{K_{i-1}}$ für $i=1, \ldots ,r$ eine Radikalerweiterung ist.
\end{itemize}
% subsection 21.3 (end)

\subsection[Bemerkung zur Auflösbarkeit von Galoisgruppen]{Bemerkung} % (fold)
\label{sub:21.4}
Ist $\nicefrac{L}{K}$ galoissch und $\nicefrac{E}{L}$ endlich mit $\nicefrac{E}{K}$ galoissch und $\Gal(\nicefrac{E}{K})$ auflösbar, so ist $\Gal(\nicefrac{L}{K})$
auflösbar, da nach dem Hauptsatz der Galoistheorie (\ref{sub:178}) gilt
\[
	\Gal(\nicefrac{L}{K}) \cong \nicefrac{\Gal(\nicefrac{E}{K})}{\Gal(\nicefrac{E}{L})}
\]
Also gilt für $\nicefrac{L}{K}$ endlich galoissch: $\nicefrac{L}{K} \text{ auflösbar } \iff \Gal(\nicefrac{L}{K}) \text{ auflösbar}$
% subsection 21.4 (end)

\subsection[Satz: Durch Radikale auflösbar ist äquivalent zu auflösbar]{Satz} % (fold)
\label{sub:215}
Sei $\nicefrac{L}{K}$ eine endliche Körpererweiterung. Dann sind äquivalent
\begin{multicols}{2}
	\begin{enumerate}[(1)]
		\item $\nicefrac{L}{K}$ ist durch Radikale auflösbar 
		\item $\nicefrac{L}{K}$ ist auflösbar.
	\end{enumerate}
\end{multicols}
% subsection 215 (end)

\subsection[Korollar: Die Galois-Gruppen von durch Radikale auflösbaren KE sind auflösbar]{Korollar} % (fold)
\label{sub:216}
Sei $\nicefrac{L}{K}$ eine endliche Galoiserweiterung. Dann gilt
\[
	\nicefrac{L}{K} \text{ ist durch Radikale auflösbar } \iff \Gal(\nicefrac{L}{K}) \text{ ist auflösbar}
\] 
\minisec{Beweis}
Folgt aus \ref{sub:21.4} und \ref{sub:215}. \bewende
% subsection 216 (end)

\subsection[Lemma 1: Auflösbarkeit in einem kommutierenden Diagramm]{Lemma 1} % (fold)
\label{sub:217}
\begin{minipage}{0.95\textwidth}
	Seien $L,K,F$ Unterkörper von $\mathds{C}$ mit $K \subseteq L$, $K \subseteq F$, $\nicefrac{L}{K}$ endlich. Sei $E := F(L)$. Dann gilt
	\begin{enumerate}[a)]
		\item Ist $\nicefrac{L}{K}$ durch Radikale auflösbar, so ist auch $\nicefrac{E}{F}$ durch Radikale auflösbar.
		\item Ist $\nicefrac{L}{K}$ auflösbar, so ist auch $\nicefrac{E}{F}$ auflösbar.
	\end{enumerate}
\end{minipage}
\begin{minipage}{0.1\textwidth}
	\[
		\begin{tikzcd}
			L \rar[hook]&  F(L)\\
			K \uar[hook] \rar[hook] & F \uar[hook]
		\end{tikzcd}
	\]
\end{minipage}
\minisec{Beweis}
\begin{enumerate}[a)]
	\item Sei $\nicefrac{L}{K}$ durch Radikale auflösbar. Dann existiert $K= K_0 \subseteq K_1 \subseteq \ldots \subseteq K_r$ mit 
	$L \subseteq K_r$ und $\exists \alpha_i \in K_i, d_i \in \mathds{N}$ mit $K_i= K_{i-1}(\alpha_i)$ und $\alpha_i^{d_i} \in K_{i-1}$.
	Es folgt $K_r = K(\alpha_1, \ldots , \alpha_r)$ und
	\[
		E=F(L) \subseteq F(K_r) = F\big(K(\alpha_1, \ldots , \alpha_r)\big) = F(\alpha_1, \ldots , \alpha_r)
	\]
	Setze nun $F_i := F(\alpha_1, \ldots , \alpha_i)$. Dann $F=F_0 \subseteq F_1 \subseteq \ldots \subseteq F_r$ mit $E \subseteq F_r$, $F_i = F_{i-1}(\alpha_i)$ mit
	$\alpha_i^{d_i} \in K_{i-1} \subseteq F_{i-1}$. Daher ist auch $\nicefrac{E}{F}$ durch Radikale auflösbar.
	\item Sei $\nicefrac{L}{K}$ auflösbar. Wir können o.B.d.A. annehmen, dass $\nicefrac{L}{K}$ galoissch ist mit $\Gal(\nicefrac{L}{K})$ auflösbar. Insbesondere ist 
	$L$ der Zerfällungskörper eines Polynoms $p \in K[X]$ (siehe \ref{sub:154}). Dann ist $E=F(L)$ der Zerfällungskörper von $p$ über $F$. Insbesondere ist 
	$\nicefrac{E}{F}$ normal, endlich und damit galoissch. Es bleibt zu zeigen: $\Gal(\nicefrac{E}{F})$ ist auflösbar.
	
	Da $\nicefrac{L}{K}$ normal ist, erhalten wir durch Einschränkung auf $L$ einen Gruppenhomomorphismus 
	$\Aut (\nicefrac{E}{K}) \to\Aut(\nicefrac{L}{K}) = \Gal(\nicefrac{L}{K})$. Da $\Gal(\nicefrac{E}{F}) \le \Aut(\nicefrac{E}{K})$ erhalten wir einen Gruppenhomomorphismus
	\[
		\psi : \Gal(\nicefrac{E}{F}) \to \Gal(\nicefrac{L}{K}), \quad \sigma \mapsto \sigma\big|_L
	\]
	Ist $\sigma|_L = \id$ so folgt $\sigma= \id$, da $\sigma$ ein $F$-Homomorphismus ist und $E=F(L)$. Damit ist $\psi$ injektiv und 
	\[
		\Gal(\nicefrac{E}{F}) \cong \bild \psi \le \Gal(\nicefrac{L}{K}).
	\]
	Damit ist $\Gal(\nicefrac{E}{F})$ isomorph zu einer Untergruppe einer auflösbaren Gruppe also auflösbar.\bewende
\end{enumerate}
% subsection 217 (end)

\subsection[Lemma 2: Transitivität der beiden Auflösbarkeitsbegriffe]{Lemma 2} % (fold)
\label{sub:218}
Seien $K \subseteq L \subseteq M$ Unterkörper von $\mathds{C}$. Sei $\nicefrac{M}{K}$ endlich. Dann gilt
\begin{enumerate}[a)]
	\item $\nicefrac{M}{K}$ durch Radikale auflösbar $\iff \nicefrac{L}{K}$ und $\nicefrac{M}{L}$ durch Radikale auflösbar.
	\item $\nicefrac{M}{K}$ auflösbar $\iff \nicefrac{L}{K}$ und $\nicefrac{M}{L}$ auflösbar.
\end{enumerate}
\minisec{Beweis}
\begin{enumerate}[a)]
	\item\glqq$\Longrightarrow $\grqq{} folgt direkt aus der Definition.
	
	\glqq$\Longleftarrow$\grqq{} Ist $\nicefrac{L}{K}$ durch Radikale auflösbar, so gibt es eine Folge $K=K_0 \subseteq K_1 \subseteq \ldots \subseteq K_r$ von 
	Radikalerweiterungen mit $L \subseteq K_r$. Mit $\nicefrac{M}{L}$ ist nach Lemma 1 (\ref{sub:217}) auch $\nicefrac{M(K_r)}{K_r}$ durch Radikale auflösbar. Es gibt also
	eine Folge $K_r \subseteq K_{r+1} \subseteq \ldots \subseteq K_R$ von Radikalerweiterungen mit $M \subseteq M(K_r) \subseteq K_R$. Insgesamt beweist die Folge
	\[
		K=K_0 \subseteq K_1 \subseteq \ldots \subseteq K_r \subseteq K_{r+1} \subseteq \ldots \subseteq K_R,
	\]
	dass $\nicefrac{M}{K}$ durch Radikale auflösbar ist.
	\item \glqq$\Longrightarrow$\grqq{} Sei $\nicefrac{M}{K}$ auflösbar. Wir können o.B.d.A annehmen, dass $\nicefrac{M}{K}$ eine endliche Galoiserweiterung mit 
	auflösbarer Galois-Gruppe $\Gal(\nicefrac{M}{K})$ ist (Sonst vergrößern wir $M$). Dann ist auch $\nicefrac{L}{K}$ auflösbar. Mit $\nicefrac{M}{K}$ ist auch 
	$\nicefrac{M}{L}$ eine endliche Galois-Erweiterung. Weiter ist $\Gal(\nicefrac{M}{L})$ als Untergruppe der auflösbaren Gruppe $\Gal(\nicefrac{M}{K})$ auflösbar.
	
	
	
	\begin{minipage}{0.75\textwidth}
		\glqq$\Longleftarrow$\grqq{} Ist $\nicefrac{L}{K}$ auflösbar, so gibt es $\nicefrac{L'}{L}$ endlich mit $\nicefrac{L'}{K}$ galoissch und $\Gal(\nicefrac{L'}{K})$
		auflösbar. Mit $\nicefrac{M}{L}$ ist nach Lemma 1 (\ref{sub:217}) auch $\nicefrac{L'(M)}{L'}$ auflösbar. Es gibt also $\nicefrac{M'}{L'(M)}$ endlich mit 
		$\nicefrac{M'}{L'}$ galoissch mit auflösbarer Galois-Gruppe $\Gal(\nicefrac{M'}{L'})$.
		Angenommen $\nicefrac{M'}{K}$ ist galoissch. Dann $\Gal(\nicefrac{M'}{L'}) \unlhd \Gal(\nicefrac{M'}{K})$ und
		\[
			\nicefrac{\Gal(\nicefrac{M'}{K})}{\Gal(\nicefrac{M'}{L'})} \unlhd  \Gal( \nicefrac{L'}{K}).
		\]
	\end{minipage}\quad 
	\begin{minipage}{0.17\textwidth}
		\[
			\begin{tikzcd}
						M \rar[hook]&  M'\\
						  & L' \uar[hook] \\
						K \arrow[hook]{uu} \rar{=} & K \uar[hook]
			\end{tikzcd}
		\]
	\end{minipage}
	
	
	Da $\nicefrac{M'}{K}$ endlich und separabel ist, gibt es $\alpha_1 \in M'$ mit $M'=K(\alpha_1)$ (Satz vom primitiven Element, \ref{sub:1615}). Sei $Z$ der 
	Zerfällungskörper von $p_{\alpha_1} \in K[X]$. Also $p_{\alpha_1} = (X-\alpha_1) \cdot \ldots \cdot (X-\alpha_r)$ und $Z=K(\alpha_1, \ldots , \alpha_r)$.
	Es ist $M' \subseteq Z$ und $\nicefrac{Z}{K}$ ist eine endliche Galois-Erweiterung. Es ist $L'$ ein Zwischenkörper von $\nicefrac{Z}{K}$. Da $\nicefrac{L'}{K}$ normal
	ist, ist nach dem Hauptsatz der Galoistheorie (\ref{sub:178}) die Restriktionsabbildung $\Gal(\nicefrac{Z}{K}) \twoheadrightarrow \Gal(\nicefrac{L'}{K})$ surjektiv mit 
	Kern $\Gal(\nicefrac{Z}{L'})$. Es ist $\Gal(\nicefrac{L'}{K})$ auflösbar.
	
	Bleibt zu zeigen: $\Gal(\nicefrac{Z}{L'})$ ist auflösbar. Nach Fortsetzungssatz gibt es $\sigma_j \in \Gal(\nicefrac{Z}{K})$ mit $\sigma_j(\alpha_1) = \alpha_j$.
	Sei $\sigma_j(M') =: M_j'$. Da $\nicefrac{L'}{K}$ normal ist, ist $\sigma_j(L') = L'$. $\nicefrac{M'}{L'}$ normal impliziert $\nicefrac{M_j'}{L'}$ normal.
	Aus $\Gal(\nicefrac{M'}{L'})$ auflösbar folgt $\Gal(\nicefrac{M_j'}{L'})$ auflösbar. Für jedes $j$ erhalten wir Restriktionsabbildungen
	\[
		\psi_j : \Gal (\nicefrac{Z}{L'}) \to \Gal(\nicefrac{M_j'}{L'})
	\]
	Sei $\psi : \Gal(\nicefrac{Z}{L'}) \to \Gal(\nicefrac{M_1'}{L'}) \times \ldots \times \Gal(\nicefrac{M_r'}{L'})$ das Produkt der $\psi_j$. Da $\alpha_i \in M_i'$ und
	$Z=K(\alpha_1, \ldots , \alpha_r)$ ist $\psi$ injektiv. Damit ist $\Gal(\nicefrac{Z}{L'})$ auflösbar. \bewende
\end{enumerate}
% subsection 218 (end)

\subsection{Beweis von \ref{sub:215}} % (fold)
\label{sub:219}
\begin{description}
	\item[(1) $\Rightarrow$ (2):] Sei $\nicefrac{L}{K}$ durch Radikale auflösbar. Dann gibt es $K=K_0 \subseteq K_1 \subseteq \ldots \subseteq K_n$ und $\alpha_1 \in K_i$,
	$d_i \in \mathds{N}$ mit $K_i = K_{i-1}(\alpha_i)$, $\alpha_i^{d_i} \in K_{i-1}$ und $L \subseteq K_n$. Sei $N:= d_1 \cdot \ldots \cdot d_n$. Sei 
	$\zeta \in PE_n(\mathds{C})$. Betrachte 
	\[
		K = K_0 \subseteq K_1(\zeta) \subseteq \ldots \subseteq K_n(\zeta)
	\]
	Es ist $\nicefrac{\mathds{Q}_N}{\mathds{Q}}$ auflösbar, da $\Gal(\nicefrac{\mathds{Q}_N}{\mathds{Q}}) \cong \enbrace*{\nicefrac{\mathds{Z}}{n \mathds{Z}}}^\times $
	abelsch ist. Mit Lemma 1 (\ref{sub:217}) folgt, dass $\nicefrac{K_0(\zeta)}{K}$ auflösbar ist. Wegen Lemma 3 (\ref{sub:2110}) sind 
	$\nicefrac{K_1(\zeta)}{K_0(\zeta)}, \ldots , \nicefrac{K_n(\zeta)}{k_{n-1}(\zeta)}$ auflösbar. Mit Lemma 2 b) und einer Induktion folgt, dass
	$\nicefrac{K_n(\zeta)}{K}$ auflösbar ist. Damit ist auch $\nicefrac{L}{K}$ auflösbar, denn $L \subseteq K_n \subseteq K_n(\zeta)$. \bewende
	\item[(2) $\Rightarrow $ (1):] Siehe Bosch "{}Algebra"{} S. 265
\end{description}
% subsection 219 (end)

\subsection[Lemma 3: Auflösbarkeit einer Radikalerweiterung, die alle Einheitswurzeln enthält]{Lemma 3} % (fold)
\label{sub:2110}
Sei $L=K(\alpha)$ mit $\alpha^n \in K, \alpha \not= 0$. Weiter gelte $E_n(\mathds{C}) \subseteq L$. Dann ist $\nicefrac{L}{K}$ auflösbar.
\minisec{Beweis}
Die Nullstellen von $X^n- \alpha^n \in K[X]$ sind genau die $\zeta \cdot \alpha$ mit $\zeta \in E_n(\mathds{C})$. Da $E_n(\mathds{C}) \subseteq K \subseteq K(\alpha)$ ist
$K(\alpha)$ der Zerfällungskörper von $X^n-\alpha^n$ und damit endlich und galoissch. Die Galois-Gruppe $\Gal(\nicefrac{L}{K})$ wirkt auf den Nullstellen von 
$X^n-\alpha^n$. Zu $\sigma \in \Gal(\nicefrac{L}{K})$ gibt es also $\zeta_\sigma \in E_n(\mathds{C})$ mit $\sigma(\alpha) = \zeta_\sigma \cdot \alpha$.

Behauptung: $\psi : \Gal(\nicefrac{L}{K}) \to E_n{\mathds{C}}, \sigma \mapsto \zeta_\sigma$ ist ein Gruppenhomomorphismus.\\
$\sigma \tau (\alpha) = \sigma( \zeta_\tau \cdot  \alpha) = \zeta_\tau \cdot \sigma(\alpha) = \zeta_\tau \cdot  \zeta_\sigma  \cdot \alpha$. Also $\zeta_{\sigma \tau} = \zeta_\sigma\cdot \zeta_\tau$.
Wegen $L=K(\alpha)$ ist $\psi$ injektiv. Da $E_n(\mathds{C})$ abelsch ist, ist auch $\Gal(\nicefrac{L}{K})$ abelsch und damit auflösbar. \bewende
% subsection 2120 (end)
% section 21 (end)
\newpage

\section{Eine nicht-auflösbare Gleichung} % (fold)
\label{sec:22}

\subsection[Definition: Alternierende Gruppe $A_n$]{Definition} % (fold)
\label{sub:221}
Der Kern der \Index{Signum-Abbildung} $\sgn : S_n \to \set{\pm 1} $ heißt die \bet{$n$-te alternierende Gruppe $A_n$}\index{n-te alternierende Gruppe A_n@$n$-te alternierende Gruppe $A_n$}.
% subsection 221 (end)

\subsection[Definition: Kommutatoruntergruppe]{Definition} % (fold)
\label{sub:222}
Sei $G$ eine Gruppe. Die kleinste Untergruppe von $G$, die alle Elemente der Form $g h g ^{-1} h ^{-1} =: [g,h]$, $g,h \in G$ enthält, heißt die 
\Index{Kommutatoruntergruppe} von $G$ und wird mit $[G,G]$ bezeichnet.
% subsection 222 (end)

\subsection[{Lemma: $[G,G] \subseteq \ker \varphi$, wenn $\varphi$ auf abelsche Gruppe abbildet}]{Lemma} % (fold)
\label{sub:223}
Sei $\varphi : G \to A$ ein Gruppenhomomorphismus mit $A$ abelsch. Dann $[G,G] \subseteq \ker \varphi$.
\minisec{Beweis}
\begin{align*}
	\varphi \enbrace*{[g,h]} = \varphi \enbrace*{g h g ^{-1} h ^{-1}} = \varphi(g) \cdot \varphi(h) \cdot \varphi(g) ^{-1}  \cdot \varphi(h) ^{-1} = 
	\varphi(g) \cdot \varphi(g) ^{-1} \cdot \varphi(h) \cdot \varphi(h) ^{-1} = e \bewende
\end{align*}
% subsection 223 (end)

\subsection[Satz: Kommutatoruntergruppe von $A_n$ für $n \ge 5$]{Satz} % (fold)
\label{sub:224}
Für $n\ge 5$ ist $[A_n,A_n] = A_n$.\marginnote{$\nicefrac{S_n}{A_n} \cong \set{\pm 1} $}
\minisec{Beweis}
$S_n$ wird von den \Index{Transpositionen} $(x \, y)$, $x \not= y \in \{1,\ldots ,n\}$ erzeugt, d.h. jedes Element in $S_n$ ist ein Produkt von Transpositionen. 
$\sigma \in A_n$ gilt genau dann, wenn $\sigma$ ein Produkt einer geraden Anzahl von Transpositionen ist. Es genügt also zu zeigen
\[
	\forall x \not= y, y' \not= z \text{ ist } (x \, y) (y' \, z) \in [A_n,A_n]
\]
O.B.d.A. Sei $x \not= z$. Ist $y=y'$, so gibt es $a \not= b \in \{1, \ldots , n\} \setminus \{x,y,z\}$, da $n \ge 5$. Dann gilt
\begin{align*}
	(x\, y) (y \, z) &= [(x \, y)(y \, a) , (x \, z) (z \, b)]  \in [A_n, A_n]\\
	&= (x \, y) (y \, a) (x \, z) (z \, b) (y \, a) (x \, y) (z \, b) (x \, z)
\end{align*}
Ist $y \not= y'$, so 
\begin{align*}
	(x \, y) (y' \, z) = \underbrace{(x \, y) (y \, y')}_{\in [A_n, A_n]} \underbrace{(y \, y') (y' \, z)}_{\in [A_n, A_n]} \in [A_n, A_n]\bewende
\end{align*}
% subsection 224 (end)

\subsection[Korollar: Auflösbarkeit von $A_n$ für $n\ge 5$]{Korollar} % (fold)
\label{sub:225}
$A_n$ ist für $n \ge 5$ nicht auflösbar.
\minisec{Beweis}
Angenommen doch: $\{e\} = N_0 \unlhd N_1 \unlhd N_2 \unlhd \ldots \unlhd N_k = A_n$ mit $\nicefrac{N_i}{N_{i-1}}$ abelsch. \marginnote{$\unlhd$ durchstreichen}
Insbesondere ist $\nicefrac{A_n}{N_{k-1}}$ abelsch und wir haben einen surjektiven Gruppenhomomorphismus $\pi : A_n \twoheadrightarrow \nicefrac{A_n}{N_{k-1}}$.
Mit Lemma \ref{sub:223} folgt
\[
	[A_n,A_n] \subseteq \Kern \pi \Longrightarrow A_n  \subseteq \Kern \pi  \text{ und } \pi \text{ trivial } \lightning
\]
% subsection 225 (end)

\subsection[Korollar: $S_n$ ist für $n \ge 5$ nicht auflösbar]{Korollar} % (fold)
\label{sub:226}
Für $n \ge 5$ ist $S_n$ nicht auflösbar.
% subsection 226 (end)

\subsection[{Satz: Die Galois-Gruppe von $X^5- 4X+2 \in \mathds{Q}[X]$ ist $S_5$}]{Satz} % (fold)
\label{sub:227}
Die Galois-Gruppe von $X^5- 4X+2 \in \mathds{Q}[X]$ ist $S_5$.
\minisec{Beweis}
Sei $L$ der Zerfällungskörper von $p$. Seien $T= \set{t_0, t_1, t_2, z, \overline{z} } $ die Nullstellen von $p$.\marginnote{$z, \overline{z} \not\in \mathds{R} $, 
$t_0,t_1, t_2 \in \mathds{R}$} Sei $G := \Gal(\nicefrac{L}{\mathds{Q}})$. Da $L=\mathds{Q}(T)$ und $G$ auf $T$ wirkt, können wir $G$ als Untergruppe von
$S_T (\cong S_5)$ auffassen. Nach Eisenstein ist $p$ irreduzibel (teste die Primzahl 2). Mit dem Fortsetzungssatz (\ref{sub:149}) folgt: Zu $\alpha, \beta \in T$ gibt es
$\sigma \in G$ mit $\sigma(\alpha)=\beta$. Es gibt also nur eine Bahn (siehe \ref{sub:35}) für die Wirkung von $G$ auf $T$. Es folgt mit der Bahnengleichung (\ref{sub:310})
\[
	\frac{\abs*{G} }{\abs*{G_{t_0}}} = \abs*{G t_0} = \abs*{T} = 5  
\]
Es folgt $5 \mid \abs*{G}$. Damit hat $G$ eine nichttriviale $5$-Sylowgruppe. Da $G \le S_T$ gilt $\abs*{G} \mid \abs*{S_T} = \abs*{S_5} = 120 $. Daher hat die $5$-Sylowgruppe
Ordnung $5$. Damit hat jedes nichttriviales Element in ihr Ordnung $5$. Insbesondere enthält $G$ ein Element $\sigma$ von Ordnung $5$. Weiter ist die komplexe Konjugation
ein Element der Ordnung $2$.
% subsection 227 (end)

\subsection[{Korollar: Der Zerfällungskörper von $X^5-4X +2 \in \mathds{Q}[X]$}]{Korollar} % (fold)
\label{sub:228}
Der Zerfällungskörper von $X^5-4X +2 \in \mathds{Q}[X]$ ist nicht durch Radikale auflösbar.
% subsection 228 (end)

\subsection[Lemma: Nullstellen von $X^5- 4X +2$]{Lemma} % (fold)
\label{sub:229}
$p=X^5-4X+2$ hat genau $3$ reelle Nullstellen und $2$ weitere Nullstellen, $z$ und $\overline{z} $ in $\mathds{C} \setminus \mathds{R}$.
\minisec{Beweis}
Betrachte $p$ als differenzierbare Abbildung $\mathds{R} \to \mathds{R}$. Es gilt $\lim_{ t \to \infty}p(t) = \infty $ und $\lim_{ t \to -\infty} = - \infty $.
Es ist $p' = 5 X^4 -4 = 5 (X^2- \frac{2}{\sqrt{5}  } ) (X^2+ \frac{2}{\sqrt{5}})$. Also hat $p'$ genau zwei reelle Nullstellen $\pm \frac{\sqrt{2}  }{\sqrt[4]{5}}$.
\[
	p \enbrace*{\frac{\sqrt{2}  }{\sqrt[4]{5}  } } < 0 \qquad p \enbrace*{ - \frac{\sqrt{2}  }{\sqrt[4]{5}  } } >0
\]
Daher hat $p$ genau $3$ reelle Nullstellen. Da $p' \not= 0$ in diesen $3$ Nullstellen sind die reellen Nullstellen alle einfache Nullstellen. Da $\deg p =5$ gibt es noch
mindestens eine weitere Nullstelle $z \in \mathds{C} \setminus \mathds{R}$. Dann ist auch $\overline{z} \not= z $ eine Nullstelle: $p(\overline{z} ) = \overline{p(z)} =0.$
\bewende
% subsection 229 (end)

\subsection[Lemma: Erzeuger von $S_5$]{Lemma} % (fold)
\label{sub:230}
Seien $\tau, \sigma \in S_5$ mit $\ord(\sigma)=5$ und $\tau$ eine Transposition. Dann erzeugen $\sigma$ und $\tau$ die Gruppe $S_5$, d.h. $S_5$ ist die einzige Untergruppe
von $S_5$, die $\sigma$ und $\tau$ enthält.
\minisec{Beweis}
o.B.d.A. sei $\tau = (1 \, 2)$. Da $\ord(\sigma)=5$ gibt es keine echte nicht leere Teilmenge von $T$ von $\set{1,\ldots ,5} $ mit $\sigma(T)=T$. Es folgt 
\[
	\set{1,\ldots ,5} = \set{1, \sigma(1), \sigma^2(1), \sigma^3(1), \sigma^4(1)}.  
\]
Indem wir $\sigma$ durch eine geeignete Potenz von $\sigma$ ersetzen dürfen wir annehmen $\sigma(1)=2$. Nach Umordnung von $3,4,5$ können wir annehmen
\[
	\sigma(i) = \begin{cases}
		i+1, &\text{ falls }i \not= 5\\
		1 , &\text{ falls } i =5
	\end{cases}
\]
Nun ist $ (2 \, 3)= \sigma (1 \, 2) \sigma ^{-1}$, $(3 \, 4) = \sigma (2 \,3) \sigma ^{-1}$, $(4 \, 5) = \sigma (3 \, 4) \sigma ^{-1}$. Also liegen
$(1\, 2), (2 \, 3), (3 \, 4), (4 \, 5)$ in der von $\sigma$ und $\tau$ erzeugten Untergruppe von $S_5$. Durch diese Transpositionen wird aber schon die ganze $S_5$ erzeugt:
Da sich jede endliche Folge von Zahlen durch sukzessive Vertauschungen von benachbarten Zahlen ordnen lässt. \bewende
% subsection 230 (end)

% section 22 (end)
\cleardoubleoddemptypage
\pagenumbering{Alph}
\setcounter{page}{1}
\printindex
\listoffigures
\listoftables
% \todototoc
% \listoftodos[Todo's und andere Baustellen]
\end{document}
